% concatenated from SIMS.tex
%%
%% This is file `sample-sigconf.tex',
%% generated with the docstrip utility.
%%
%% The original source files were:
%%
%% samples.dtx  (with options: `all,proceedings,bibtex,sigconf')
%% 
%% IMPORTANT NOTICE:
%% 
%% For the copyright see the source file.
%% 
%% Any modified versions of this file must be renamed
%% with new filenames distinct from sample-sigconf.tex.
%% 
%% For distribution of the original source see the terms
%% for copying and modification in the file samples.dtx.
%% 
%% This generated file may be distributed as long as the
%% original source files, as listed above, are part of the
%% same distribution. (The sources need not necessarily be
%% in the same archive or directory.)
%%
%%
%% Commands for TeXCount
%TC:macro \cite [option:text,text]
%TC:macro \citep [option:text,text]
%TC:macro \citet [option:text,text]
%TC:envir table 0 1
%TC:envir table* 0 1
%TC:envir tabular [ignore] word
%TC:envir displaymath 0 word
%TC:envir math 0 word
%TC:envir comment 0 0
%%
%% The first command in your LaTeX source must be the \documentclass
%% command.
%%
%% For submission and review of your manuscript please change the
%% command to \documentclass[manuscript, screen, review]{acmart}.
%%
%% When submitting camera ready or to TAPS, please change the command
%% to \documentclass[sigconf]{acmart} or whichever template is required
%% for your publication.
%%
%%

%快速写蓝色字体
\newcommand{\edit}[1]{\textcolor{black}{#1}}
\newcommand{\fei}[1]{\textcolor{red}{[Fei: #1]}}
\newcommand{\zbedit}[1]{\textcolor{black}{#1}}
% \newcommand{\add}[1]{\textcolor{blue}{#1}}
\newenvironment{add}
{\color{black}}
{}

%快速注释footnote
\renewcommand{\footnote}[1]{}
\newif\ifcycomment
%\cycommenttrue   % 打开 cy 注释脚注
\cycommentfalse % 注释掉（关闭）cy 注释脚注

\newcommand{\cyfoot}[1]{\ifcycomment\footnote{#1}\fi}

\documentclass[sigconf]{acmart}
%,anonymous,review

\usepackage{amsthm}
\usepackage{multirow}
\newtheorem{assumption}{Assumption}
\usepackage{algorithm}
\usepackage{algpseudocode}
\algrenewcommand\algorithmicrequire{\textbf{Input:}}
\algrenewcommand\algorithmicensure{\textbf{Output:}}

%%
%% \BibTeX command to typeset BibTeX logo in the docs
\AtBeginDocument{%
  \providecommand\BibTeX{{%
    Bib\TeX}}}

%% Rights management information.  This information is sent to you
%% when you complete the rights form.  These commands have SAMPLE
%% values in them; it is your responsibility as an author to replace
%% the commands and values with those provided to you when you
%% complete the rights form.
\setcopyright{cc}
\setcctype{by}
\copyrightyear{2026}
\acmYear{2026}
\acmDOI{10.1145/3770855.3818042}
%% These commands are for a PROCEEDINGS abstract or paper.
\acmConference[KDD 2026] {Proceedings of the 32nd ACM SIGKDD Conference on Knowledge Discovery and Data Mining V.2}{August 9--13, 2026}{Jeju Island, Republic of Korea.}
\acmBooktitle{Proceedings of the 32nd ACM SIGKDD Conference on Knowledge Discovery and Data Mining V.2 (KDD 2026), August 9--13, 2026, Jeju Island, Republic of Korea}
%%
%%  Uncomment \acmBooktitle if the title of the proceedings is different
%%  from ``Proceedings of ...''!
%%
%%\acmBooktitle{Woodstock '18: ACM Symposium on Neural Gaze Detection,
%%  June 03--05, 2018, Woodstock, NY}
\acmISBN{979-8-4007-2259-2/2026/08}


%%
%% Submission ID.
%% Use this when submitting an article to a sponsored event. You'll
%% receive a unique submission ID from the organizers
%% of the event, and this ID should be used as the parameter to this command.
%%\acmSubmissionID{123-A56-BU3}

%%
%% For managing citations, it is recommended to use bibliography
%% files in BibTeX format.
%%
%% You can then either use BibTeX with the ACM-Reference-Format style,
%% or BibLaTeX with the acmnumeric or acmauthoryear sytles, that include
%% support for advanced citation of software artefact from the
%% biblatex-software package, also separately available on CTAN.
%%
%% Look at the sample-*-biblatex.tex files for templates showcasing
%% the biblatex styles.
%%

%%
%% The majority of ACM publications use numbered citations and
%% references.  The command \citestyle{authoryear} switches to the
%% "author year" style.
%%
%% If you are preparing content for an event
%% sponsored by ACM SIGGRAPH, you must use the "author year" style of
%% citations and references.
%% Uncommenting
%% the next command will enable that style.
%%\citestyle{acmauthoryear}


%%
%% end of the preamble, start of the body of the document source.
\settopmatter{printacmref=true}
\begin{document}

%%
%% The "title" command has an optional parameter,
%% allowing the author to define a "short title" to be used in page headers.
\title{SIMS: Scale-Invariant Merit-Function-Based Scalarization for Multi-Task Learning}

%%
%% The "author" command and its associated commands are used to define
%% the authors and their affiliations.
%% Of note is the shared affiliation of the first two authors, and the
%% "authornote" and "authornotemark" commands
%% used to denote shared contribution to the research.
\author{Zebin Chen}
\authornote{Both authors contributed equally to this research.}
\email{12432660@mail.sustech.edu.cn}
\affiliation{%
  \institution{Southern University of Science and Technology}
  \city{Shenzhen}
  \country{China}
}

%\authornotemark[1]
\author{Fei Xing}
\authornotemark[1]
\email{fxing8-c@my.cityu.edu.hk}
\affiliation{%
  \institution{City University of Hong Kong}
  \city{Hong Kong SAR}
  \country{China}}


\author{Yang Chen}
\email{cheny2023@mail.sustech.edu.cn}
\affiliation{%
  \institution{Southern University of Science and Technology}
  \city{Shenzhen}
  \country{China}
}

\author{Hua Liu}
\email{liuh5@sustech.edu.cn}
\affiliation{%
 \institution{Southern University of Science and Technology}
  \city{Shenzhen}
  \country{China}
}

\author{Andy H.F. Chow}
\email{andychow@cityu.edu.hk}
\affiliation{%
  \institution{City University of Hong Kong}
  \city{Hong Kong SAR}
  \country{China}}

\author{Yuhua Qian}
\email{jinchengqyh@126.com}
\affiliation{%
  \institution{Shanxi University}
  \city{Taiyuan}
  \country{China}}


\author{Yu Zhang}
\email{yu.zhang.ust@gmail.com}
\authornote{Yu Zhang is the corresponding author.}
\affiliation{%
  \institution{Southern University of Science and Technology}
  \city{Shenzhen}
  \country{China}}




%%
%% By default, the full list of authors will be used in the page
%% headers. Often, this list is too long, and will overlap
%% other information printed in the page headers. This command allows
%% the author to define a more concise list
%% of authors' names for this purpose.
\renewcommand{\shortauthors}{Zebin Chen, Fei Xing, Yang Chen, Hua Liu, Andy H.F. Chow, Yuhua Qian, Yu Zhang.}

%%
%% The abstract is a short summary of the work to be presented in the
%% article.
\begin{abstract}
% Multi-task learning (MTL) is an important paradigm in machine learning, yet its training dynamics must navigate unavoidable trade-offs among competing objectives.
Multi-task learning (MTL) requires navigating unavoidable trade-offs among competing objectives.
% This perspective naturally casts MTL as multi-objective optimization (MOO). Among MOO methods, scalarization is particularly appealing because it reduces an MOO problem to a single-objective one that can be solved using existing optimization techniques.
This paradigm is frequently formulated as multi-objective optimization (MOO), where the scalarization is favored to reduce an MOO problem to a single objective.
% \fei{reduces a MOO problem to a single-objective one that can be solved using existing optimization techniques}.
We empirically find that existing merit-function-based scalarization approaches are sensitive to the relative scales of different objectives in practical MTL, where task losses commonly differ by orders of magnitude. 
The optimization process often favors objectives with larger scales even though the underlying Pareto  optimal solutions remains invariant to rescaling (i.e., multiplying an objective by a positive constant).
% Although positive rescaling (i.e., multiplying an objective by a positive constant) preserves the Pareto set, it can drive 
% \fei{bias changes to drive} 
% optimization toward solutions dominated by the larger-scale objective.
To address this issue, we propose \textbf{S}cale-\textbf{I}nvariant \textbf{M}erit-function-based \textbf{S}calarization (\textbf{SIMS}) for MTL. Specifically, SIMS adopts a transformation-induced merit function to convert the MOO problem of MTL to a single objective \edit{that renders optimization invariant to} the magnitudes of losses.
% , derived from our transformation-induced merit-function-based scalarization.
Theoretically, we prove that the requirement for scale invariance uniquely determines this transformation to be logarithmic. We further show that \edit{this} general transformation-induced merit function preserves weak Pareto optimality and admits a smooth surrogate with controllable approximation error.
% Theoretically, we show that the induced merit function preserves weak Pareto optimality and admits a smooth surrogate with controllable approximation error. Moreover, we prove that scale invariance uniquely determines the transformation to be logarithmic, under which the formulation exactly recovers SIMS.
% \fei{add 'for MTL'}, 
% \fei{ a special case of our transformation-induced merit-based reformulation. Theoretically, we show that the induced merit function preserves weak Pareto optimality and admits a smooth surrogate with controllable approximation error. Moreover, we prove that scale invariance uniquely determines the transformation to be logarithmic, under which the formulation exactly recovers SIMS. Empirically,......} 
Extensive experiments on representative multi-task benchmarks demonstrate that SIMS consistently outperforms existing scalarization methods and achieves state-of-the-art performance.
\end{abstract}

%%
%% The code below is generated by the tool at http://dl.acm.org/ccs.cfm.
%% Please copy and paste the code instead of the example below.
%%
% \begin{CCSXML}
% <ccs2012>
%  <concept>
%   <concept_id>00000000.0000000.0000000</concept_id>
%   <concept_desc>Do Not Use This Code, Generate the Correct Terms for Your Paper</concept_desc>
%   <concept_significance>500</concept_significance>
%  </concept>
%  <concept>
%   <concept_id>00000000.00000000.00000000</concept_id>
%   <concept_desc>Do Not Use This Code, Generate the Correct Terms for Your Paper</concept_desc>
%   <concept_significance>300</concept_significance>
%  </concept>
%  <concept>
%   <concept_id>00000000.00000000.00000000</concept_id>
%   <concept_desc>Do Not Use This Code, Generate the Correct Terms for Your Paper</concept_desc>
%   <concept_significance>100</concept_significance>
%  </concept>
%  <concept>
%   <concept_id>00000000.00000000.00000000</concept_id>
%   <concept_desc>Do Not Use This Code, Generate the Correct Terms for Your Paper</concept_desc>
%   <concept_significance>100</concept_significance>
%  </concept>
% </ccs2012>
% \end{CCSXML}

% \ccsdesc[500]{Do Not Use This Code~Generate the Correct Terms for Your Paper}
% \ccsdesc[300]{Do Not Use This Code~Generate the Correct Terms for Your Paper}
% \ccsdesc{Do Not Use This Code~Generate the Correct Terms for Your Paper}
% \ccsdesc[100]{Do Not Use This Code~Generate the Correct Terms for Your Paper}
\begin{CCSXML}
<ccs2012>
   <concept>
       <concept_id>10010147.10010257.10010258.10010262</concept_id>
       <concept_desc>Computing methodologies~Multi-task learning</concept_desc>
       <concept_significance>500</concept_significance>
       </concept>
 </ccs2012>
\end{CCSXML}

\ccsdesc[500]{Computing methodologies~Multi-task learning}

%%
%% Keywords. The author(s) should pick words that accurately describe
%% the work being presented. Separate the keywords with commas.
\keywords{Multi-Task Learning, Multi-Objective Optimization, Merit Function}
%% A "teaser" image appears between the author and affiliation
%% information and the body of the document, and typically spans the
%% page.
% \begin{teaserfigure}
%   \includegraphics[width=\textwidth]{sampleteaser}
%   \caption{Seattle Mariners at Spring Training, 2010.}
%   \Description{Enjoying the baseball game from the third-base
%   seats. Ichiro Suzuki preparing to bat.}
%   \label{fig:teaser}
% \end{teaserfigure}

% \received{20 February 2007}
% \received[revised]{12 March 2009}
% \received[accepted]{5 June 2009}

%%
%% This command processes the author and affiliation and title
%% information and builds the first part of the formatted document.
\maketitle
\newcommand\kddavailabilityurl{https://doi.org/10.5281/zenodo.20377485}
\ifdefempty{\kddavailabilityurl}{}{
\begingroup\small\noindent\raggedright\textbf{Resource Availability:}\\
The source code of this paper has been made publicly available at
\url{\kddavailabilityurl}. The corresponding GitHub repository is available at
\url{https://github.com/Chen-zb/SIMS}.
\endgroup
}
\section{Introduction}

Multi-task learning (MTL) aims to develop a model capable of learning from multiple related tasks simultaneously~\cite{caruana1997multitask,zhang2021survey}. By leveraging shared representations across tasks, MTL enhances \edit{data efficiency and generalization~\cite{caruana1997multitask,ruder2017overview}}
% \footnote{***cy: add reference. check robustness. \edit{delete robustness and have added}}
, and has established itself as a standard paradigm in applications such as dense prediction\edit{~\cite{vandenhende2021multi}}
% \footnote{***cy: add ref. \edit{Done}}
, \edit{autonomous driving~\cite{chowdhuri2019multinet,wang2025review} and recommendation systems~\cite{zhao2019recommending}}
%\footnote{***cy: add ref. \edit{Done}}
. However, tasks in MTL are rarely perfectly aligned. The optimization of one task may degrade the performance of another because the learning process often requires navigating unavoidable trade-offs among competing objectives.

A significant body of research addresses this challenge through architectural design~\cite{liu2019end,misra2016cross}, loss weighting~\cite{kendall2018multi,vandenhende2021multi}, or gradient balancing~\cite{chen2018gradnorm,guo2018dynamic}\cyfoot{***cy: Correlate the references with architectural design, loss weighting, or gradient balance. \edit{Done}}. While these methods are effective in many applications, they are typically heuristic, relying on optimizing a proxy objective that \edit{lacks theoretical guarantees for converging to the optimal solution}\footnote{***at current position, MOO is not mentioned. Suddenly mentioning `Pareto optimal solutions' seems a bit strange. \edit{Put it in another way}} when tasks are competing.
% and introduce additional design choices or hyperparameters that may be difficult to justify from a unified optimality principle.
To address these limitations, recent efforts have sought to 
%formalize the competition between tasks.
%A natural way is to 
%formalize the trade-offs between tasks and thus 
cast the objective function of MTL as a multi-objective optimization (MOO) problem, where each task loss serves as an individual objective with the goal of reaching \edit{Pareto optimal solutions~\cite{sener2018multi,marler2004survey}}\cyfoot{***cy: check `compromise'. Maybe `finding the Pareto optimal solutions'? \edit{have changed to "the Pareto optimal solutions"}}\cyfoot{***cy: add ref. \edit{Done}}. MOO-based MTL methods are generally categorized into the \edit{adaptive gradient} approach~\cite{sener2018multi} and scalarization approach~\cite{lin2024smooth}.  \edit{Unlike adaptive gradient approach, scalarization approach}\cyfoot{***cy: Scalarization-based approaches? \edit{Done, all change to scalarization approaches}} is particularly appealing because it transforms the MOO problem into standard single-objective training by optimizing an aggregated objective, \edit{thereby avoiding the high computational and memory overheads associated with manipulating per-task gradients.}\footnote{***this reason cannot demonstrate that the scalarization approach is appealing. \edit{The advantages of the scalarization approaches over the adaptive gradient approach have been added}}

In this paper, we focus on the merit-function-based\cyfoot{***cy: `merit-function based' or `merit-based'. Be consistent. \edit{merit-function-based}} scalarization approach, which eliminates the need for manually tuning task weighting while maintaining a principled connection to the Pareto optimality. However, in practice, the scales of task losses in various tasks often differ by orders of magnitude. Based on a motivating example shown in Section \ref{sec:motivation}, we observe that existing merit-function-based scalarization methods are often sensitive to the scales of task losses.
%To illustrate the impact of such scale gaps, we study a two-task toy example in Fig.~\ref{fig:toy-scale}.
% and observe that existing merit-based scalarizations are often scale-sensitive.
Although the change of scales of objectives via the rescaling \edit{(i.e., multiplying each objective by a positive constant, akin to a change of units)} does not change the Pareto optimal solutions, 
% rescaling can change optimization trajectories and systematically bias the obtained Pareto solution.
\edit{it can substantially affect the optimization of scalarization methods. Specifically, the rescaling changes the relative magnitudes of objectives and also gradients, which can alter the optimization trajectory and may lead to unsatisfactory solutions,  
%steer convergence toward extreme trade-off solutions, meaning Pareto solutions 
which strongly favor only a few objectives while sacrificing others, with worse average performance across tasks. %Such biased solutions often exhibit worse average performance across tasks, for example when evaluated by the mean of per-task metrics. 
Moreover, we find that the change of the scales of objective functions may even make the optimization procedure divergent. }\cyfoot{***cy: what is positive rescaling? Give a definition. \edit{Done}}\cyfoot{***cy: More explanations of the phenomena in the toy example. \edit{have added more details}}\cyfoot{***cy: why bias the obtained Pareto point is bad. Explain it. \edit{I conclude that it will lead to extreme trade-off solutions}}

To alleviate the above issues in scalarization methods and achieve the scale-invariance in objective functions, we propose the \textbf{S}cale-\textbf{I}nvariant \textbf{M}erit-function-based \textbf{S}calarization (SIMS) method, which \edit{preserves a principled characterization of weak Pareto optimality, }\cyfoot{***cy: check `Pareto-relevant'. \edit{change to "preserves a principled characterization of weak Pareto optimality"}}\footnote{***our method cannot find all the Pareto-optimal solutions. So here we cannot say `the full Pareto set'. Moreover, `Pareto set' means `Pareto front'? \edit{have changed to "preserves a principled characterization of weak Pareto optimality"}} while remaining intrinsically \edit{invariant to the scales of task losses}.\cyfoot{***cy: Change all `robust' to another word in this paper. \edit{change to "invariant to task-wise loss scales"}} 
Specifically, in SIMS, we \edit{propose} a \edit{transformation}-induced\cyfoot{***cy: `Deformation' is mostly used to refer to physical deformation. \edit{change to transformation}} merit function, \edit{which} applies a monotonically increasing transformation function to each task loss. Theoretically, we prove that, to achieve the scale-invariance, the transformation function in SIMS should be logarithmic functions. We further provide theoretical analysis to show that the transformation-induced merit function in SIMS maintains a valid characterization of weak Pareto optimality and admits a smooth surrogate with controllable approximation error.
% \edit{Furthermore, we prove that scale invariance uniquely determines the transformation to be logarithmic}. 
%Thus, SIMS achieves consistent optimization regardless of the initial task scales\cyfoot{***cy: You can modify this sentence if it doesn't suit. \edit{Done}}, \edit{thereby avoiding extreme trade-off solutions}. 
Empirically, extensive experiments demonstrate that the proposed SIMS method consistently outperforms existing scalarization approaches and achieves state-of-the-art performance.
% These findings motivate us to develop a merit-function-based scalarization that (i) preserves Pareto-relevant guarantees and (ii) is intrinsically robust to task-wise rescaling. Specifically, we derive a deformation-induced merit function by applying a monotone deformation to each task loss, and we provide a theoretical analysis showing that this transformed formulation still yields a valid characterization of weak Pareto optimality. Building on this result, by further exploiting the affine invariance of the deformation under positive rescaling, we obtain \textbf{Scale-Invariant Merit-based Scalarization (SIMS)}.
% Empirically, we evaluate SIMS on several representative multi-task benchmarks and show that it consistently improves overall performance over existing scalarization baselines, achieving state-of-the-art results.

The main contributions of this work are summarized as follows.
\begin{itemize}
% \item We empirically demonstrate that existing merit-function-based scalarizations can be strongly scale-sensitive, and show via a toy example that scale invariance is necessary to obtain robust trade-off solutions.
\item We empirically demonstrate that existing \edit{merit-function-based }scalarization methods are often sensitive to the scales of task losses and show that scale-invariance\cyfoot{***cy: Is `scale invariance' a general term? Otherwise, you need to define it. \edit{will define it on Motivating examples}} is essential for obtaining consistent solutions.\footnote{***we are the first to make such observation for scalarization methods? If not, then do not list as a contribution. \edit{Added the attributive of the scalarization methods}}
\item We propose the SIMS method, a scalarization approach which can \edit{characterize weak Pareto optimality while ensuring the scale-invariant property}.%\cyfoot{***cy: Describe the method in one sentence. \edit{Done}}, thereby \edit{avoiding extreme trade-off solutions.}\cyfoot{***cy: Describe the efficacy of the method in one sentence. \edit{Done}}
\item We theoretically analyze the proposed SIMS method, including the formulation of the transformation function, the weak Pareto optimality, and the approximation analysis.% prove that the requirement for scale invariance uniquely determines the transformation to be logarithmic. We further demonstrate that the general transformation-induced merit function, derived from any strictly increasing mapping, preserves weak Pareto optimality and admits a smooth surrogate with controllable approximation error.
\item Extensive experiments on multi-task benchmarks demonstrate the effectiveness of the proposed SIMS method.% as it consistently outperforms state-of-the-art baselines.
% \item Building on this result, we propose SIMS, a scale- merit scalarization, and show it achieves state-of-the-art performance on standard multi-task learning benchmarks.
\end{itemize}



\section{Related work}
% \subsection{Loss Balancing in MTL}
% Multi-task learning often suffers from task imbalance, where naively summing task losses lets one task dominate due to scale disparities~\cite{kirchdorfer2024analytical}. A rich line of work addresses this by dynamically weighting or normalizing losses to stabilize training. For example, uncertainty-based weighting learns task-specific homoscedastic uncertainties to down-weight noisy or inherently harder tasks~\cite{kendall2018multi}, dynamic weight averaging (DWA) assigns higher weight to tasks with slower loss decrease~\cite{liu2019end}, and GradNorm balances training rates by equalizing gradient magnitudes~\cite{chen2018gradnorm}. Other strategies include geometric mean loss formulations~\cite{chennupati2019multinet++}, task difficulty-based prioritization~\cite{guo2018dynamic}, or learnable loss rescaling~\cite{liu2021towards}. While these loss-level methods can improve multitask optimization, they are largely heuristic or data-driven, lacking a unified theoretical interpretation, and their relative efficacy often depends on careful tuning~\cite{xin2022current}. 
% On the other hand, gradient-based balancing approaches explicitly embrace a multi-objective optimization perspective, which tackling task interference by directly manipulating gradients to find update directions that jointly improve multiple tasks, such as 
% Multiple-Gradient Descent Algorithm (MGDA)~\cite{sener2018multi} and many methods that adjust gradients to avoid conflicts~\cite{yu2020gradient,liu2021conflict,chen2018gradnorm}. These methods typically rely on per-task gradient computations and additional projection or optimization steps, which introduce substantial computational overhead that scales unfavorably with the number of tasks, limiting their practicality in large-scale multi-task settings~\cite{kurin2022defense}. Such limitations have in turn motivated research to cast MTL more systematically as a MOO problem and to develop loss-level balancing mechanisms, aiming to retain principled trade-off handling without the heavy per-task gradient machinery.

% \subsection{MOO in MTL}


% Multi-task learning can be naturally formulated as a vector optimization problem in which each task loss is treated as an objective, and the goal is to obtain solutions that are Pareto optimal with respect to task-wise improvements. This viewpoint connects MTL to classical multi-objective optimization, where Pareto optimality and scalarization have been extensively studied as fundamental concepts and tools \citep{miettinen1999nonlinear,marler2004survey,marler2009weightedsum}.

% Early work that explicitly promoted this connection in modern MTL cast joint training as multi-objective optimization and adapted gradient-based MOO algorithms to compute update directions that satisfy common descent conditions \citep{sener2018mtlmoo}. Building on this formulation, subsequent research explored how to obtain either a single Pareto solution or a set of trade-off solutions. ParetoMTL decomposes the MTL problem into multiple constrained subproblems associated with different trade-off preferences, yielding a collection of well-distributed Pareto solutions \citep{lin2019paretomtl}. In parallel, a line of methods focused on gradient-level updates that reconcile task objectives during training. PCGrad proposed a projection-based procedure that modifies per-task gradients when conflicts arise \citep{yu2020pcgrad}. CAGrad introduced a conflict-averse update rule that regularizes the training trajectory while maintaining convergence properties under its formulation \citep{liu2021cagrad}.

% Beyond computing a single compromise point, several works aim to parameterize or approximate a continuum of Pareto trade-offs within a single training framework. Pareto HyperNetworks introduced a hypernetwork-based mechanism that maps a preference vector to a corresponding model on the Pareto front, enabling post hoc selection of operating points \citep{navon2021phn}. Related efforts studied controllable or switchable trade-offs within one system for multi-task settings \citep{ma2020controllablepareto}. Game-theoretic and fairness-inspired formulations further enriched the MOO perspective on MTL. Nash-MTL interpreted gradient aggregation as a bargaining game and derived an update procedure from the Nash bargaining solution \citep{navon2022bargaining}. Impartial MTL studied criteria that encourage balanced training progress across tasks by controlling gradient and loss scales \citep{liu2021impartial}. More recently, the role of scalarization in MTL has been revisited from empirical and theoretical angles, including analyses that question when sophisticated multi-objective optimizers are necessary and studies that characterize the limitations and capabilities of scalarization for exploring trade-offs \citep{kurin2022unitary,wang2023revisitscalarization}.

% \subsection{Merit-based Optimization in MTL}
% In multi-objective optimization, a merit function maps multiple objectives to a single scalar value or criterion that can be optimized~\cite{marler2004survey}. Classical examples include weighted sums of task losses, weighted Chebyshev norms, and $\epsilon$-constraint methods that treat some objectives as constraints~\cite{miettinen1999nonlinear,mahapatra2020multi}


% Merit-based optimization casts MTL as multi-objective optimization and chooses a merit function that induces either a scalar training objective or a descent criterion over task gradients~\cite{marler2004survey}. In principle, this design space is broad, ranging from classical scalarizations (weighted sums, Tchebycheff norms, $\epsilon$-constraint) to modern welfare- and preference-driven criteria that target specific Pareto trade-offs\footnote{***}~\cite{miettinen1999nonlinear,mahapatra2020multi,navon2022multi}. In deep MTL\footnote{***}, however, the key bottleneck is not expressiveness but usability, as the merit must be differentiable and first-order tractable so that it can be optimized efficiently with SGD~\cite{chen2025gradient,peitz2025multi}. This motivates a recent line of deep-friendly\footnote{***} merit formulations that explicitly engineer\footnote{***} smooth or penalized surrogates with easy-to-compute gradients. STCH~\cite{lin2024smooth} smooths the classical Tchebycheff scalarization into a lightweight surrogate suitable for gradient-based MOO, while FOOPS~\cite{chen2025efficient} turns Pareto constraints into a single-objective penalty whose gradient is straightforward to evaluate. These smooth merit formulations are attractive for deep MTL because they yield a single coherent training signal compatible with standard gradient-based training. However, they often rely on extra optimization signals\footnote{***} (e.g., temperatures, reference points, preferences) to shape this update, which can interact with loss rescaling and make training scale-sensitive; achieving better Pareto trade-offs therefore calls for scale-robust merit design.

% Merit functions were originally developed as a bridge from structured equilibrium problems to unconstrained scalar minimization. In variational inequalities and complementarity problems, gap functions and related merit constructions yield nonnegative scalars whose zero level sets coincide with the solution set, enabling optimization-based solution procedures \citep{auchmuty1989gap,facchinei2003vi}. This concept has been extended to multi-objective optimization, where merit or gap-type functions can be designed to characterize weak Pareto optimality through a single scalar objective while retaining a direct relationship to Pareto solution sets \citep{tanabe2023meritmoo}.

% In multi-objective learning and MTL, merit-based ideas appear in formulations that impose Pareto optimality as a constraint and convert it into an optimizable objective through merit constructions and penalty reformulations. One representative direction studies optimization on the Pareto set under preference guidance and develops merit functions with tractable gradients to enable first-order algorithms \citep{chen2025opsmerit}. In addition to explicitly merit-constrained formulations, recent work has proposed smooth and first-order compatible scalarization mechanisms that serve a similar operational role in training pipelines. Smooth Tchebycheff scalarization provides a differentiable approximation to classical Tchebycheff-style constructions and is used as a practical scalar objective for gradient-based multi-objective learning, including multi-task settings \citep{lin2024stch}. Taken together, these developments reflect a progression from classical merit and gap functions in optimization theory to differentiable, training-compatible merit or merit-like scalar objectives used in contemporary multi-task learning.

% If the content exceeds the page limit, use \noindent \textbf{MTL as MOO.}
% If it is less than the page limit, use \subsection{MTL as MOO}
\noindent \textbf{MOO for MTL.} 
%To develop principled MTL methods with provable theoretical guarantees, MTL can be naturally formulated as an MOO problem, where 
To formulate the objective function of MTL as an MOO problem, each task loss corresponds to an objective. To better handle task conflicts, the adaptive gradient approach has been widely studied to construct a descent direction that balances multiple gradients of tasks. For example, MGDA \cite{sener2018multi} finds a common descent direction to satisfy the Pareto stationarity \citep{desideri2012multiple}.\footnote{***add references \edit{Done}} Subsequent methods such as CAGrad \cite{liu2021conflict}, PCGrad \cite{yu2020gradient}, Nash-MTL \cite{navon2022multi}, and FairGrad~\cite{ban2024fair} improve the stability or fairness via respective gradient aggregation operations. Although these adaptive gradient methods provide strong theoretical guarantees, they typically require accessing or manipulating per-task gradients and storing them during optimization, which leads to high computational and memory overhead. 
Due to these issues, we focus on the scalarization approach that combines multiple task losses to a single objective.

\vspace{0.5\baselineskip}
\noindent \textbf{Scalarization approach.}
In MOO, the scalarization approach converts multiple objectives to a single scalar objective to be optimized~\cite{marler2004survey}. 
% Classical approaches include linear weighted sum, weighted Chebyshev norms, and $\epsilon$-constraint methods that treat some objectives as constraints~\cite{miettinen1999nonlinear,mahapatra2020multi}. 
A substantial body of prior work in MTL is built upon the linear scalarization, where different task losses are linearly combined to form the objective function. %multi-task objective is expressed as a weighted combination of individual task losses. 
%Importantly, although the task weights in these methods may change dynamically during training, the underlying objective remains a linear aggregation. 
Representative methods under this paradigm include Equal Weight (EW) with an identical task weight, Uncertainty Weighting (UW)~\cite{kendall2018multi}, Dynamic Weight Averaging (DWA)~\cite{liu2019end}, and GradNorm~\cite{chen2018gradnorm}. 
However, the linear scalarization could miss all the solutions on the non-convex part of the Pareto front~\cite{das1997closer}. 
% To solve this, Weighted Tchebycheff (TCH) scalarization is often used as it can theoretically identify non-convex solutions. 
To alleviate this issue, the Smooth Tchebycheff (STCH) method~\cite{lin2024smooth} introduces a smooth formulation to resolve the non-differentiability of the classical weighted Tchebycheff method, which enables efficient gradient-based optimization, while preserving the theoretical capability to capture non-convex Pareto fronts. Moreover, the FOOPS method~\cite{chen2025efficient} leverages a merit function for MTL to provide an explicit scalar measure of the Pareto optimality. %, thereby offering a clearer characterization of the entire Pareto set without relying on linear scalarization. 
\zbedit{Some recent methods have also considered scale mismatch in MTL. For example, GLS \cite{chennupati2019multinet++} reduces the influence of loss magnitudes through geometric loss, but it does not explicitly characterize whether the current solution admits further joint improvement or reaches Pareto optimality. Other methods, such as Nash-MTL \cite{navon2022multi} and DB-MTL \cite{lin2025dual} address related issues from the gradient-level perspective.}
Unlike these methods, the proposed SIMS method aims to solve the scale-sensitive issue of existing merit functions.%we propose a Scale-Invariant Merit-function-based Scalarization for MTL, which preserves the ability to characterize the full Pareto set while avoiding extreme trade-off solutions \edit{driven by differing orders of magnitude in task losses}\cyfoot{***cy: change `cross-task loss-scale disparities'. \edit{Done}}.

\section{Preliminary}
\subsection{MTL as MOO}
An MTL problem involves a set of $m$ related tasks and can be formulated as an MOO problem as
\begin{equation}
    {\min}_{\theta\in \mathbb{R}^d} \; \boldsymbol{\mathcal{L}}(\theta)
    = \bigl( \mathcal{L}_1(\theta), \mathcal{L}_2(\theta), \ldots, \mathcal{L}_m(\theta) \bigr)^\top,
    \label{eq:1}
\end{equation}
where $\mathcal{L}_i(\theta):\mathbb{R}^d\to\mathbb{R}_+$ denotes the task loss of the $i$-th task and $\theta$ denotes the shared model parameters across tasks. %An MTL algorithm aims to optimize all tasks simultaneously by exploiting shared representations and information across tasks.
For problem~\eqref{eq:1}, there is usually no single solution that can simultaneously minimize all objectives. Instead, we consider the Pareto optimality for problem~\eqref{eq:1}, which is defined as follows.
\begin{definition}[Pareto optimality]\label{def:pareto_opt}
A solution $\theta^\star \in \mathbb{R}^d$ is \emph{Pareto optimal} to problem~\eqref{eq:1}, if there does not exist any other solution $\theta \neq \theta^\star$ such that $\mathcal{L}_i(\theta) \le \mathcal{L}_i(\theta^\star)$ for all $i \in [m]$ and $\mathcal{L}_j(\theta) < \mathcal{L}_j(\theta^\star)$ for at least one index $j$.
\end{definition}

%A \edit{solution}\footnote{***point changes to solution \edit{Done}} is Pareto optimal if it cannot be improved in any objective without sacrificing another. 
\edit{The set of all Pareto optimal solutions is called the Pareto set. The image of the Pareto set in the loss space is called the Pareto front.} We also consider a weaker notion of Pareto optimality, namely weak Pareto optimality.\footnote{***briefly introduce the Pareto front, which is already used in this paper.\edit{Done}}

\begin{definition}[Weak Pareto optimality]\label{def:weak_pareto_opt}
A solution $\theta^\star \in \mathbb{R}^d$ is \emph{weakly Pareto optimal} to problem~\eqref{eq:1}, if there does not exist any other solution $\theta \in \mathbb{R}^d$ such that $\mathcal{L}_i(\theta) < \mathcal{L}_i(\theta^\star)$ for all $i \in [m]$.
\end{definition}

% A \edit{solution}\footnote{***} is weakly Pareto optimal, if there is no other point that improves all of the objective functions simultaneously. Obviously, every Pareto optimal \edit{solution}\footnote{***} is weakly Pareto optimal.
% \begin{definition}[Pareto Dominance]
% Let $\theta^a, \theta^b \in \mathbb{R}^d$ be two parameter vectors.
% We say that $\theta^a$ \emph{dominates} $\theta^b$, denoted by
% $\theta^a \prec \theta^b$, if and only if $\mathcal{L}_i(\theta^a) \le \mathcal{L}_i(\theta^b), \forall i \in [m],$ and there exists $j \in [m]$ such that $\mathcal{L}_j(\theta^a) < \mathcal{L}_j(\theta^b)$. 
% % Moreover, $\theta^a$ is said to \emph{strictly dominate} $\theta^b$, denoted by
% % $\theta^a \prec_{\mathrm{strict}} \theta^b$, if and only if $\mathcal{L}_i(\theta^a) < \mathcal{L}_i(\theta^b), \forall i \in [m]$.
% \end{definition}

% \begin{definition}[(Weakly) Pareto Optimality]
% A point $\theta^\star \in \mathbb{R}^d$ is called \emph{Pareto optimal} if there does not exist
% $\hat{\theta} \in \mathbb{R}^d$ such that $\hat{\theta} \prec \theta^\star$.
% A point $\theta^\star \in \mathbb{R}^d$ is called \emph{weakly Pareto optimal} if there does not exist
% $\hat{\theta} \in \mathbb{R}^d$ such that $\mathcal{L}_i(\hat{\theta}) < \mathcal{L}_i(\theta^\star),\forall i \in [m]$.
% \end{definition}

% In this paper, we focus on identifying a set of representative Pareto solutions that effectively approximate the Pareto front.  
% This idea, along with a comparison between the proposed method and two existing approaches, is illustrated in Figure~2.

\begin{figure*}[!hpbt]
  \centering
  \includegraphics[width=0.97\linewidth]{figs/toy_5.pdf}
  \caption{The visualization of optimization trajectories of different methods in the loss space. Gray: Pareto front. Black: initializations. Colored curves: training trajectories. Under both objective-scale settings (Scale 10:1 and 1:10), the Pareto solution set is unchanged. EW, STCH, and FOOPS show clear sensitivity to objective scaling. \zbedit{GLS may get trapped in local optima and converge away from the Pareto front}. In contrast, SIMS (Ours) behaves consistently across scales, producing similar trajectories and converging to comparable regions.
  }
  \label{fig:toy-scale}
\end{figure*}

\subsection{Merit Functions}
Merit functions~\cite{patriksson1997merit,charitha2013note,pappalardo2016merit} provide a surrogate function that connects MOO with single-objective optimization. Merit functions are defined as follows.
\begin{definition}[Merit Function]
\label{def_merit_function}
A function $p:\mathbb{R}^d\to\mathbb{R}$ is a merit function for problem~\eqref{eq:1}
if it satisfies the following properties:
\begin{itemize}
    \item $p(\theta)\ge 0$ for all $\theta\in\mathbb{R}^d$;
    \item $\theta^\star$ is a weakly Pareto optimal solution of problem~\eqref{eq:1} if and only if $p(\theta^\star)=0$.
\end{itemize}
\end{definition}

%Such functions are designed to be nonnegative over the feasible set and equal to zero if and only if the point is a solution\footnote{***Delete gap function~\edit{Done}}. This reformulation enables solving the original problem by globally minimizing the merit function, and has been extensively studied due to desirable properties such as continuity and differentiability under suitable conditions~\cite{facchinei1997new,pappalardo2016merit}.

%More recently, merit-function-based formulations have also been extended to MOO and vector optimization, where 
 
%We now give a formal definition of merit function in MOO as follows~\cite {pappalardo2016merit}.



%\footnote{***introduce some examples of merit function.\edit{Done}}

%We give some examples of merit functions for MOO. 
A simple example of the merit function is the maximum-loss merit function~\cite{miettinen1999nonlinear}, which is defined as
\[
p_{\max}(\theta)
=
\max_{i\in[m]} \big( \mathcal{L}_i(\theta) - \mathcal{L}_i^{\inf} \big),
\]
where $[m]$ denotes the set of integers from 1 to $m$, and $\mathcal{L}_i^{\inf}= \inf_{\theta\in\mathbb{R}^d} \mathcal{L}_i(\theta)$ denote the minimal attainable value of the $i$-th objective. 
%This function measures the worst deviation of the objective vector from the ideal loss vector. 
We can prove that the maximum-loss merit function satisfies Definition \ref{def_merit_function}. %It is always nonnegative and vanishes if and only if $\theta$ is weakly Pareto optimal. 
Moreover, a weighted extension of the maximum-loss merit function leads to the Chebyshev merit function, which is used in MOO~\cite{ehrgott2005multicriteria,lin2024smooth}. Specifically, the Chebyshev merit function\footnote{***Chebyshev-type merit function or Chebyshev merit function? to be consistent. \edit{Done}} is defined as
\[
p_{\mathrm{CH}}(\theta)
=
\max_{i\in[m]} w_i \big( \mathcal{L}_i(\theta) - \mathcal{L}_i^{\inf} \big).
\]
Compared with $p_{\max}(\theta)$, $p_{\mathrm{CH}}(\theta)$ preserves the properties of the merit function, while allowing different objectives to be emphasized through the weights.

%they offer a \emph{weight-free} characterization of Pareto optimality~\cite{tanabe2024new,chen2025efficient}. This perspective is particularly appealing in MTL, in which task weights are often difficult to specify and may vary across datasets or training stages.

%In this paper, we will adopt a smoothed merit function as a surrogate objective for problem~\eqref{eq:1}. 
%Since the resulting problem is generally nonconvex, our goal is to compute a stationary point of the surrogate function, defined as follows.
To analyze the merit function, we need the stationary property, which is defined as follows.

\begin{definition}
A solution $\theta$ is an $\epsilon$-stationary point ($\epsilon \ge 0$) of a differentiable function $p$ if its gradient $\nabla p(\theta)$ satisfies $\|\nabla p(\theta)\| \le \epsilon$, where $\|\cdot\|$ is the Euclidean norm. If $\epsilon = 0$, then $\theta$ is a stationary point.
\end{definition}


% Below we give two examples of merit functions for MOO. A classical example is the Chebyshev-type merit function, derived from the Chebyshev scalarization in MOO~\cite{miettinen1999nonlinear}. Let $\mathcal{L}_i^{\inf} := \inf_{\theta} \mathcal{L}_i(\theta)$ denote the minimal attainable value of the $i$-th objective. The Chebyshev merit function is defined as
% \[
% p_{\mathrm{CH}}(\theta) := \max_{i\in[m]} \big( \mathcal{L}_i(\theta) - \mathcal{L}_i^{\inf} \big),
% \]
% which measures the worst deviation of the objective vector from the ideal loss vector. This merit function vanishes if and only if $\theta$ is weakly Pareto optimal, and is strictly positive otherwise.

% A smooth counterpart can be obtained using the log-sum-exp smoothing technique, leading to the smoothed Chebyshev merit function
% \[
% p_{\mathrm{SCH}}(\theta) := \tau \ln \left( \sum_{i=1}^m 
% \exp\!\left( \frac{\mathcal{L}_i(\theta) - \mathcal{L}_i^{\inf}}{\tau} \right) \right), 
% \quad \tau > 0,
% \]
% which provides a differentiable approximation of $p_{\mathrm{CH}}(\theta)$ and preserves the merit function property.



% Below we give two examples of merit functions for MOO. A simple example is the maximum-loss merit function
% \[
% p_{\max}(\theta) := \max_{i\in[m]} \mathcal{L}_i(\theta),
% \]
% which focuses on the worst-performing objective and encourages balanced improvement across tasks. A smooth counterpart can be obtained via the log-sum-exp smoothing technique, leading to the merit function \[
% p_{\mathrm{LSE}}(\theta)
% := \tau \ln\!\left(\sum_{i=1}^m
% \exp\!\left(\frac{\mathcal{L}_i(\theta)}{\tau}\right)\right), \quad \tau > 0,
% \]
% which is amenable to gradient-based methods. Both functions map multiple objectives to a single scalar while emphasizing the improvement of poorly performing objectives, making them suitable surrogates for MOO.


% \paragraph{Why merit functions?}
% Compared with scalarization-based approaches, merit functions offer two key advantages.
% First, they avoid explicitly specifying task or objective weights, which may be subjective or unstable across training stages.
% Second, they enable direct theoretical analysis by translating Pareto optimality into properties of a scalar function (e.g., nonnegativity, zero-value characterization, and error bounds), thereby allowing one to leverage standard tools from smooth and nonsmooth optimization~\cite{Nesterov2005,BeckTeboulle2012}. 

% \subsection{Loss Pre-normalization}
% In multi-task learning, task-specific losses often differ substantially in scale due to heterogeneous task types (e.g., classification vs.\ regression) or targets measured in distinct units. Such scale disparities distort the geometry of the loss space and may severely bias subsequent merit-based aggregation or task-balancing mechanisms. Therefore, previous works usually introduce a \emph{loss pre-normalization} step, which can be interpreted as a \emph{reparameterization of the loss space} prior to optimization. Formally, it defines a normalized loss
% \begin{equation}
%     \tilde{\mathcal{L}}_i(\theta) = \psi \left(\mathcal{L}_i(\theta)\right),
% \end{equation}
% where $\psi : \mathbb{R}_{+} \rightarrow \mathbb{R}$ is a monotone and order-preserving transformation designed to align loss scales across tasks. The normalized losses $\{\tilde{\mathcal{L}}_i\}$ are then used as
% inputs to downstream task-balancing approaches.

% In literature, there are mainly two types of loss pre-normalization methods. The first is called \emph{rescaled normalization}, which normalizes each task loss by its initial value:
% \begin{equation}
%     \tilde{\mathcal{L}}_i(\theta) = \frac{\mathcal{L}_i(\theta)}{\mathcal{L}_i(\theta_0)},
% \end{equation}
% where $\mathcal{L}_i(\theta_0)$ denotes the initial loss of task $i$ at $\theta_0$. This normalization anchors all task losses to a shared reference point at initialization, allowing $\tilde{\mathcal{L}}_i(\theta)$ to reflect the \emph{relative training progress} of each task rather than its absolute magnitude. Rescaled normalization is simple to implement and is effective when task losses are of comparable scale and differ mainly by constant multiplicative factors. For more severe scale heterogeneity, we usually consider
% \emph{logarithmic normalization}, defined as
% \begin{equation}
%     \tilde{\mathcal{L}}_i(\theta) = \ln\left(\mathcal{L}_i(\theta)\right).
% \end{equation}
% From a mathematical perspective, logarithmic normalization maps multiplicative scale differences into additive offsets, thereby compressing the loss dynamic range and mitigating scale dominance. 


\section{Motivating Examples}
\label{sec:motivation}

% In multi-task learning (MTL), losses from different tasks are rarely commensurate.
% Magnitude discrepancies can arise from task heterogeneity (e.g., difficulty and noise), label/output spaces (e.g., regression vs.\ classification), and loss definitions or reduction conventions (e.g., sum vs.\ mean, normalization by sequence length/pixels/anchors), among other implementation details.
% As a result, aggregation rules that operate on absolute loss values may implicitly favor large-scale tasks, even when all tasks are intended to be treated equally. \textcolor{blue}{A: add some citation. when you say some phenomenon, results or not very obvious facts, you should cite some papers to support it.}
% In multi-task learning, the loss terms associated with different tasks are often not directly comparable in scale. \textcolor{green}{we don't consider noise and uncentainty, so we would better not mention that.}Such discrepancies can stem from intrinsic task heterogeneity, e.g., varying noise levels and uncertainty—as well as from differences in output spaces and physical units when mixing regression and classification objectives~\cite{kendall2018multi}. In practice, implementation choices further affect loss magnitudes, including the reduction convention (e.g., mean vs. sum) and task-dependent normalization (e.g., averaging over pixels, anchors, or sampled RoIs)~\cite{ren2015faster,he2017mask}. Consequently, aggregation strategies that operate on absolute loss values can implicitly prioritize large-scale objectives,\footnote{***} prior work~\cite{chen2018gradnorm} shows that naive equal-weight summation may cause training to be dominated by tasks with larger loss scales, undermining the intended equal treatment across task.\footnote{***simplify this paragraph}
%\footnote{***in this section, you need to clarify why merit-function-based scalarization methods need the scale-invariant property.}
In MTL, the task losses across tasks are often in different scales. This phenomenon can arise from heterogeneous loss functions and distinct types of tasks (e.g., regression versus classification tasks)~\cite{kendall2018multi}, and is further enlarged by implementation choices such as the reduction convention \edit{(e.g., sum versus mean aggregation)}\footnote{***what does reduction convention mean? \edit{Added explanations}} and task-specific normalization~\cite{ren2015faster,he2017mask}. 
Although such scale differences across tasks are common, they \edit{do not theoretically affect}\footnote{***revise this phrase \edit{Done}} the Pareto optimality because any rescaling of objectives will not change the Pareto optimal solutions. 
% Although such scale discrepancies are common, they are inconsequential for dominance in the multi-objective view that any positive rescaling of objectives leaves Pareto dominance unchanged. 
In contrast, scalarization methods are often sensitive to the scales of task losses because rescaling changes the gradient magnitudes across different tasks, which can alter the optimization trajectory to a different solution.  %\cyfoot{***cy: why convergence at a different solution is bad. explain it. \edit{have added}}
%This issue is particularly problematic because Pareto optimality does not necessarily imply balanced performance across tasks. 
When task losses differ by orders of magnitude, \edit{existing merit-function-based }scalarization methods\footnote{***you need to narrow down to merit-function-based scalarization methods. \edit{Done}} can be dominated by objectives with larger scales, pushing the optimization toward Pareto-optimal yet extreme trade-off solutions, where the model favors a few tasks while substantially sacrificing the remaining ones. This often yields unsatisfactory performance for MTL.

% \paragraph{Scale invariance as a desideratum.}
% A basic requirement for robust MTL is \emph{scale invariance}: if each task loss is rescaled by an arbitrary positive constant,
% \(\ell_t(\theta) \mapsto c_t \ell_t(\theta)\) with \(c_t>0\),
% the training objective should not change in any essential way (e.g., should differ by at most an additive constant).
% The logarithm provides a canonical illustration.
% Consider the log-sum scalarization
% \begin{equation}
% m_{\ln}(\theta) = \sum_{t=1}^T w_t \ln(\ell_t(\theta)+\epsilon),
% \label{eq:log_scalar}
% \end{equation}
% where \(\epsilon>0\) is a small constant.
% Under rescaling, \(\ln(c_t \ell_t)=\ln \ell_t + \ln c_t\), hence \(m_{\ln}\) changes only by \(\sum_t w_t \ln c_t\), which does not affect optimization.
% Equivalently, Eq.~\eqref{eq:log_scalar} optimizes a weighted product (geometric mean) of losses and depends on \emph{relative} rather than absolute progress. 

% \paragraph{Scale invariance.}
% Let \(\boldsymbol{\ell}(\theta)=(\ell_1(\theta),\dots,\ell_T(\theta))\in\mathbb{R}_+^T\) be the vector of task losses. Scale imbalance is determined by \emph{relative} task scales, so w.l.o.g.\ we measure rescaling relative to task~1 and fix it as the reference. Specifically, consider
% \begin{equation}
%     (\mathbf{c}\odot\boldsymbol{\ell})(\theta)=(c_1\ell_1(\theta),\,c_2\ell_2(\theta),\,\dots,\,c_T\ell_T(\theta)),
% \end{equation}
% where \(c_1=1\), $c_t>0$.

% \begin{definition}[Strong scale invariance]
% A differentiable scalarization \(m:\mathbb{R}_+^T\to\mathbb{R}\) is \emph{(strongly) scale-invariant} if for any \(\mathbf{c}\in\mathbb{R}_+^T\) with \(c_1=1\), there exists
% \(\kappa(\mathbf{c})\) independent of \(\theta\) such that
% \begin{equation}
% m\!\left(\mathbf{c}\odot\boldsymbol{\ell}(\theta)\right)
% =
% m\!\left(\boldsymbol{\ell}(\theta)\right)
% +\kappa(\mathbf{c})
% \quad\text{for all }\theta.
% \label{eq:strong_scale_inv}
% \end{equation}
% \end{definition}
% Differentiating \eqref{eq:strong_scale_inv} yields \(\nabla_\theta m(\mathbf{c}\odot\boldsymbol{\ell}(\theta))=\nabla_\theta m(\boldsymbol{\ell}(\theta))\), i.e., the induced gradients are scale-free.

% A canonical example satisfying the strong definition is the log-sum scalarization
% $
% m_{\ln}(\theta)=\sum_{t=1}^T w_t \ln \ell_t(\theta),
% $
% since \(m_{\ln}(\mathbf{c}\odot\boldsymbol{\ell})=m_{\ln}(\boldsymbol{\ell})+\sum_t w_t\ln c_t\).

% \begin{definition}[Weak scale invariance]
% A differentiable scalarization \(m\) is \emph{weakly scale-invariant} if for any \(\mathbf{c}\) with \(c_1=1\) there exists a scalar
% \(\alpha(\mathbf{c})>0\), independent of \(\theta\), such that
% \begin{equation}
% \nabla_\theta\, m\!\left(\mathbf{c}\odot\boldsymbol{\ell}(\theta)\right)
% =
% \alpha(\mathbf{c})\,
% \nabla_\theta\, m\!\left(\boldsymbol{\ell}(\theta)\right)
% \quad\text{for all }\theta.
% \label{eq:weak_scale_inv}
% \end{equation}
% In this sense, the closer \(\alpha(\mathbf{c})\) is to \(1\), the less sensitive the induced optimization dynamics are to relative loss rescaling.
% \end{definition}

% A common and simple heuristic for \emph{relative} scale imbalance\cite{dai2023improvable} is to pre-normalize each task loss by a reference value and then apply linear aggregation:
% \begin{equation}
% \tilde{\ell}_t(\theta)=\frac{\ell_t(\theta)}{\ell_t(\theta_{\mathrm{ref}})+\epsilon},\qquad
% m_{\mathrm{LS}}(\theta)=\sum_{t=1}^T w_t\,\tilde{\ell}_t(\theta),
% \label{eq:prenorm}
% \end{equation}
% where \(\theta_{\mathrm{ref}}\) is typically initialization or an early checkpoint, $\mathrm{LS}$ means the linear scalarization method.
% Under a relative rescaling \(\ell_t\mapsto c_t\ell_t\) (with \(c_1=1\)),
% \[
% \tilde{\ell}_t(\theta)\ \mapsto\ \frac{c_t\ell_t(\theta)}{c_t\ell_t(\theta_{\mathrm{ref}})+\epsilon}\ \approx\ \frac{\ell_t(\theta)}{\ell_t(\theta_{\mathrm{ref}})}
% \quad(\epsilon\ \text{small}),
% \]
% so the dependence on \(c_t\) is only \emph{approximately} canceled; equivalently, the induced gradient scaling factor \(\alpha(\mathbf{c})\) in \eqref{eq:weak_scale_inv} is hoped to remain close to \(1\).
% However, the reference \(\ell_t(\theta_{\mathrm{ref}})\) is only a proxy for task scales: as relative scales drift over training (e.g., due to different learning speeds or saturation), this proxy becomes stale and the cancellation can be distorted, causing \(\alpha(\mathbf{c})\) to deviate from \(1\) and potentially reintroducing scale dominance.

% \paragraph{Scale invariance.}
% Let \(\hat{\boldsymbol{\ell}}(\theta)=(\hat\ell_1(\theta),\dots,\hat\ell_T(\theta))\in\mathbb{R}_+^T\) denote the \emph{scale-aligned} (i.e., comparable-unit) task losses.
% To model relative scale mismatch, we assume the observed losses \(\boldsymbol{\ell}(\theta)\) are related to \(\hat{\boldsymbol{\ell}}(\theta)\) by an unknown task-wise rescaling
% \begin{equation}
% \boldsymbol{\ell}(\theta)=\mathbf{c}\odot\hat{\boldsymbol{\ell}}(\theta)
% =\big(c_1\hat\ell_1(\theta),\,c_2\hat\ell_2(\theta),\,\dots,\,c_T\hat\ell_T(\theta)\big),
% \label{eq:scale_model}
% \end{equation}
% where \(c_t>0\). Since only \emph{relative} scales matter, w.l.o.g.\ we fix task~1 as the reference and set \(c_1=1\); then \(c_t\) (for \(t\ge2\)) represents the scale ratio of task \(t\) relative to task~1.

% \begin{definition}[Strong scale invariance]
% A differentiable scalarization \(m:\mathbb{R}_+^T\to\mathbb{R}\) is \emph{scale-invariant} if for any \(\mathbf{c}\in\mathbb{R}_+^T\) with \(c_1=1\), there exists
% \(\kappa(\mathbf{c})\) independent of \(\theta\) such that for all $\theta$,
% \begin{equation}
% m\!\left(\boldsymbol{\ell}(\theta)\right)
% =
% m\!\left(\mathbf{c}\odot\hat{\boldsymbol{\ell}}(\theta)\right)
% =
% m\!\left(\hat{\boldsymbol{\ell}}(\theta)\right)
% +\kappa(\mathbf{c}).
% \label{eq:strong_scale_inv}
% \end{equation}
% \end{definition}
% Differentiating \eqref{eq:strong_scale_inv} yields
% \(\nabla_\theta m(\mathbf{c}\odot\hat{\boldsymbol{\ell}}(\theta))=\nabla_\theta m(\hat{\boldsymbol{\ell}}(\theta))\), i.e., the training signal is invariant to relative rescaling.

% A canonical example is the log-sum scalarization
% \(
% m_{\ln}(\boldsymbol{\ell})=\sum_{t=1}^T w_t \ln \ell_t,
% \)
% Under a relative rescaling $\boldsymbol{\ell}=\mathbf{c}\odot\hat{\boldsymbol{\ell}}$, we have 
% \(
% m_{\ln}(\boldsymbol{\ell})=\sum_t w_t\ln(c_t\hat\ell_t)
% = m_{\ln}(\hat{\boldsymbol{\ell}})+\sum_t w_t\ln c_t
% \)
% and thus \eqref{eq:strong_scale_inv} holds.

% \begin{definition}[Weak scale invariance]
% A differentiable scalarization \(m\) is \emph{weakly scale-invariant} if for any \(\mathbf{c}\) with \(c_1=1\) there exists a scalar
% \(\alpha(\mathbf{c})>0\), independent of \(\theta\), such that for all $\theta$,
% \begin{equation}
% \begin{aligned}
% \nabla_\theta\, m\!\left(\boldsymbol{\ell}(\theta)\right)
% &=
% \nabla_\theta\, m\!\left(\mathbf{c}\odot\hat{\boldsymbol{\ell}}(\theta)\right)\\
% &=
% \alpha(\mathbf{c})\,
% \nabla_\theta\, m\!\left(\hat{\boldsymbol{\ell}}(\theta)\right).
% \label{eq:weak_scale_inv}
% \end{aligned}
% \end{equation}
% In this sense, the closer \(\alpha(\mathbf{c})\) is to \(1\), the less sensitive the induced optimization dynamics are to relative scale mismatch.
% \end{definition}

% A common heuristic for relative scale imbalance~\cite{dai2023improvable} is to estimate the unknown scale ratios \(\mathbf{c}\) from a reference point and then apply linear scalarization:
% \begin{equation}
% \tilde{\ell}_t(\theta)=\frac{\ell_t(\theta)}{\ell_t(\theta_{\mathrm{ref}})+\epsilon},\qquad
% m_{\mathrm{LS}}(\theta)=\sum_{t=1}^T w_t\,\tilde{\ell}_t(\theta),
% \label{eq:prenorm}
% \end{equation}
% where \(\theta_{\mathrm{ref}}\) is typically initialization or an early checkpoint.
% Under the model \eqref{eq:scale_model}, this corresponds (up to a \(\theta\)-independent factor and ignoring \(\epsilon\)) to replacing
% \(\ell_t(\theta)\) by an approximately scale-aligned quantity \(\hat\ell_t(\theta)\) using
% \(
% \ell_t(\theta)/\ell_t(\theta_{\mathrm{ref}})\approx \hat\ell_t(\theta)/\hat\ell_t(\theta_{\mathrm{ref}}),
% \)
% so the dependence on \(\mathbf{c}\) is only \emph{approximately} removed.
% As relative scales drift over training, the reference-based estimate becomes stale, which can distort the implied \(\mathbf{c}\) and drive the effective \(\alpha(\mathbf{c})\) away from \(1\), reintroducing scale dominance.

% A common heuristic for relative scale imbalance~\cite{dai2023improvable} is to first estimate a per-task inverse-scale factor from a reference point,
% \begin{equation}
% a_t \;\triangleq\; \frac{1}{\ell_t(\theta_{\mathrm{ref}})+\epsilon},
% \qquad a_1 \text{ normalized to } 1,
% \label{eq:a_est}
% \end{equation}
% and then apply a weighted linear objective to the resulting scale-corrected losses:
% \begin{equation}
% m_{\mathrm{LS}}(\theta)\;=\;\sum_{t=1}^T w_t\, a_t\, \ell_t(\theta).
% \label{eq:ls_a}
% \end{equation}
% Under the scale model \(\ell_t(\theta)=c_t\,\hat{\ell}_t(\theta)\) with \(c_1=1\), this amounts to using \(a_t\) as an estimate of \(1/c_t\), so that
% \(
% a_t\,\ell_t(\theta)\approx \hat{\ell}_t(\theta)
% \)
% and \(m_{\mathrm{LS}}(\theta)\approx \sum_t w_t\,\hat{\ell}_t(\theta)\).
% However, \(a_t\) is only a heuristic estimate: since \(\ell_t(\theta_{\mathrm{ref}})\) depends on both \(c_t\) and the (unknown) aligned value \(\hat{\ell}_t(\theta_{\mathrm{ref}})\), the inferred scaling can be biased; moreover, as relative scales drift over training, a fixed reference makes \(a_t\) stale. Consequently, the effective gradient scaling \(\alpha(\mathbf{c})\) in \eqref{eq:weak_scale_inv} may deviate from \(1\), and scale dominance can re-emerge.


%%%%%%%%%%%%%%%%%%%
% \begin{definition}[Weak scale invariance (via residual relative scales)]
% Assume the scale model \(\boldsymbol{\ell}(\theta)=\mathbf{c}\odot\hat{\boldsymbol{\ell}}(\theta)\) with \(c_1=1\).
% For a differentiable scalarization \(m:\mathbb{R}_+^T\to\mathbb{R}\), define its induced training signal
% \[
% \nabla_\theta m(\boldsymbol{\ell}(\theta))
% =\sum_{t=1}^T \underbrace{\Big(c_t\,\partial_{\ell_t} m(\mathbf{c}\odot\hat{\boldsymbol{\ell}}(\theta))\Big)}_{\text{effective coefficient on }\nabla_\theta \hat\ell_t}\;\nabla_\theta \hat\ell_t(\theta).
% \]
% Using task~1 as the reference, let
% \[
% \alpha_m(\mathbf{c},\hat{\boldsymbol{\ell}})\;\triangleq\;
% \frac{\partial_{\ell_1} m(\mathbf{c}\odot\hat{\boldsymbol{\ell}})}{\partial_{\ell_1} m(\hat{\boldsymbol{\ell}})},
% \qquad
% a_{m,t}(\mathbf{c},\hat{\boldsymbol{\ell}})\;\triangleq\;
% \frac{c_t\,\partial_{\ell_t} m(\mathbf{c}\odot\hat{\boldsymbol{\ell}})}{\alpha_m(\mathbf{c},\hat{\boldsymbol{\ell}})\,\partial_{\ell_t} m(\hat{\boldsymbol{\ell}})},
% \]
% so that \(a_{m,1}=1\). We call \(m\) \emph{weakly scale-invariant} if the induced \emph{residual relative scales}
% \(\mathbf{a}_m(\mathbf{c},\hat{\boldsymbol{\ell}})\) are closer to \(\mathbf{1}\) than \(\mathbf{c}\), e.g.,
% \begin{equation}
% \big\|\ln \mathbf{a}_m(\mathbf{c},\hat{\boldsymbol{\ell}})\big\|_\infty
% \;\le\;
% \rho\,
% \big\|\ln \mathbf{c}\big\|_\infty,
% \qquad \text{for some }\rho\in[0,1),
% \label{eq:weak_scale_inv_residual}
% \end{equation}
% uniformly over \(\hat{\boldsymbol{\ell}}\) in the regime of interest.
% Smaller \(\rho\) means stronger suppression of relative scale mismatch.
% \end{definition}

%%%%%%%%%%%%%
% \begin{definition}[Exact scale invariance (relative)]
% A differentiable scalarization $m:\mathbb{R}_+^T\to\mathbb{R}$ is
% \emph{(exactly) scale-invariant} w.r.t.\ the relative rescaling model
% $\boldsymbol\ell(\theta)=\mathbf c\odot \hat{\boldsymbol\ell}(\theta)$ (with $c_1=1$)
% if for any such $\mathbf c$ there exist scalars
% $\alpha(\mathbf c)>0$ and $\beta(\mathbf c)$, both independent of $\theta$, such that
% for all $\theta$,
% \begin{equation}
% m\!\left(\mathbf c\odot \hat{\boldsymbol\ell}(\theta)\right)
% =
% \alpha(\mathbf c)\, m\!\left(\hat{\boldsymbol\ell}(\theta)\right)
% +\beta(\mathbf c).
% \label{eq:exact_scale_inv}
% \end{equation}
% \end{definition}

% \paragraph{Implication.}
% Differentiating~\eqref{eq:exact_scale_inv} yields
% $\nabla_\theta m(\mathbf c\odot \hat{\boldsymbol\ell}(\theta))
% =
% \alpha(\mathbf c)\, \nabla_\theta m(\hat{\boldsymbol\ell}(\theta))$,
% so the \emph{update direction} is invariant to relative rescaling (up to a global step-size factor).

% \begin{definition}[$\gamma$-weak scale invariance / scale attenuation]
% Let $\tilde w_t(\theta;\mathbf c):=
% c_t\frac{\partial m}{\partial \ell_t}(\mathbf c\odot \hat{\boldsymbol\ell}(\theta))$
% be the effective coefficient on $\nabla_\theta \hat\ell_t(\theta)$.
% We say $m$ is \emph{$\gamma$-weakly scale-invariant} if there exists $\gamma\in[0,1]$
% such that for any $\mathbf c$ with $c_1=1$ and all $\theta$,
% \begin{equation}
% \mathrm{spread}\!\Big(\tilde{\mathbf w}(\theta;\mathbf c)\oslash \tilde{\mathbf w}(\theta;\mathbf 1)\Big)
% \;\le\;
% \gamma\, \mathrm{spread}(\mathbf c),
% \label{eq:weak_scale_att}
% \end{equation}
% where $\oslash$ denotes elementwise division.
% \end{definition}

% \paragraph{Scale mismatch and trajectory bias.}
% Let $\hat{\boldsymbol{\ell}}(\theta)=(\hat\ell_1(\theta),\dots,\hat\ell_T(\theta))\in\mathbb{R}_+^T$
% denote the \emph{scale-aligned} (comparable-unit) task losses that reflect the intrinsic
% multi-objective trade-offs. In practice, however, each task loss is often defined up to an
% arbitrary positive unit (e.g., different normalization, data cardinality, or loss type),
% so the observed losses can be modeled as an unknown task-wise rescaling:
% \begin{equation}
% \boldsymbol{\ell}(\theta)=\mathbf{c}\odot\hat{\boldsymbol{\ell}}(\theta),
% \qquad c_t>0,\quad c_1=1,
% \label{eq:scale_model}
% \end{equation}
% where $c_t$ (for $t\ge2$) captures the \emph{relative} scale ratio of task $t$ to task~1.

% Although positive rescaling does not change Pareto dominance itself, it can substantially
% distort the optimization \emph{trajectory} induced by a scalarization $m(\boldsymbol{\ell}(\theta))$:
% tasks with larger $c_t$ may dominate the aggregated training signal, steering the iterate toward
% solutions that are artifacts of unit choices rather than reflecting the desired trade-offs on the
% aligned objectives. This motivates a scalarization principle:
% \emph{the induced optimization dynamics should be insensitive to relative rescaling.}

% \paragraph{Effective task influence under a scalarization.}
% For any differentiable scalarization $m:\mathbb{R}_+^T\to\mathbb{R}$, the training signal satisfies
% \begin{equation}
% \nabla_\theta m(\boldsymbol{\ell}(\theta))
% =
% \sum_{t=1}^T
% \frac{\partial m}{\partial \ell_t}\big(\boldsymbol{\ell}(\theta)\big)\,
% \nabla_\theta \ell_t(\theta).
% \label{eq:chain_rule}
% \end{equation}
% Under the rescaling model $\ell_t(\theta)=c_t\hat\ell_t(\theta)$, we have
% $\nabla_\theta \ell_t(\theta)=c_t\nabla_\theta\hat\ell_t(\theta)$, hence
% \begin{equation}
% \nabla_\theta m\!\left(\mathbf{c}\odot\hat{\boldsymbol{\ell}}(\theta)\right)
% =
% \sum_{t=1}^T
% \underbrace{\left(
% c_t\,\frac{\partial m}{\partial \ell_t}\!\left(\mathbf{c}\odot\hat{\boldsymbol{\ell}}(\theta)\right)
% \right)}_{\triangleq~\tilde w_t(\theta;\mathbf{c})}
% \,\nabla_\theta\hat\ell_t(\theta),
% \label{eq:effective_weight}
% \end{equation}
% where $\tilde w_t(\theta;\mathbf{c})$ is the \emph{effective task influence} (coefficient) on the
% aligned gradient $\nabla_\theta\hat\ell_t(\theta)$.

% \paragraph{Scale invariance as scale attenuation.}
% We quantify how much relative rescaling $\mathbf{c}$ perturbs the induced influence pattern
% $\tilde{\mathbf{w}}(\theta;\mathbf{c})=(\tilde w_1,\dots,\tilde w_T)$.
% For multiplicative scale factors, a natural dispersion measure is the log-spread
% \[
% \mathrm{spread}(\mathbf{v})
% ~\triangleq~
% \max_t \ln v_t - \min_t \ln v_t.
% \]
% We compare the perturbation induced by $\mathbf{c}$ to the baseline $\mathbf{1}$ via the ratio
% $\tilde{\mathbf{w}}(\theta;\mathbf{c})\oslash \tilde{\mathbf{w}}(\theta;\mathbf{1})$.

% \begin{definition}[$\gamma$-scale invariance (scale attenuation)]
% A differentiable scalarization $m$ is \emph{$\gamma$-scale-invariant} (w.r.t.~\eqref{eq:scale_model})
% if there exists $\gamma\in[0,1]$ such that for any $\mathbf{c}\in\mathbb{R}_+^T$ with $c_1=1$ and all $\theta$,
% \begin{equation}
% \mathrm{spread}\!\Big(
% \tilde{\mathbf{w}}(\theta;\mathbf{c})\oslash \tilde{\mathbf{w}}(\theta;\mathbf{1})
% \Big)
% ~\le~
% \gamma\,\mathrm{spread}(\mathbf{c}).
% \label{eq:gamma_scale_inv}
% \end{equation}
% \end{definition}

% The parameter $\gamma$ directly captures \emph{how much} the scalarization compresses relative scale mismatch:
% $\gamma=0$ means the influence pattern is invariant to relative rescaling (no trajectory bias from units),
% while $\gamma=1$ means the mismatch is fully inherited (no attenuation).

% \paragraph{Why the log transform is the natural optimal choice.}
% A canonical choice is the log-sum scalarization
% \begin{equation}
% m_{\ln}(\boldsymbol{\ell})
% ~=~
% \sum_{t=1}^T w_t \ln \ell_t,
% \qquad w_t>0.
% \label{eq:mlog}
% \end{equation}
% It yields
% $\frac{\partial m_{\ln}}{\partial \ell_t}=w_t/\ell_t$, and thus from \eqref{eq:effective_weight},
% \[
% \tilde w_t(\theta;\mathbf{c})
% =
% c_t\cdot \frac{w_t}{c_t\hat\ell_t(\theta)}
% =
% \frac{w_t}{\hat\ell_t(\theta)},
% \]
% which is \emph{independent of $\mathbf{c}$}. Therefore
% \[
% \tilde{\mathbf{w}}(\theta;\mathbf{c})\oslash\tilde{\mathbf{w}}(\theta;\mathbf{1})=\mathbf{1},
% \quad\Rightarrow\quad
% \mathrm{spread}(\cdot)=0,
% \]
% so \eqref{eq:gamma_scale_inv} holds with $\gamma=0$.

% Moreover, within the broad and widely-used class of \emph{separable} scalarizations
% $m(\boldsymbol{\ell})=\sum_t w_t\,\phi(\ell_t)$, the log transform is essentially the unique
% choice that can completely cancel multiplicative rescaling in the sense of \eqref{eq:effective_weight}.
% Indeed, requiring $c\,\phi'(c x)$ to be independent of $c$ for all $c>0$ implies
% $\phi'(y)\propto 1/y$, hence $\phi(y)\propto \ln y$ (up to an additive constant).
% This explains why log-based objectives arise naturally as the ``optimal'' way to remove unit-induced
% trajectory bias.

% \paragraph{Heuristic scale correction via a reference point.}
% A common heuristic is to estimate an inverse scale from a reference iterate $\theta_{\mathrm{ref}}$:
% \begin{equation}
% a_t \triangleq \frac{1}{\ell_t(\theta_{\mathrm{ref}})+\epsilon},
% \qquad a_1~\text{normalized to }1,
% \label{eq:a_est}
% \end{equation}
% and then optimize a weighted linear objective on the corrected losses:
% \begin{equation}
% m_{\mathrm{LS}}(\theta)=\sum_{t=1}^T w_t\,a_t\,\ell_t(\theta).
% \label{eq:ls_a}
% \end{equation}
% For \eqref{eq:ls_a}, we have $\frac{\partial m_{\mathrm{LS}}}{\partial \ell_t}=w_t a_t$ and thus
% $\tilde w_t(\theta;\mathbf{c})=c_t w_t a_t$.
% If $a_t$ were proportional to $1/c_t$, then $\tilde w_t(\theta;\mathbf{c})$ would be approximately
% independent of $\mathbf{c}$, leading to a small $\gamma$ in \eqref{eq:gamma_scale_inv}
% (i.e., attenuating scale mismatch).

% However, under \eqref{eq:scale_model},
% \[
% a_t
% =
% \frac{1}{c_t\hat\ell_t(\theta_{\mathrm{ref}})+\epsilon},
% \]
% so $c_t a_t \approx 1/\hat\ell_t(\theta_{\mathrm{ref}})$ still depends on the (unknown and task-varying)
% aligned losses at the reference point, which introduces bias in the inferred scale ratios.
% Moreover, as training progresses, $\hat\ell_t(\theta)$ (and hence effective scales) drift, making a fixed
% $\theta_{\mathrm{ref}}$ stale. Consequently, the resulting influence pattern
% $\tilde{\mathbf{w}}(\theta;\mathbf{c})$ can gradually inherit more of the original mismatch
% (i.e., $\gamma$ increases), and scale dominance may re-emerge.


%%%%%%%%%%%%%%%%%%%%%%
% \paragraph{Scale mismatch and trajectory bias.}
% In multi-task learning, the loss terms associated with different tasks are often not directly comparable in scale.
% Such discrepancies can stem from intrinsic task heterogeneity—e.g., varying noise levels and uncertainty—as well as
% from differences in output spaces and physical units when mixing regression and classification objectives
%~\cite{kendall2018multi}. In practice, implementation choices further affect loss magnitudes, including the reduction
% convention (e.g., mean vs.\ sum) and task-dependent normalization (e.g., averaging over pixels, anchors, or sampled
% RoIs)~\cite{ren2015faster,he2017mask}. Consequently, aggregation strategies that operate on absolute loss values can
% implicitly prioritize large-scale objectives; prior work~\cite{chen2018gradnorm} shows that naive equal-weight summation
% may cause training to be dominated by tasks with larger loss scales, undermining the intended treatment across tasks.

% To formalize this issue, we distinguish between (i) the \emph{intrinsic} multi-objective trade-offs we care about and
% (ii) \emph{extrinsic} multiplicative rescalings induced by unit or implementation choices\footnote{***}. Let
% $\hat{\boldsymbol{\mathcal{L}}}(\theta)=(\hat{\mathcal{L}}_1(\theta),\dots,\hat{\mathcal{L}}_T(\theta))\in\mathbb{R}_{++}^T$\footnote{***what is the difference with $\mathbb{R}_{+}$ \edit{It means non-negative}}
% denote a latent \emph{canonical} loss vector\footnote{***} under some fixed (but arbitrary) choice of measurement units.
% Importantly, \emph{canonical does not mean balanced:} the magnitudes of $\hat{\mathcal{L}}_t(\theta)$ may still differ
% substantially across tasks due to intrinsic heterogeneity (difficulty, noise, data regime, etc.).
% We model unit-induced or implementation-induced scale mismatch by assuming the observed losses
% $\boldsymbol{{\mathcal{L}}}(\theta)$ are related to $\hat{\boldsymbol{{\mathcal{L}}}}(\theta)$ via an unknown task-wise rescaling:
% \begin{equation}
% \boldsymbol{{\mathcal{L}}}(\theta)=\mathbf{c}\odot\hat{\boldsymbol{{\mathcal{L}}}}(\theta)
% =\big(c_1\hat{\mathcal{L}}_1(\theta),\,c_2\hat{\mathcal{L}}_2(\theta),\,\dots,\,c_T\hat{\mathcal{L}}_T(\theta)\big),
% \label{eq:scale_model}
% \end{equation}
% where $c_t>0$. Since only relative scales matter, we fix task~1 as the reference and set $c_1=1$ w.l.o.g.;
% then $c_t$ (for $t\ge2$) represents the scale ratio of task $t$ relative to task~1.\footnote{***delete this paragraph}

% While Pareto dominance itself is invariant to positive rescaling, a scalarization-based optimizer can exhibit
% trajectory bias: under a rescaled observation $\boldsymbol{{\mathcal{L}}}=\mathbf{c}\odot\hat{\boldsymbol{{\mathcal{L}}}}$,
% the training signal produced by a scalarization $m(\boldsymbol{{\mathcal{L}}}(\theta))$ may tilt toward large $c_t$ tasks\footnote{***},
% leading to a solution that is an artifact of unit choices rather than reflecting the intended trade-off on the
% canonical objectives. This motivates the following principle:
% \emph{the induced optimization dynamics should be insensitive to relative rescaling.}\footnote{***delete this paragraph}

% \paragraph{Effective task influence induced by a scalarization.}
% Let $m:\mathbb{R}_{++}^T\to\mathbb{R}$ be differentiable. By the chain rule,
% \begin{equation}
% \nabla_\theta m(\boldsymbol{{\mathcal{L}}}(\theta))
% =
% \sum_{t=1}^T
% \frac{\partial m}{\partial {\mathcal{L}}_t}\big(\boldsymbol{{\mathcal{L}}}(\theta)\big)\,
% \nabla_\theta {\mathcal{L}}_t(\theta).
% \label{eq:chain_rule}
% \end{equation}
% Under \eqref{eq:scale_model}, $\nabla_\theta {\mathcal{L}}_t(\theta)=c_t\,\nabla_\theta \hat{\mathcal{L}}_t(\theta)$, hence
% \begin{equation}
% \nabla_\theta m\!\left(\mathbf{c}\odot\hat{\boldsymbol{{\mathcal{L}}}}(\theta)\right)
% =
% \sum_{t=1}^T
% \underbrace{\left(
% c_t\,\frac{\partial m}{\partial {\mathcal{L}}_t}\!\left(\mathbf{c}\odot\hat{\boldsymbol{{\mathcal{L}}}}(\theta)\right)
% \right)}_{\triangleq~\tilde w_t(\theta;\mathbf{c})}
% \,\nabla_\theta\hat{\mathcal{L}}_t(\theta).
% \label{eq:effective_weight}
% \end{equation}
% We refer to $\tilde w_t(\theta;\mathbf{c})$ as the \emph{effective influence coefficient} of task~$t$ on the
% aligned gradient $\nabla_\theta \hat{\mathcal{L}}_t(\theta)$ under rescaling $\mathbf{c}$.

% \paragraph{Scale invariance.}
% We quantify scale-robustness by measuring how much the relative rescaling $\mathbf{c}$ perturbs the influence pattern
% $\tilde{\mathbf{w}}(\theta;\mathbf{c})=(\tilde w_1,\dots,\tilde w_T)$ compared to the unscaled case $\mathbf{1}$.
% For multiplicative effects, a natural dispersion measure is the log-spread\footnote{***what is $v_t$? \edit{It means a vector}} \edit{defined for a vector $\mathbf{v}\in \mathbb{R}^T$ as} 
% \[
% \mathrm{spread}(\mathbf{v})
% ~\triangleq~
% \max_{t}\ln v_t - \min_{t}\ln v_t.
% \]
% We consider the distortion ratio
% $\tilde{\mathbf{w}}(\theta;\mathbf{c})\oslash \tilde{\mathbf{w}}(\theta;\mathbf{1})$,
% where $\oslash$ denotes elementwise division.

% \begin{definition}[$\gamma$-scale invariance]
% % A differentiable scalarization $m$ is \emph{$\gamma$-scale-invariant} w.r.t.\ the rescaling model
% % Eq.~\eqref{eq:scale_model} if there exists $\gamma\in[0,1]$ such that for any $\mathbf{c}\in\mathbb{R}_+^T$ with $c_1=1$
% % and all $\theta$,
% A differentiable scalarization $m$ is said to be \emph{$\gamma$-scale-invariant} if there exists
% $\gamma \ge 0$ such that, for any task-wise rescaling vector
% $\mathbf{c} \in \mathbb{R}_+^T$ with $c_1 = 1$, and for all $\theta$, the following holds:
% \begin{equation}
% \mathrm{spread}\!\Big(
% \tilde{\mathbf{w}}(\theta;\mathbf{c}) \oslash \tilde{\mathbf{w}}(\theta;\mathbf{1})
% \Big)
% ~\le~
% \gamma\,\mathrm{spread}(\mathbf{c}).
% \label{eq:gamma_scale_inv}
% \end{equation}
% \end{definition}

% The parameter $\gamma$ characterizes how sensitively the induced influence pattern responds
% to task-wise rescaling. In particular, $\gamma=0$ indicates perfect scale invariance,
% $\gamma=1$ corresponds to a scale-preserving behavior, and $\gamma>1$ implies that the
% relative rescaling is amplified in the induced influence pattern.

% \paragraph{Log transform.}
% A canonical scalarization is the log-sum objective
% \begin{equation}
% m_{\ln}(\boldsymbol{{\mathcal{L}}})=\sum_{t=1}^T w_t \ln {\mathcal{L}}_t,
% \qquad w_t>0,
% \label{eq:mlog}
% \end{equation}
% which is well-defined on $\mathbb{R}_{++}^T$. Since
% $\frac{\partial m_{\ln}}{\partial {\mathcal{L}}_t}=w_t/{\mathcal{L}}_t$,
% \[
% \tilde w_t(\theta;\mathbf{c})
% =
% c_t\cdot \frac{w_t}{c_t\hat{\mathcal{L}}_t(\theta)}
% =
% \frac{w_t}{\hat{\mathcal{L}}_t(\theta)},
% \]
% which is \emph{independent of $\mathbf{c}$}. Consequently,
% $\tilde{\mathbf{w}}(\theta;\mathbf{c})\oslash \tilde{\mathbf{w}}(\theta;\mathbf{1})=\mathbf{1}$ and the left-hand
% side of Eq.~\eqref{eq:gamma_scale_inv} is identically zero; thus $m_{\ln}$ achieves $\gamma=0$. So the log transform is the natural optimal choice. 

% % Moreover, log is essentially the unique way to \emph{fully cancel} multiplicative rescaling within the widely-used
% % class of separable scalarizations $m(\boldsymbol{{\mathcal{L}}})=\sum_t w_t\,\phi({\mathcal{L}}_t)$:
% % requiring $c\,\phi'(c x)$ to be independent of $c$ for all $c>0$ forces $\phi'(y)\propto 1/y$ and hence
% % $\phi(y)\propto \ln y$ (up to an additive constant). This explains why log-based objectives arise naturally as an
% % ``optimal'' mechanism for eliminating unit-induced trajectory bias.


% \paragraph{Heuristic scale correction via a reference point.}
% A common heuristic for mitigating scale imbalance ~\cite{dai2023improvable} is to estimate a per-task inverse-scale
% factor from a reference iterate,
% \begin{equation}
% a_t \triangleq \frac{1}{{\mathcal{L}}_t(\theta_{\mathrm{ref}})+\epsilon},
% \qquad a_1 \text{ normalized to }1,
% \label{eq:a_est}
% \end{equation}
% and then optimize a weighted linear objective on the corrected losses:
% \begin{equation}
% m_{\mathrm{LS}}(\theta)=\sum_{t=1}^T w_t\,a_t\,{\mathcal{L}}_t(\theta).
% \label{eq:ls_a}
% \end{equation}
% Intuitively, if $a_t$ accurately reflected the inverse rescaling (i.e., $a_t \propto 1/c_t$), then
% $a_t\,{\mathcal{L}}_t(\theta)$ would approximately remove unit-induced mismatch and reduce scale dominance in training.
% However, this estimate is only heuristic: ${\mathcal{L}}_t(\theta_{\mathrm{ref}})$ entangles the unknown scale factor and the
% (task- and state-dependent) canonical loss value, so $a_t$ can be biased as a proxy for $1/c_t$.
% Moreover, a fixed reference point could become unsatisfactory as training evolves and relative loss magnitudes drift.
% As a result, the intended scale correction may be inaccurate or degrade over time, potentially re-introducing (or even
% amplifying) disparity in task influence and leading to larger trajectory differences under rescaling.\footnote{***delete this paragraph}


%\noindent \textbf{Illustrative Example.}
To see that, we study a two-task synthetic problem~\cite{navon2022multi}, where the Pareto front is available in closed-form solutions as indicated by the gray curve in Fig.~\ref{fig:toy-scale}. 
% We apply different positive scaling factors to the two tasks to simply emulate the mismatch in measurement units across objectives. 
We construct this toy example by introducing two task weights of the objectives.
Specifically, we multiply the two task losses by different scaling factors, i.e., $10:1$ and $1:10$. This setup emulates a mismatch in measurement units across task losses, allowing us to examine how various methods behave under different task scales. %different produce distinct optimization trajectories and converge to different solutions under unit changes.
Note that this rescaling does not change the Pareto front. 
We run from multiple initializations (in black dots) and compare SIMS with linear scalarization method with equal-weight (EW), geometric loss strategy (GLS) \cite{chennupati2019multinet++} and two merit-function-based scalarization methods (i.e., STCH \cite{lin2024smooth} and FOOPS \cite{chen2025efficient}). More experimental details are provided in Appendix ~\ref{appdix:example}.

% Figure~\ref{fig:toy-scale} highlights the qualitative difference induced by scale handling. With EW (first plot), trajectories are systematically steered toward regions preferred by absolute-scale dominance, and the iterates collapse onto a narrow subset of the Pareto front. In contrast, the log-transformed objective (second plot) is \emph{strongly scale-invariant} under our definition, and its optimization dynamics are therefore insensitive to the injected rescaling. Correspondingly, trajectories approach the front in a more balanced manner and converge to solutions that are markedly less affected by the initial scale mismatch. Although deliberately simple, this example isolates the mechanism predicted by our invariance criterion that enforcing strong scale invariance can mitigate spurious bias introduced purely by loss magnitudes.


% The discussion above suggests that scale invariance should be enforced directly at the level of the training signal. In what follows, we make this intuition precise by introducing a scale-invariant advantage and a corresponding merit function, which are invariant to task-wise rescaling while remaining fully compatible with standard gradient-based optimization.

% The discussion above suggests that scale invariance should be enforced directly at the level of the training signal\footnote{***}.
% The proposed $\gamma$-scale invariance notion further provides a concrete design principle for scalarization:
% a desirable merit function should attenuate the influence distortion induced by task-wise rescaling, ideally achieving
% a small $\gamma$. In what follows, we make this intuition precise by introducing a scale-invariant advantage and a
% corresponding merit function, which are invariant to task-wise rescaling while remaining fully compatible with
% standard gradient-based optimization.

% Under the two scaling settings, 
Fig.~\ref{fig:toy-scale} shows that although the Pareto front remains unchanged, the EW, STCH, and FOOPS methods exhibit significant sensitivity to the scales of task losses. Specifically, under these two scaling factors (10:1 and 1:10), EW, STCH, and FOOPS exhibit different optimization trajectories, respectively, and they consistently drift toward the task loss with the larger scale. %These methods often yields to scale-dependent trade-offs and tend to converge to solutions dominated by a single task.
% that are close to the optimum of a single task.
We further observe that the FOOPS method exhibits optimization instability and even fails to converge to a stationary solution in some cases (i.e., the point circled in red). %\cyfoot{***cy: Circle this point with a red circle. \edit{Done}}. 
\zbedit{Although GLS is invariant to positive task-wise loss rescaling due to its geometric-loss formulation, it still suffers from local optima in this case and converges to solutions away from the Pareto front under both scaling settings. This indicates that scale insensitivity alone is insufficient without a merit-function-based criterion that explicitly characterizes the possibility of further joint improvement.}
In contrast, %\cyfoot{***cy: which method is log-based objective? and you haven't used or defined the term `log-based objective' before \edit{turn to describe our method}}
the proposed SIMS method exhibits a scale-invariant property due to the proposed transformation-induced merit function, %by design, since we intentionally adopt a merit-function-based scalarization that is insensitive to loss scaling,} 
leading to a more consistent convergence behavior under different scalings. \edit{This aligns with our goal of obtaining solutions that remain unchanged under different loss scales.} %\cyfoot{***cy: don't use robustness or robust in this paper. \edit{Now this phenomenon has been described}} 
\edit{Overall, to mitigate extreme trade-off solutions and achieve better average performance across tasks, scalarization methods in MTL should exhibit a certain degree of scale invariance.}%\cyfoot{***cy: why the conclusion here is `a certain degree of scale invariance is essential'. The toy example here needs to rewrite. \edit{More discussions were added}}
% \fei{Is the example want to show that when you use a transformation (scaling transformation in this example), under your methods, the results is invariant. If so, you should express that for a ‘good’ MTL method, its results should be invariant to scaling transformation, even for other transformation.}



\section{Methodology}

In this section, we present the proposed SIMS method.

\subsection{Formulation}
\label{section:formulation}
To achieve the scale invariance where the induced optimization dynamics is less affected by arbitrary task-wise scales, we propose a transformation-induced merit function. %an explicit  of each task loss into the merit function. 
Specifically, we first apply a mapping $\psi(\cdot)$ to the original task loss $\mathcal{L}_i(\theta)$, and %optimize the resulting transformed losses. To this end, we 
define the transformed task loss as
\begin{align}
\tilde{\mathcal{L}}_i(\theta)=\psi\!\left(\mathcal{L}_i(\theta)\right)\notag,
\end{align}
where $\psi(\cdot)$ is a transformation function shared by different tasks. 
%that can transform a nonnegative scalar to another nonnegative scalar. 
Moreover, $\psi(\cdot)$ should be a monotonically increasing function so that %it preserves the order of the original losses; that is, 
minimizing $\tilde{\mathcal{L}}_i(\theta)$ is equivalent to minimizing $\mathcal{L}_i(\theta)$.


% Moreover, $\psi$ should be a monotonically increasing\footnote{***cy: check increasing.} function as minimizing $\tilde{\mathcal{L}}_i(\theta)$ will make $\mathcal{L}_i(\theta)$ small.
% a monotone non-decreasing function designed to align loss scales across tasks.

Then, based on %\edit{given}\footnote{***given \edit{Done}}
the MOO formulation in problem~\eqref{eq:1}, %\footnote{***problem \eqref{eq:1} \edit{Done}}, 
the new MOO problem based on $\{\tilde{\mathcal{L}}_i(\theta)\}$ is formulated as
\begin{equation}
\min_{\theta \in \mathbb{R}^d} \! \tilde{\boldsymbol{\mathcal{L}}}(\theta)\!
    =\!
    \bigl(
        \psi(\mathcal{L}_1(\theta)),\!
        \psi(\mathcal{L}_2(\theta)),\!
        \ldots,
        \psi(\mathcal{L}_m(\theta))
    \bigr)^\top\!.
    \label{eq:2}
\end{equation}
According to the definition of the merit function in Definition \ref{def_merit_function}, a merit function associated with problem~\eqref{eq:2} should be non-negative and attain zero if and only if the solution is (weakly) Pareto optimal. Under this requirement and assuming lower semicontinuity of $\{\mathcal{L}_i\}$, the proposed transformation-induced merit function is formulated as
\begin{equation}
    \bar{u}(\theta)=\sup_{\theta' \in \mathbb{R}^d}\min_{i \in [m]}
    \bigl\{\psi(\mathcal{L}_i(\theta))-\psi(\mathcal{L}_i(\theta'))\bigr\}.
    \label{eq:ubar}
\end{equation}
$\bar{u}(\theta)$ measures the largest improvement that some other feasible solution can simultaneously achieve across all objectives over $\theta$. 

\subsection{Properties}
\label{sec5}

In the following theorem, we prove that $\bar{u}$ is a valid merit function in the sense of weak Pareto optimality. 
% In particular, $\bar{u}(\theta)=0$ if and only if $\theta$ is a weakly Pareto optimal solution of problem~\eqref{eq:1}.\footnote{***delete this sentence.}
\begin{theorem}
\label{theo:u0_merit}
Suppose that $\psi(\cdot)$ is a monotonically increasing function. Then we have $\bar{u}(\theta)\ge 0$ for all $\theta \in \mathbb{R}^d$. Moreover, $\theta \in \mathbb{R}^d$ is weakly Pareto optimal for problem~(1) if and only if $\bar{u}(\theta)= 0$. Hence, $\bar{u}$ is a merit function.
\end{theorem}

As shown in the motivating example, a scalarization method for MTL is expected to be invariant to the scale of task losses. We formalize the definition of this scale-invariance requirement as follows.
% As shown in the motivating example, a reasonable scalarization or merit-function-based formulation for MTL should be invariant to task-wise positive rescaling of losses. We formalize this scale-invariance requirement below.

\begin{definition}[Scale-Invariance]
\label{def:scale_invariance}
Consider an MTL problem with task losses $\{\mathcal{L}_i\}_{i=1}^m$. For any positive constants $c_i > 0$, the rescaled losses is defined as $\bar{\mathcal{L}}_i(\theta) = c_i\,\mathcal{L}_i(\theta)$. Then, a scalarization MTL method is said to be \emph{invariant to the scaling transformation} if the scalarized objective induced by $\{ \bar{\mathcal{L}}_i(\theta)\}_i^m$ differs from the original scalarized objective only by an additive constant that is independent of \(\theta\).
% Then an MTL solution $\theta^\star$ is said to be \emph{invariant to the scaling transformation} if $\theta^\star$ is optimal under $\{\mathcal{L}_i\}_{i=1}^m$ if and only if it is optimal under $\{\bar{\mathcal{L}}_i\}_{i=1}^m$.
\end{definition}

Definition~\ref{def:scale_invariance} requires scalarization MTL methods to be invariant to task-wise rescaling of losses, so that changing the units or magnitudes of individual task losses does not affect the preference of the results. Without such a property, a task may dominate the optimization merely due to the large scale of its task loss. 

For the functional form of the transformation function $\psi(\cdot)$, in the following theorem, we prove that the scale invariance property of the transformation-induced merit function $\bar u(\theta)$ makes $\psi(\cdot)$ a logarithmic function.

\begin{theorem}%[Scale invariance characterizes logarithmic transforms]
\label{thm:log_characterization}
When solving an MTL problem by minimizing $\bar{u}(\theta)$ in Eq.~\eqref{eq:ubar}, if the optimal solution is required to be invariant to the scaling transformation, then $\psi$ must be of the form $\psi(x) = a \ln x + b$ for some constants $a > 0$ and $b \in \mathbb{R}$.
% \footnote{***change $\log$ to $\ln$ in the whole paper \edit{Done}}
\end{theorem}

The constants $a$ and $b$ in $\psi(\cdot)$ only introduce a positive affine rescaling and therefore do not change the ordering of objective values. Consequently, they do not affect the solution set of problem~\eqref{eq:ubar}. Hence, there is no need to tune $a$ and $b$ in practice. Without loss of generality, we adopt a simple choice in experiments, that is, $\psi(x) = \ln x$, which satisfies the required monotonicity and scale-invariance properties.

\subsection{Approximation}
\label{sec:app}
Since $\bar{u}(\theta)$ has been proved in Theorem \ref{theo:u0_merit} to be a merit function of problem \eqref{eq:2}, we can minimize $\bar{u}(\theta)$ to obtain a weakly Pareto optimal solution. 
%In this subsection, we develop a practical approach for solving problem~\eqref{eq:ubar} using first-order methods. Importantly, the proposed procedure is applicable to \emph{any} transformation function $\psi$, including the logarithmic transform we discussed above. According to Theorem~\ref{theo:u0_merit}, problem~\eqref{eq:1} can be reformulated as minimizing the scalar-valued function $\bar{u}(\theta)$. 
However, due to the pointwise $\min$ operator, $\bar{u}$ is generally nonsmooth and non-differentiable~\cite{boyd2004convex}. 
This prevents the direct use of gradient descent and instead necessitates subgradient-based methods, which typically converge more slowly~\cite{goffin1977convergence}. %\footnote{***most of this paragraph should be put in the next section. \edit{Done}}


To this end, we construct a smoothed and regularized surrogate of $\bar{u}(\theta)$ inspired by classical smoothing techniques (i.e.,log-sum-exp smoothing~\cite{nesterov2005smooth,beck2017first}) for nonsmooth optimization. 
Note that the proposed optimization procedure is applicable to any transformation function $\psi$. 
% e.g., log-sum-exp smoothing~\cite{beck2017first}, Nesterov smoothing~\cite{nesterov2005smooth,boyd2004convex}, and inf-conv smoothing~\cite{beck2012smoothing})
Specifically, we define%\footnote{***no subject in this sentence. \edit{Done} change to "we define"} 
\begin{align}
h_{\lambda, \tau}(\theta,\theta')
=&\tau \ln \left(\sum_{i=1}^{m}\exp\left(\frac{\psi(\mathcal{L}_i(\theta'))-\psi(\mathcal{L}_i(\theta))}{\tau}\right)\right)\notag+\frac{\lambda}{2}\|\theta-\theta'\|^2 \nonumber\\
v_{\lambda, \tau}(\theta)=&-\min_{\theta'\in \mathbb{R}^d}h_{\lambda,\tau}(\theta,\theta'),\label{eq:v}
\end{align}
where $\|\cdot\|$ denotes the Euclidean norm, $\tau>0$ is a smoothing parameter, and $\lambda\ge 0$ is a regularization parameter. Therefore, problem~\eqref{eq:1} is converted to minimizing $v_{\lambda,\tau}$, which is equivalent to solving the following saddle-point problem as
\begin{equation}
    % \min_{\theta\in \mathbb{R}^d} v_{\lambda,\tau}(\theta)
    % \;=\;
    % -
    \max_{\theta\in \mathbb{R}^d}\min_{\theta'\in \mathbb{R}^d}
    h_{\lambda,\tau}(\theta,\theta').
    \label{eq:vltau}
\end{equation}
Note that when $\lambda=0$, $v_{0,\tau}$ is a log-sum-exp smoothing approximation of $\bar{u}$, with the approximation accuracy controlled by $\tau$~\cite{nesterov2005smooth,boyd2004convex}. Introducing $\lambda>0$ adds a quadratic regularization term that improves the curvature of the inner problem, which will be shown to yield a well-conditioned surrogate with stable gradients. %Importantly, this max–min reformulation allows simultaneous gradient updates of the outer variable $\theta$ and the inner variable $\theta'$, making gradient-based saddle-point methods directly applicable~\cite{lin2020gradient,xu2023unified,zhang2024generalization}.


% According to Theorem~\ref{theo:u0_merit}, it is natural to convert the original multi-objective problem into  minimizing the scalar-valued function $\bar{u}(\theta)$. However, the pointwise $\min$ operator makes $\bar{u}$ generally nonsmooth and thus non-differentiable~\cite{boyd2004convex}. 
% Therefore, first-order methods can only rely on subgradient methods rather than gradient descent. It is well known that minimizing a nonsmooth convex function using subgradient methods typically requires $\mathcal{O}(1/\epsilon^{2})$ iterations to reach $\epsilon$-optimal solution~\cite{goffin1977convergence}, whereas for smooth convex objectives, gradient descent achieves a faster rate of $\mathcal{O}(1/\epsilon)$. To enable efficient gradient-based optimization, we next reformulate~\eqref{eq:ubar} into a max–min problem with a strongly convex inner objective, which admits well-defined gradients and can directly benefit from gradient descent.\footnote{***most of this paragraph should be put in the next section. \edit{Done}}

% Motivated by classical smoothing techniques for nonsmooth optimization (e.g., Nesterov smoothing~\cite{nesterov2005smooth,boyd2004convex} and inf-conv smoothing~\cite{beck2012smoothing}), we construct a smooth and regularized surrogate of $\bar{u}(\theta)$ in Eq.~\eqref{eq:ubar}. Specifically, we define

% This max-min structure enables simultaneous updates of the outer variable $\theta$ and the inner variable $\theta'$ using only first-order information, making gradient-based saddle-point methods directly applicable~\cite{lin2020gradient,xu2023unified,zhang2024generalization}.


% When $\lambda=0$, the log-sum-exp term in \edit{Eq.}\footnote{***change to Eq. \edit{Done}}~\eqref{eq:hltau} provides a smooth approximation to $\bar{u}(\theta)$ with accuracy controlled by $\tau$~\cite{nesterov2005smooth,boyd2004convex}. With a suitable $\lambda$, the quadratic term  can improve the inner minimization structure\footnote{***what structure?} and, as shown later, leads to a well-conditioned surrogate that admits stable gradients.

% The construction in Eqs.~\eqref{eq:hltau}--\eqref{eq:v} has several desirable properties. First, $v_{0,\tau}(\theta)$ serves as a proper smooth approximation to the original nonsmooth merit function $\bar{u}(\theta)$, and it uniformly approaches $\bar{u}(\theta)$ as $\tau\rightarrow 0$. Second, the regularized surrogate $v_{\lambda,\tau}(\theta)$ preserves the keymerit-function characterization of weak Pareto optimality in an approximate sense: its value remains close to that of $\bar{u}(\theta)$ up to a controllable error depending on $\tau$, while the regularization improves the inner optimization structure and enables tractable first-order analysis. These advantages make $v_{\lambda,\tau}$ particularly appealing, as it provides a scalar objective that is (i) differentiable, (ii) theoretically analyzable, and (iii) avoids explicitly introducing preference weights, thereby offering a principled alternative to classical scalarization-based formulations.

%\footnote{***delete this paragraph since the following analysis reveals the properties.~\edit{Rewrite briefly and delete former vision. Done}}

In the following, we show that $v_{\lambda,\tau}$ preserves the key properties of $\bar{u}$, while quantifying the approximation gap. We begin by showing that $v_{0,\tau}$ is a smooth approximation of $\bar{u}$.
\begin{proposition}[Smooth Approximation]
\label{prop:smooth_approx}
The function $v_{0, \tau}(\theta)$ is a smoothing approximation of $\bar{u}(\theta)$. That is, for any $\tau>0$, $v_{0, \tau}(\theta)$ is continuously differentiable on
$\mathbb{R}^d$ and satisfies the properties:
\begin{enumerate}
    \item \emph{Consistency:}\;
    $\lim_{\tau\to 0} v_{0,\tau}(\theta)=\bar{u}(\theta)$ for all $\theta\in \mathbb{R}^d$;
    \item \emph{Smoothness:}\;
    there exist constants $C>0$ and $\alpha>0$, independent of $\tau$, such that
    $v_{0,\tau}$ has a Lipschitz gradient on $\mathbb{R}^d$ with the Lipschitz constant
    $L_{v_{0,\tau}} = C + \alpha/\tau$.
\end{enumerate}
\end{proposition}

\noindent
In other words, Proposition~\ref{prop:smooth_approx} shows that $v_{0, \tau}(\theta)$ shares similar properties with the transformation-induced merit function $\bar{u}(\theta)$, while being much easier to be optimized via gradient-based methods due to its differentiable nature. Moreover, it is easy to see that $\bar{u}(\theta)$ is properly upper- and lower-bounded by $v_{0,\tau}(\theta)$.
\begin{proposition}[Bounded Approximation]
\label{prop:bounded_approx}
For any $\theta\in \mathbb{R}^d$, we have $v_{0,\tau}(\theta) \le \bar{u}(\theta)\leq v_{0,\tau}(\theta) + \tau\ln m $.
\end{proposition}
Proposition~\ref{prop:bounded_approx} implies that $v_{0,\tau}$ provides a \emph{uniform} smooth approximation to $\bar{u}$ with a gap controlled by $\tau$. In particular, choosing $\tau$ sufficiently small makes the bound tight uniformly over $\mathbb{R}^d$.

Next, we show that the regularization $\lambda\geq 0$ may yield a strongly-convex inner structure in problem~\eqref{eq:vltau}, which ensures that the inner minimizer is well-defined (i.e., it exists and is unique) and facilitates the first-order analysis. We first make the weak convexity assumption, which is widely used in nonconvex optimization~\cite{lin2020gradient,xu2023unified,chen2025efficient}.
\begin{assumption}[Weak convexity]
\label{ass:weak_convexity}
For each $i\in[m]$, the function $\tilde{\mathcal{L}}_i:\mathbb{R}^d \to\mathbb{R}$ is assumed to be $\mu_i$-weakly convex on $\mathbb{R}^d$.
That is, $\tilde{\mathcal{L}}_i(\theta)+\frac{\mu_i}{2}\|\theta\|^2$ is convex on $\mathbb{R}^d$.
\end{assumption}

\begin{lemma}[Uniqueness of the inner minimizer]
\label{lemmsc}
Suppose that Assumption~\ref{ass:weak_convexity} holds. If $\lambda>\bar{\mu}$, where $\bar{\mu}:=\max_{i\in[m]}\mu_i$, then for any fixed $\theta\in \mathbb{R}^d$, function $h_{\lambda,\tau}(\theta,\theta')$ is strictly convex with respect to $\theta'$. 
Consequently, the inner problem $\min_{\theta'\in \mathbb{R}^d} h_{\lambda,\tau}(\theta,\theta')$ admits a unique minimizer, denoted by $
\theta^{\star}_{\lambda,\tau}(\theta)
=
\arg\min_{\theta'\in\mathbb{R}^d} h_{\lambda,\tau}(\theta,\theta')$.
\end{lemma}

Lemma~\ref{lemmsc} shows that when $\lambda > \bar{\mu}$, the inner problem becomes uniformly strongly-convex, which guarantees a unique minimizer and enables stable gradient-based updates with linear convergence~\cite{boyd2004convex}. Moreover, through the implicit minimizer mapping $\theta^{\star}_{\lambda,\tau}(\theta)$, the outer objective $v_{\lambda,\tau}$ inherits favorable smoothness properties. Therefore, we will impose the condition $\lambda > \bar{\mu}$ throughout the remainder of the analysis.

Finally, we show that the surrogate $v_{\lambda,\tau}(\theta)$ preserves the weak Pareto optimal property up to a controllable tolerance.

\begin{proposition}
\label{theov}
Suppose that Assumption~\ref{ass:weak_convexity} holds. Then $v_{\lambda,\tau}(\theta)$ satisfies the following properties :
\begin{enumerate}
    \item (Lower bound) $v_{\lambda,\tau}(\theta) \ge -\tau \ln m$ for any $\theta \in \mathbb{R}^d$.
    \item (Necessity) If $\theta$ is weakly Pareto optimal for problem~\eqref{eq:1}, then $v_{\lambda,\tau}(\theta) \le 0$.
    \item (Approximate sufficiency) If $\lambda = 0$ and $v_{0,\tau}(\theta) \le 0$, then $\theta$ is $\epsilon$-weakly Pareto optimal with $\epsilon = \tau \ln m$.
    \item (Sufficiency under convexity) If $\lambda > 0$, $v_{\lambda,\tau}(\theta) \le -\tau \ln m$, and each $\mathcal{L}_i$ is convex at $\theta$, then $\theta$ is weakly Pareto optimal.
\end{enumerate}
\end{proposition}
% \begin{proposition}
% \label{theov}
% Suppose that Assumption~\ref{ass:weak_convexity} holds. % and let $v_{\lambda,\tau}$ be defined in Eq.~\eqref{eq:v}. 
% Then $v_{\lambda,\tau}(\theta)\ge -\tau\ln m$ for all $\theta\in \mathbb{R}^d$. Moreover, if $\theta$ is weakly Pareto optimal for problem~\eqref{eq:1}, then $v_{\lambda,\tau}(\theta)\le 0$. Conversely,
% \begin{enumerate}
%     \item if $\lambda=0$ and $v_{0,\tau}(\theta)\le 0$, then $\theta$ is $\epsilon$-weakly Pareto optimal with $\epsilon=\tau\ln m$;
%     \item if $\lambda>0$, $v_{\lambda,\tau}(\theta)\le -\tau\ln m$, and each $\mathcal{L}_i$ is convex at $\theta$,
%     then $\theta$ is weakly Pareto optimal.
% \end{enumerate}
% \end{proposition}
When $\lambda = 0$, the smoothing introduces at most an $\epsilon = \tau \ln m$ error\footnote{***an approximation error of $\epsilon = \tau \ln m$? \edit{Done} Yes, relaxation changes error}, meaning that minimizing $v_{0,\tau}$ yields an $\epsilon$-weakly Pareto optimal solution to the original problem. This result justifies $v_{\lambda,\tau}$ as a principled surrogate for $\bar{u}(\theta)$ with a clear trade-off between smoothness and Pareto optimality gap\footnote{***what does Pareto accuracy mean?\edit{Done} changes to optimality gap} controlled by $\tau$.


% Together, Propositions~\ref{prop:smooth_approx}-\ref{theov} justify optimizing $v_{\lambda,\tau}$ as a principled, differentiable surrogate of the original nonsmooth merit function $\bar{u}$, with controllable approximation error and preserved Pareto interpretation. 
% To visually illustrate the effect of the smoothing and regularization parameters,
% Figure~2 presents a simple example comparing $\bar{u}$ and $v_{\lambda,\tau}$.
% As expected, smaller values of $\tau$ and $\lambda$ yield a closer approximation to $\bar{u}$,
% while larger values produce a smoother but more biased surrogate.
% In particular, $\tau$ primarily controls the smoothness--accuracy trade-off of the log-sum-exp approximation,
% whereas $\lambda>0$ further stabilizes the inner minimization by improving its curvature,
% which in turn leads to better-conditioned gradients for first-order optimization.
% We therefore consider the smooth optimization problem
% \begin{align}
%     \min_{\theta \in \mathbb{R}^d} 
%     v_{\lambda,\tau}(\theta)\notag.
% \end{align}

% Recalling Eq.~\eqref{eq:v}, minimizing $v_{\lambda,\tau}$ is equivalent to solving the saddle-point problem
% \begin{equation}
%     \min_{\theta\in \mathbb{R}^d} v_{\lambda,\tau}(\theta)
%     \;=\;
%     -\max_{\theta\in \mathbb{R}^d}\min_{\theta'\in \mathbb{R}^d}
%     h_{\lambda,\tau}(\theta,\theta').
%     \label{eq:vltau}
% \end{equation}
% This max-min structure enables simultaneous updates of the outer variable $\theta$ and the inner variable $\theta'$ using only first-order information, making gradient-based saddle-point methods directly applicable~\cite{lin2020gradient,xu2023unified,zhang2024generalization}.

\subsection{Algorithm and Convergence Analysis}\label{analysis}
% $\lambda>\bar{\mu}$, 
To solve the saddle-point problem~\eqref{eq:vltau}, 
%is a smooth nonconvex-strongly-convex problem. This naturally motivates
we adopt the two-time-scale gradient descent-ascent (TTGDA) method~\cite{lin2020gradient,xu2023unified}, where the inner variable $\theta'$ is updated with a larger learning rate to rapidly track the inner minimizer, while the outer variable $\theta$ is updated more conservatively to ensure stable descent. % of the surrogate objective $v_{\lambda,\tau}$.
Algorithm~\ref{alg1} summarizes the TTGDA method for solving problem~\eqref{eq:vltau}.

To compute the gradients, we define the weights of gradient aggregation as
\[
w_i(\theta,\theta') :=
\frac{\exp\!\left(\frac{\tilde{\mathcal{L}}_i(\theta')-\tilde{\mathcal{L}}_i(\theta)}{\tau}\right)}
{\sum_{j=1}^m \exp\!\left(\frac{\tilde{\mathcal{L}}_j(\theta')-\tilde{\mathcal{L}}_j(\theta)}{\tau}\right)}, 
\quad i\in[m].
\]
These weights arise directly from the log-sum-exp smoothing and form a probability simplex. They emphasize tasks whose improvement from $\theta$ to $\theta'$ is relatively small, thereby adaptively reweighting task gradients based on local changes in the objective values.\footnote{***what does `local first-order information' mean? \edit{Done} changes to local changes in the objective values}. %This mechanism is similar to softmax-based scalarization~\cite{lin2019pareto} and adaptive reweighting schemes~\cite{lin2024smooth,lu2025learning} in MOO, but here it emerges naturally from the smooth saddle-point formulation rather than being imposed heuristically. 
Using these weights, the gradients of $h_{\lambda,\tau}(\theta,\theta')$ with respect to $\theta$ and $\theta'$ are computed as
\begin{align}
\nabla_{\theta'} h_{\lambda,\tau}(\theta,\theta')
&=
\sum_{i=1}^m w_i(\theta,\theta')\nabla \tilde{\mathcal{L}}_i(\theta')
+
\lambda(\theta'-\theta),
\label{eq:grad_1}
\\\!
\nabla_{\theta} h_{\lambda,\tau}(\theta,\theta')
&=\!
-\sum_{i=1}^m w_i(\theta,\theta')\!\nabla \tilde{\mathcal{L}}_i(\theta)
\!+\! \lambda(\theta-\theta').
\label{eq:grad_2}
\end{align}


\begin{algorithm}[t]
\caption{TTGDA Algorithm}
\label{alg1}
\begin{algorithmic}[1]
\Require Iterations $K$, smoothing parameters $\tau>0$, $\lambda\ge 0$,
stepsizes $\eta_{\theta'} \gg \eta_{\theta}$,
initial points $\theta^{0},\theta'^{0}\in\mathbb{R}^d$.
\For{$k=0,1,\ldots,K-1$}
\State Compute gradients $\nabla_{\theta'} h_{\lambda,\tau}(\theta^{k},\!\theta'^{k})$ and $\nabla_{\theta} h_{\lambda,\tau}(\theta^{k},\!\theta'^{k})$ via Eqs.~\eqref{eq:grad_1} and~\eqref{eq:grad_2}
    \State Update
    \begin{align}
        &\theta'^{k+1} \leftarrow \theta'^{k} - \eta_{\theta'}\, \nabla_{\theta'} h_{\lambda,\tau}\!\left(\theta^{k},\theta'^{k}\right),\notag \\
        &\theta^{k+1}\leftarrow \theta^{k} + \eta_{\theta}\, \nabla_{\theta} h_{\lambda,\tau}\!\left(\theta^{k},\theta'^{k}\right). \notag
    \end{align}
\EndFor
\State Randomly draw $\hat{\theta}$ uniformly from $\{\theta^k\}_{k=0}^{K}$.
\Ensure $\hat{\theta}$.
\end{algorithmic}
\end{algorithm}
%We analyze the iteration complexity of Algorithm~\ref{alg1} for solving~\eqref{eq:vltau} in the next subsection. The analysis follows the standard framework for TTGDA methods in smooth nonconvex--strongly-convex minimax problems~\cite{lin2020gradient,lin2025two}.

Next, we will analyze the convergence of the TTGDA algorithm. We begin with some basic assumptions.

\begin{assumption}[Lipschitz continuity]
\label{ass:lipschitz}
For each $i\in[m]$, the function $\tilde{\mathcal{L}}_i:\mathbb{R}^d \to\mathbb{R}$ is
$\ell_i$-Lipschitz continuous on $\mathbb{R}^d$. That is, for all
$\theta_1,\theta_2\in \mathbb{R}^d$, $\big|\tilde{\mathcal{L}}_i(\theta_1)-\tilde{\mathcal{L}}_i(\theta_2)
\big|\le \ell_i \|\theta_1-\theta_2\|$.
\end{assumption}


\begin{assumption}[Smoothness]
\label{ass:smoothness}
For each $i\in[m]$, the function $\tilde{\mathcal{L}}_i$ is $\zeta_i$-smooth on $\mathbb{R}^d$. That is, for all $\theta_1,\theta_2\in \mathbb{R}^d$, $\|\nabla \tilde{\mathcal{L}}_i(\theta_1)-\nabla \tilde{\mathcal{L}}_i(\theta_2)\|
\le \zeta_i \|\theta_1-\theta_2\|$.
\end{assumption}

The Lipschitz continuity and smoothness assumptions in Assumptions~\ref{ass:lipschitz} and~\ref{ass:smoothness} are widely adopted in MOO~\cite{ sener2018multi,lin2024smooth,chen2025gradient}. Based on these assumptions, we can characterize the smoothness of the saddle function $h_{\lambda,\tau}(\theta,\theta')$.
\begin{lemma}\label{smoothnessh}
Under Assumptions~\ref{ass:lipschitz} and~\ref{ass:smoothness}, define $\bar{\ell}:=\max_{i\in[m]}\ell_i$ and $\bar{\zeta}:=\max_{i\in[m]}\zeta_i$. Then $h_{\lambda,\tau}(\theta,\theta')$ is $\ell$-smooth jointly in $(\theta,\theta')$, with $\ell=\frac{\bar{\ell}^2}{\tau}+\bar{\zeta}+\lambda$.
\end{lemma}

\noindent
% Since $v_{\lambda,\tau}(\theta)
% =-\min_{\theta'} h_{\lambda,\tau}(\theta,\theta')$ is differentiable, our goal is to find an $\epsilon$-stationary point of $v_{\lambda,\tau}$ using only first-order information of $h_{\lambda,\tau}$. 
We define $\Delta_v$ and $\delta^0$ as
\[
\Delta_v =v_{\lambda,\tau}(\theta^0) - \min_{\theta\in\mathbb{R}^d} v_{\lambda,\tau}(\theta), \quad \delta^0:=\bigl\|\theta^{\star}_{\lambda,\tau}(\theta^0)-\theta'^0\bigr\|^2.
\]
Then in the following theorem, we analyze the convergence property of Algorithm~\ref{alg1} for solving problem~\eqref{eq:vltau}.

\begin{theorem}
\label{thm:gda_complexity}
Under Assumptions~\ref{ass:weak_convexity},~\ref{ass:lipschitz} and~\ref{ass:smoothness}, we set  $\lambda>\bar{\mu}$, $\eta_{\theta'}=\Theta(1/\ell)$, and $\eta_{\theta}=\Theta(1/(\kappa^2\ell))$, where $\kappa=\frac{\ell}{\lambda-\bar{\mu}}>0$ denote the condition number~\cite{lin2020gradient}. The iteration complexity (also the gradient complexity) of Algorithm~\ref{alg1} to obtain an $\epsilon$-stationary point is bounded by
\[
\mathcal{O}\!\left(
\frac{\kappa^2\ell\,\Delta_v
+
\kappa\ell^2 (\delta^0)^2}{\epsilon^2}
\right).
\]
\end{theorem}

Theorem~\ref{thm:gda_complexity} establishes the convergence guarantee of TTGDA for problem~\eqref{eq:vltau}. The condition number $\kappa$ captures the intrinsic asymmetry between the outer and inner variables and leads to a two-time-scale stepsize separation $\eta_{\theta'}/\eta_{\theta}=\Theta(\kappa^2)$. Such a separation is necessary for stability when using gradient methods in nonconvex-strongly-convex minimax optimization~\cite{lin2020gradient,xu2023unified,zhang2024generalization}.\footnote{***any technical innovation in this analysis? if yes, highlight it.}

% Alternative approaches, such as Alt-GDA~\cite{nouiehed2019solving,lin2020gradient} or recent single-loop methods~\cite{kovalev2022first,abernethy2021last}, have been proposed for related minimax settings. However, TTGDA provides a simpler and more direct fit to the structure of~\eqref{eq:vltau}, yielding a clean complexity bound that depends only on the initial suboptimality gap $\Delta_v$ and the initial tracking error $\delta^0$. In line with classical results in nonconvex optimization, the algorithm is guaranteed to produce an $\epsilon$-stationary point within a finite number of iterations, though this point may not coincide with the final iterate. This type of guarantee is standard for gradient-based methods and is commonly implemented by selecting an iterate from the optimization path~\cite{lin2020gradient,lin2025two}.





\section{Experiments}
% \fei{Just say 'In this section, we evaluate SIMS on three standard multi-task benchmarks to demonstrate its effectiveness across different MTL problems.'Others should put in latter}
In this section, we evaluate SIMS on three standard multi-task benchmarks
% comprising NYUv2~\cite{silberman2012indoor}, CityScapes~\cite{cordts2016cityscapes}, and PASCAL-Context~\cite{everingham2010pascal}
to demonstrate its effectiveness across diverse task compositions and objective scales. Furthermore, we conduct ablation studies on the smoothing parameter $\tau$ and the impact of the regularization parameter $\lambda$. In addition, we investigate different transformation functions $\psi$ and demonstrate the superiority of SIMS over heuristic loss normalization strategies.
% Furthermore, we analyze the impact of different transformation functions $\psi$ in the ablation studies.
% \vspace{-pt}



\subsection{Experimental Setup}
\noindent \textbf{Datasets.}
% \fei{Add Dataset with a brief intro for each here}
The following datasets are used: (i) NYUv2~\cite{silberman2012indoor}, which is an indoor
scene understanding dataset. It has 3 tasks (13-class semantic segmentation, depth estimation, and
surface normal prediction) with 795 training and 654 testing images. 
(ii) CityScapes~\cite{cordts2016cityscapes}, which is an urban scene understanding dataset. It has 2 tasks (7-class semantic segmentation and depth estimation) with 2,975 training and 500 testing images.
And (iii) PASCAL-Context~\cite{everingham2010pascal}, which is a challenging scene understanding dataset for images in the wild. We conduct experiments on 4 tasks (21-class semantic segmentation, 7-class human parts segmentation, saliency estimation, and surface normal estimation) with 4,998 training and 5,105 testing images. These benchmarks were selected for their diverse task definitions and supervision statistics. This diversity induces large disparities in objective scales, serving as a rigorous setting to demonstrate the scale-invariant property of our proposed method.
\vspace{0.5\baselineskip}

\noindent \textbf{Baselines.}
We compare the proposed SIMS with other scalarization methods, including Equal Weight (EW), STCH~\cite{lin2024smooth}, and FOOPS~\cite{chen2025efficient}. To assess competitiveness against state-of-the-art multi-task architectures, we also compare with representative methods across three categories. For CNN-based methods, we include Hard-Parameter Sharing (HPS), employing a shared backbone with task-specific output heads, Cross-Stitch~\cite{misra2016cross}, Multi-Task Attention Network (MTAN)~\cite{liu2019end}, and NDDR-CNN~\cite{gao2019nddr}. Regarding Transformer-based methods, we include VTAGML~\cite{bhattacharjee2023vision}, SwinMTL~\cite{taghavi2024swinmtl}, and DenseMTL \cite{lopes2023cross}. Finally, we consider foundation models utilizing parameter-efficient fine-tuning  \cite{hu2022lora}, 
% \footnote{***cy: add ref. \edit{Done}}
such as MultiLoRA~\cite{wang2023multilora} and MTSAM~\cite{wang2025mtsam}.
% state of the art methods that are commonly used for multi task learning with deep backbones, including CNN-based methods: Hard-Parameter Sharing (HPS), Cross-Stitch~\cite{misra2016cross}, Multi-Task Attention Network (MTAN)~\cite{liu2019end} and NDDR-CNN~\cite{gao2019nddr};Transformer-based approaches: VTAGML~\cite{bhattacharjee2023vision}, SwinMTL~\cite{taghavi2024swinmtl} and DenseMTL~\cite{lopes2023cross}; Foundation-model-based approaches equipped with parameter-efficient fine-tuning: MultiLoRA~\cite{wang2023multilora} and MTSAM~\cite{wang2025mtsam}.
\vspace{0.5\baselineskip}

\noindent \textbf{Evaluation metric.}
% For each experiment,
We report the common evaluation metrics for each task and provide comprehensive details in Appendix~\ref{appdix:metric}. Following the experimental setup in~\cite{maninis2019attentive},  we use the average of the relative improvement ($\Delta_b$) of each task over the baseline method (i.e., HPS) to quantify overall performance, which is formulated as $\Delta_b = \frac{1}{T} \sum_{i=1}^{T} \frac{1}{K_i} \sum_{j=1}^{K_i}
\frac{(-1)^{s_{i,j}} \left( M_{i,j}^{p} - M_{i,j}^{\text{base}} \right)}
{M_{i,j}^{\text{base}}}$, where 
% \fei{in previous section, we have use $m$ for task number, better to choose another notation}
$T$ denotes the number of tasks and $K_i$ is the number of evaluation metrics associated with task $i$.
$M_{i,j}^{p}$ and $M_{i,j}^{\text{base}}$ represent the performance of method $p$ and the baseline method on the $j$-th metric of task $i$, respectively.
The indicator $s_{i,j} \in \{0,1\}$ specifies the optimization direction of each metric, where $s_{i,j}=1$ indicates that lower values correspond to better performance, and $s_{i,j}=0$ otherwise.
\vspace{0.5\baselineskip}



\noindent \textbf{Model.}
% We adopt a modern MTL setup
% % \fei{should be cited} 
% that combines a pretrained backbone with parameter-efficient fine-tuning. 
% Given that our benchmarks consist of dense prediction tasks,
We adopt the paradigm of adapting pretrained backbones via PEFT \cite{hu2022lora,liu2022polyhistor}. 
Specifically, we utilize the image encoder from SAM2.1-Large \cite{ravi2024sam} as the foundational feature extractor. To adapt this frozen backbone for MTL, we employ vanilla Low-Rank Adaptation (LoRA) \cite{hu2022lora}. Formally, LoRA approximates the weight update $\Delta W$ for a pretrained weight matrix $W_0 \in \mathbb{R}^{d \times k}$ as the product of two low-rank matrices $B \in \mathbb{R}^{d \times r}$ and $A \in \mathbb{R}^{r \times k}$, where the rank $r \ll \min(d, k)$.
% for a pretrained weight matrix $W_0 \in \mathbb{R}^{d \times k}$, LoRA constrains the update $\Delta W$ by representing it as the product of two low-rank matrices $B \in \mathbb{R}^{d \times r}$ and $A \in \mathbb{R}^{r \times k}$, where the rank $r \ll \min(d, k)$.
The forward pass is thus modified as $h = W_0 x + s\Delta W x = W_0 x + sB A x$, where $s$ is a scaling coefficient.
We inject these shared LoRA adapters into the Query, Key, and Value (QKV) projections of the transformer layers to learn shared representations across different downstream tasks while keeping most parameters frozen. Finally, following prior work~\cite{wang2025mtsam}, we employ a modified lightweight prompt encoder to decode these features for specific outputs.

\begin{table*}[!phbt]
\centering

\setlength{\tabcolsep}{6pt}
\renewcommand{\arraystretch}{1}
\caption{Performance on three tasks on the NYUv2 dataset. 
Results marked with $^{\dagger}$ are copied from prior work. The best results for each task are shown in \textbf{bold}, and the second best are underlined.
$\uparrow(\downarrow)$ indicates that higher (lower) is better.}
\vspace{-3pt}
\label{table:nyu}
\begin{tabular}{lcccccccccc}
\toprule
\multirow{3}{*}[-1ex]{\textbf{Method}}
& \multicolumn{2}{c}{\textbf{Segmentation}}
& \multicolumn{2}{c}{\textbf{Depth}}
& \multicolumn{5}{c}{\textbf{Surface Normal}}
& \multirow{3}{*}[-1ex]{$\Delta_b \uparrow$} \\
\cmidrule(lr){2-3}\cmidrule(lr){4-5}\cmidrule(lr){6-10}
& \multirow{2}{*}[-0.4ex]{\textbf{mIoU}$\uparrow$}
& \multirow{2}{*}[-0.4ex]{\textbf{Pix Acc}$\uparrow$}
& \multirow{2}{*}[-0.4ex]{\textbf{Abs Err}$\downarrow$}
& \multirow{2}{*}[-0.4ex]{\textbf{Rel Err}$\downarrow$}
& \multicolumn{2}{c}{\textbf{Angle Distance}}
& \multicolumn{3}{c}{\textbf{Within $t^\circ$}}
& \\
\cmidrule(lr){6-7}\cmidrule(lr){8-10}
& & & & 
& \textbf{Mean}$\downarrow$
& \textbf{Median}$\downarrow$
& \textbf{11.25}$\uparrow$
& \textbf{22.5}$\uparrow$
& \textbf{30}$\uparrow$
& \\
\midrule
HPS$^{\dagger}$          & 54.48 & 75.82 & 0.3839 & 0.1548 & 23.50 & 17.06 & 35.31 & 61.10 & 72.14 & +0.00\% \\
Cross-Stitch$^{\dagger}$ & 53.46 & 75.49 & 0.3804 & 0.1555 & 23.01 & 16.33 & 37.01 & 62.42 & 73.02 & +0.66\% \\
MTAN$^{\dagger}$         & 54.74 & 75.78 & 0.3796 & 0.1549 & 22.97 & 16.30 & 36.91 & 62.63 & 73.32 & +0.77\% \\
NDDR-CNN$^{\dagger}$     & 53.84 & 75.23 & 0.3871 & 0.1560 & 22.60 & 16.07 & 37.67 & 63.43 & 73.92 & +0.91\% \\
VTAGML$^{\dagger}$       & 58.60 & 78.63 & 0.3716 & 0.1525 & 22.05 & 15.70 & 38.14 & 64.28 & 74.50 & +4.70\% \\
DenseMTL$^{\dagger}$     & 56.65 & 77.68 & 0.3569 & 0.1391 & 22.03 & 15.87 & 37.25 & 64.67 & 75.47 & +5.88\% \\
SwinMTL$^{\dagger}$      & 64.23 & 82.78 & 0.2841 & 0.1129 & 18.94 & 13.34 & 43.35 & 71.32 & 80.89 & +19.55\% \\
MultiLoRA$^{\dagger}$        & 64.85 & 83.07 & 0.3113 & 0.1220 & 17.26 & 12.19 & 48.28 & 74.65 & 83.58 & +20.11\% \\
MTSAM$^{\dagger}$ & 65.98 & 83.42 & 0.2898 & 0.1140 & \textbf{16.34} & \textbf{11.33} & \textbf{51.22} & \textbf{77.20} & \textbf{85.51} & +23.93\% \\
\midrule
EW    & 66.64 & 84.21 & 0.2817 & 0.1107 & 17.86 & 12.66 & 45.64 & 73.39 & 82.80 & +22.34\% \\
GLS   & 66.80 & 84.44 & 0.2762 & 0.1084 & 17.23 & 12.12 & 47.50 & 74.96 & 84.03 & \underline{+23.95\%} \\
STCH  & 66.27 & 83.87 & \underline{0.2742} & \underline{0.1077} & 17.25 & 12.12 & 47.61 & 74.79 & 84.88 & +23.81\% \\
FOOPS & \textbf{67.50} & \textbf{84.80} & 0.2787 & 0.1096 & 17.60 & 12.42 & 46.62 & 74.07 & 83.30 & +23.46\% \\
SIMS (Ours)  & \underline{67.12} & \underline{84.70} & \textbf{0.2722} & \textbf{0.1050} & \underline{16.85}   & \underline{11.65}   & \underline{49.35} & \underline{75.85} & \underline{84.54} & +\textbf{25.44\%} \\
\bottomrule
\end{tabular}
\vspace{-6pt}
\end{table*}

% we build our approach on SAM2.1-Large and remove the memory module to better suit the downstream multi-task setting. Following prior work~\cite{wang2025mtsam}, we freeze the image encoder throughout training and only update a parameter-efficient subset together with a modified lightweight downstream prompt encoder. On top of this, we inject task-shared vanilla LoRA adapters into the QKV projections of the attention layers, with rank 32, dropout rate 0.1, and scaling coefficient 0.5. 
\vspace{0.5\baselineskip}

\noindent \textbf{Implementation Details.}
We optimize the model using Adam with a learning rate of $1\times10^{-4}$ and weight decay of $1\times10^{-6}$, while applying a linear warm-up for the first $10\%$ of training steps. We maintain a consistent LoRA configuration across these datasets \footnote{***cy: Check learning rate here. Maybe delete `and learning rate'. `main datasets'? Maybe delete `the main'. \edit{Done}} by setting the rank to 32, the dropout rate to 0.1, and the scaling coefficient to 0.5. Models are trained for 100 epochs on NYU and Cityscapes, and 30 epochs on PASCAL-Context. The batch size of all datasets is 4. For computational efficiency, we use the SAM2.1-Tiny backbone on NYUv2 and accordingly reduce the LoRA rank to 16 while keeping all other settings unchanged in Sec.~\ref{subsec:trans} and Sec.~\ref{subsec:prenorm}. More details are provided in Appendix~\ref{app:mtldetails}.
% on the NYUv2 dataset using SAM2.1-Tiny, and accordingly reduce the LoRA rank to 16 while keeping all other settings unchanged. 
% These choices are made solely for efficiency and do not affect the conclusions of the experiments.

We follow the configurations reported in the original papers for all baseline methods. In particular, for STCH, we adhere to the default setting that applying objective normalization proposed by~\cite{dai2023improvable} to alleviate the mismatch in objective scales across tasks. In addition, SIMS, STCH, and FOOPS all utilize log-sum-exp smoothing. Since~\cite{chen2025efficient} and our subsequent ablation studies indicate that performance is relatively insensitive to the smoothing temperature within a reasonable range, we use a shared smoothing parameter $\tau$ across these methods to ensure a fair comparison.



\begin{table}[!pbht]
\centering
% \vspace{-3pt}
% \vspace{-6pt}
\setlength{\tabcolsep}{2pt}
\renewcommand{\arraystretch}{1}
\caption{Performance on two tasks on the CityScapes dataset.
Results marked with $^{\dagger}$ are copied from prior work. The best results for each task are shown in \textbf{bold}, and the second best are underlined.
$\uparrow(\downarrow)$ indicates that higher (lower) is better.}
\vspace{-4pt}
\label{tab:cityscapes_results}
\resizebox{1.0\columnwidth}{!}{
\begin{tabular}{lcccccc}
\toprule
\multirow{2}{*}{\textbf{Method}} 
& \multicolumn{2}{c}{\textbf{Segmentation}} 
& \multicolumn{2}{c}{\textbf{Depth}} 
& \multirow{2}{*}{$\Delta_b \uparrow$} \\
\cmidrule(lr){2-3}\cmidrule(lr){4-5}
& \textbf{mIoU}$\uparrow$ & \textbf{Pix Acc}$\uparrow$ & \textbf{Abs Err}$\downarrow$ & \textbf{Rel Err}$\downarrow$ &  \\
\midrule
HPS$^{\dagger}$          & 67.40 & 90.92 & 0.0142 & 45.4262 & +0.00\% \\
STL$^{\dagger}$          & 68.13 & 91.28 & 0.0133 & 45.0390 & +2.17\% \\
Cross-Stitch$^{\dagger}$ & 68.01 & 91.29 & 0.0135 & 44.4246 & +2.11\% \\
MTAN$^{\dagger}$         & 68.97 & 91.59 & 0.0136 & 43.7508 & +2.74\% \\
NDDR-CNN$^{\dagger}$     & 68.02 & 91.25 & 0.0137 & 44.8662 & +1.51\% \\
VTAGML$^{\dagger}$       & 73.70 & 93.23 & 0.0138 & 42.8304 & +5.10\% \\
DenseMTL$^{\dagger}$     & 69.75 & 91.45 & 0.0152 & 52.1401 & -4.44\% \\
SwinMTL$^{\dagger}$      & 73.33 & 92.87 & 0.0132 & 38.0720 & +8.54\% \\
MultiLoRA$^{\dagger}$        & 87.23 & 96.67 & 0.0127 & 31.0091 & +19.51\% \\
MTSAM$^{\dagger}$ & 87.45 & 96.80 & 0.0113 & 33.0086 & +20.99\% \\
\midrule
EW & \textbf{88.24} & \textbf{97.02} & 0.0097 & 41.3157 & +19.59\% \\
GLS & 87.79 & 96.88 & \underline{0.0085} & 25.9219 & +29.97\% \\
STCH  & 88.07 & 96.95 & 0.0088 & \textbf{24.5415} & \underline{+30.33\%} \\
FOOPS & \underline{88.21} & \underline{97.00} & 0.0096 & 38.4143 & +21.35\% \\
SIMS (Ours)  & 87.91 & 96.90 & \textbf{0.0084} & \underline{24.8550} & +\textbf{30.78\%} \\
\bottomrule
\end{tabular}
}
\vspace{-10pt}
\end{table}

\begin{table}[!phbt]
\centering
\setlength{\tabcolsep}{2pt}
\renewcommand{\arraystretch}{1}
\caption{Performance on four tasks on the PASCAL-Context dataset. 
Results marked with $^{\dagger}$ are from prior work. The best results for each task are shown in \textbf{bold}, and the second best are underlined.
$\uparrow(\downarrow)$ indicates that higher (lower) is better.}
\vspace{-4pt}
\label{tab:pascal_context_results}
\resizebox{0.9\columnwidth}{!}{
\begin{tabular}{lccccc}
\toprule
\textbf{Method} & \textbf{Seg.}$\uparrow$ & \textbf{H.Parts}$\uparrow$ & \textbf{Sal.}$\uparrow$ & \textbf{Normal}$\downarrow$ & $\Delta_b \uparrow$ \\
\midrule
HPS$^{\dagger}$          & 64.77 & 57.91 & 64.10 & \textbf{14.21} & +0.00\% \\
STL$^{\dagger}$          & 65.14 & 58.58 & 65.02 & 15.94 & -2.25\% \\
Cross-Stitch$^{\dagger}$ & 64.97 & 58.63 & 64.46 & 15.32 & -1.42\% \\
MTAN$^{\dagger}$         & 64.56 & 59.08 & 64.57 & 14.74 & -0.33\% \\
NDDR-CNN$^{\dagger}$     & 65.28 & 59.18 & 65.09 & 15.57 & -1.26\% \\
MultiLoRA$^{\dagger}$    & 72.39 & 67.78 & 71.66 & 20.07 & -0.16\% \\
MTSAM$^{\dagger}$        & 74.13 & 71.04 & 76.28 & 17.10 & \underline{+8.95\%} \\
\midrule
EW                      & 74.83 & 71.92 & 74.36 & 18.42 & +6.53\% \\
GLS                     & 76.36 & \textbf{74.92} & 76.22 & 18.77 & +8.54\% \\
FOOPS                    & \underline{77.56} & 73.56 & 74.78 & \underline{18.17} & +8.89\% \\
STCH                     & 73.66 & 73.18 & \underline{76.97} & 18.18 & +8.06\% \\
SIMS (Ours)        & \textbf{77.60} & \underline{73.88} & \textbf{77.56} & 18.22 & +\textbf{10.04\%} \\
\bottomrule
\end{tabular}
}
\vspace{-6pt}
\end{table}

% \fei{Give a brief description about tasks in each dataset or put this description in the appendix}



% \begin{table}[t]
% \caption{Ablation studies on the impact of $\tau$ on the average performance on the NYUv2 dataset. Higher $\Delta_b$ indicates better overall performance.}
% \label{table:ablation_tau}
% \centering
% \setlength{\tabcolsep}{12pt}
% \renewcommand{\arraystretch}{1}
% \begin{tabular}{c c c}
% \toprule
% Method & $\tau$ & $\Delta_b \uparrow$ \\
% \midrule
% \multirow{6}{*}{SIMS} 
% & $0.1$  & +19.10\% \\
% & $0.5$  & +23.84\% \\
% & $1$    & +24.80\% \\
% & $3$    & +25.06\% \\
% & $5$    & \textbf{+25.44\%} \\
% & $10$   & +24.97\% \\
% \bottomrule
% \end{tabular}
% \end{table}

% \begin{table}[t]
% \caption{The Effect of $\lambda$ on the average performance on the NYUv2 dataset. Higher $\Delta_b$ indicates better overall performance.}
% \label{table:ablation_lambda}
% \centering
% \setlength{\tabcolsep}{8pt}
% \renewcommand{\arraystretch}{1}
% \begin{tabular}{c c c}
% \toprule
% Method & $\lambda$ & $\Delta_b \uparrow$ \\
% \midrule
% \multirow{5}{*}{SIMS}
% & 0 & +25.44\% \\
% & $1.00\times10^{-6}$ & +25.06\% \\
% & $1.00\times10^{-5}$ & \textbf{+25.57\%} \\
% & $1.00\times10^{-4}$ & +25.00\% \\
% & $1.00\times10^{-3}$ & +24.41\% \\
% \bottomrule
% \end{tabular}
% \end{table}

\begin{table*}[!pbht]
\centering
\setlength{\tabcolsep}{6pt}
\renewcommand{\arraystretch}{1}
\caption{Comparison of transformation functions and loss normalization strategies. Performance on three tasks on the NYUv2 dataset. 
The best results are shown in \textbf{bold}, and the second best are underlined.
$\uparrow(\downarrow)$ indicates that higher (lower) is better.}
\vspace{-3pt}
\label{table:trans}
\resizebox{1.85\columnwidth}{!}{
\begin{tabular}{ccccccccccc}
\toprule
\multirow{3}{*}[-1ex]{\textbf{$\psi$}}
& \multicolumn{2}{c}{\textbf{Segmentation}}
& \multicolumn{2}{c}{\textbf{Depth}}
& \multicolumn{5}{c}{\textbf{Surface Normal}}
& \multirow{3}{*}[-1ex]{$\Delta_b \uparrow$} \\
\cmidrule(lr){2-3}\cmidrule(lr){4-5}\cmidrule(lr){6-10}
& \multirow{2}{*}[-0.4ex]{\textbf{mIoU}$\uparrow$}
& \multirow{2}{*}[-0.4ex]{\textbf{Pix Acc}$\uparrow$}
& \multirow{2}{*}[-0.4ex]{\textbf{Abs Err}$\downarrow$}
& \multirow{2}{*}[-0.4ex]{\textbf{Rel Err}$\downarrow$}
& \multicolumn{2}{c}{\textbf{Angle Distance}}
& \multicolumn{3}{c}{\textbf{Within $t^\circ$}}
& \\
\cmidrule(lr){6-7}\cmidrule(lr){8-10}
& & & & 
& \textbf{Mean}$\downarrow$
& \textbf{Median}$\downarrow$
& \textbf{11.25}$\uparrow$
& \textbf{22.5}$\uparrow$
& \textbf{30}$\uparrow$
& \\
\midrule
\multicolumn{11}{c}{Our Method} \\
$\ln(x)$      & \underline{54.57} & 76.72 & \textbf{0.3614} & 0.1476 & \textbf{20.22} & \textbf{14.03} & \textbf{41.48} & \textbf{68.29} & \textbf{78.19} & \textbf{0.00\%} \\
\midrule
\multicolumn{11}{c}{Different Transformation} \\
$x^2$      & 53.15 & 76.26 & 0.3803 & 0.1548 & 21.99 & 15.79 & 37.07 & 63.85 & 74.57 & -5.09\% \\
$x$        & 54.12 & 76.68 & 0.3726 & 0.1519 & 21.24 & 15.03 & 38.90 & 65.76 & 76.12 & -2.80\% \\
$\sqrt{x}$     & 54.21 & 76.80 & 0.3663 & 0.1498 & 20.86 & 14.67 & 39.94 & 66.66 & 76.85 & -1.61\% \\
$asinh(\frac{x}{0.1})$ & 53.75 & \underline{76.81} & \underline{0.3631} & 0.1479 & 20.48 & 14.29 & 40.87 & 67.66 & 77.68 & -0.75\% \\
$asinh(\frac{x}{0.01})$& \textbf{54.99} & \textbf{77.06} & 0.3643 & 0.1478 & \underline{20.29} & 14.12 & 41.26 & \underline{68.18} & \underline{78.12} & \underline{-0.07\%} \\
\midrule
\multicolumn{11}{c}{Loss Normalization} \\
Fix@Epoch 2& 53.72 & 76.64 & 0.3639 & 0.1479 & 20.63 & 14.55 & 40.14 & 67.20 & 77.38 & -1.20\% \\
Epoch-wise&53.94 & 76.39 & \underline{0.3631} & \textbf{0.1472} & 20.25 & \underline{14.10} & \underline{41.36} & 68.15 & \underline{78.12} & -0.38\% \\
EMA      & 53.98 & 76.79 & 0.3645 & \underline{0.1475} & 20.65 & 14.60 & 39.92 & 67.07 & 77.32 & -1.15\% \\

\bottomrule
\end{tabular}
}
\vspace{-3pt}
\end{table*}

% \vspace{-3pt}
\subsection{Main Results}
% \fei{Since you have introduce dataset detail, delete this paragraph}
% We report the main results on NYUv2, Cityscapes, and PASCAL-Context, which differ markedly in both task composition and supervision statistics. NYUv2 mixes semantic segmentation with dense regression objectives such as depth and surface normals, Cityscapes focuses on outdoor driving scenes with dense semantic predictions, and PASCAL-Context features richer semantic labels with a larger label space. These differences induce substantial loss-scale mismatch across tasks due to heterogeneous loss families, label entropy, and varying numbers of valid supervised elements.
% On these datasets, prior multi-task methods~\cite{lopes2023cross} therefore rely on task-specific weighting to compensate for such scale discrepancies, which suggests that an explicitly scale-invariant design can yield a more favorable Pareto trade-off across tasks.

Our results on NYUv2, Cityscapes, and PASCAL-Context are reported in Tables~\ref{table:nyu},~\ref{tab:cityscapes_results} and~\ref{tab:pascal_context_results}.\footnote{***cy: It would be better to underline the second-best result since you are always the second-best result in Table 1. and then explain the underline in the legend. \edit{Done}} For convenience, we directly report some results from prior work~\cite{wang2025mtsam}, which are marked with $^{\dagger}$.
As presented in the results, SIMS achieves state-of-the-art (SOTA) average performance across all three benchmarks. 

Notably, EW consistently yielded the lowest performance among all evaluated scalarization methods. Crucially, the advantage of SIMS extend beyond superior average metrics. The intrinsic scale invariance of SIMS\footnote{***cy: check `scale-invariant formulation' \edit{Done}} enables it to yield a balanced performance profile that avoids the extreme trade-offs where performance on specific task is severely compromised. This superiority is particularly pronounced on the Cityscapes dataset as shown in Table~\ref{tab:cityscapes_results}, where we observe that FOOPS converges to solutions heavily dominated by the segmentation objective, neglecting depth estimation due to the scales of losses. While STCH mitigates this via explicit objective normalization and attains results comparable to SIMS, such heuristic alignment relies on statistical estimation that may face scalability challenges. As the number of tasks increases (Tables~\ref{table:nyu} and ~\ref{tab:pascal_context_results}), these normalization estimates may become increasingly unstable or inaccurate. In contrast, SIMS \edit{exhibits intrinsic scale invariance, eliminating the need for explicit alignment of the loss magnitudes.}\footnote{***cy: check robustness \edit{Done}} 
These quantitative gains are further corroborated by the visualization in Fig.~\ref{fig:toy-scale}, where SIMS maintains consistent optimization trajectories invariant to objective scaling, thereby avoiding the extreme trade-off solutions that plague baseline methods.


\begin{figure}[!tbph]
    \centering
    \vspace{-5pt}
    \includegraphics[width=1\linewidth]{figs/ablation_tau_lambda.pdf}
    \captionsetup{skip=2pt}
    \caption{Ablation study on the hyper-parameters $\tau$ and $\lambda$.}
    \label{fig:ablation_tau_lambda}
    \vspace{-10pt}
\end{figure}


\subsection{Ablation Studies}
\label{sec:ablation}
% \fei{Add 'We conduct ablation studies on the impact of the smoothing parameter $\tau$ and the regularization parameter $\lambda$.'}
We conduct ablation studies on the impact of the smoothing parameter $\tau$ and the regularization parameter $\lambda$.
\vspace{0.5\baselineskip}

\noindent \textbf{Effect of $\tau$.}\footnote{***cy: Why `Sensitivity with respect to X' but not `Effect of X' \edit{Done}}
We perform a sensitivity study on the parameter $\tau$ to examine how varying $\tau$ affects the performance of our method, with all other hyper-parameters fixed to the same settings as in previous experiments. According to the results shown in Fig.~\ref{fig:ablation_tau_lambda}(a), we can see that within a reasonable range (i.e., $[1, 10]$)\footnote{***cy: complete it. \edit{Done}} of $\tau$, our method achieves consistently stable performance, achieving competitive or state-of-the-art overall performance. The full results can be found in Appendix~\ref{appdix:tau}.
\vspace{0.8\baselineskip}

\noindent \textbf{Effect of $\lambda$. }
Recall that our method optimizes the smoothed surrogate objective
$v_{\lambda, \tau}(\theta)=-
    \min_{\theta' \in \mathbb{R}^d}
    h_{\lambda,\tau}(\theta,\theta'),$
where $\lambda \ge 0$ weights the proximal regularization term. This term improves the curvature of the inner minimization and can yield better-conditioned gradients in practice. According to Proposition~\ref{theov}, $\lambda$ also affects the Pareto interpretation of the surrogate: when $\lambda=0$, $v_{0, \tau}(\theta)\le 0$ implies an $\epsilon$-weak Pareto guarantee, whereas $\lambda > 0$ can recover weak Pareto optimality under additional mild conditions. As shown in Fig.~\ref{fig:ablation_tau_lambda}(b), we find SIMS to be insensitive to variations in $\lambda$ across several orders of magnitude, with small positive values sometimes providing modest gains. Nevertheless, to keep the approach tuning-light and to use a single uniform setting across all benchmarks, we set $\lambda =0$ in all experiments. The full results can be found in Appendix~\ref{appdix:lambda}.

\vspace{-5pt}

\subsection{Effect of Transformation Functions}
\label{subsec:trans}
% Beyond logarithmic transformations, motivated by the geometric-mean formulation~\cite{chennupati2019multinet++}, we further investigate a root-based loss aggregation, which empirically exhibits scale-invariant characteristics.
Beyond logarithmic transformations, we further investigate the optimality of our choice for the transformation function $\psi$ and its impact on the resulting multi-task trade-off. Since the transformations $\psi$ is required to be monotonically increasing, we consider (i) a power function family $\psi(x)=x^a$ with $a\in\{0.5,1,2\}$ and (ii) inverse-hyperbolic-sine transforms $\psi(x)=asinh(cx)$ with $c\in\{0.1,0.01\}$.

As shown in Table~\ref{table:trans}, in the power function family, the expansive setting ($a=2$) results in the poorest performance by aggravating discrepancies of loss scale, whereas the compressive setting ($a=0.5$) offers tangible improvements over the linear baseline ($a=1$).\footnote{***cy: improvements compared to whom. \edit{Done}} Turning to the $\text{asinh}$ family, which serves as a smooth numerical approximation to the logarithm in engineering \cite{lupton1999modified}, the setting $c=0.01$ demonstrates remarkable performance, trailing the log baseline by only $\Delta_b = -0.07\%$. Overall, our empirical results demonstrate that the logarithmic transformation consistently outperforms all alternative candidates, confirming it as the optimal design choice.
% compressive choices (e.g., $a=0.5$ and asinh) remain close to the log baseline, whereas the expansive choice $a=2$ amplifies relative scale differences and consistently degrades the overall balance across tasks. 
% This trend aligns with our $\gamma$-scale invariance principle: desirable scalarizations should suppress influence distortion induced by task-wise rescaling (small $\gamma$). In particular, the log transform achieves perfect scale invariance ($\gamma=0$) and thus serves as a natural default in our method.


\subsection{Comparison with Loss Normalization}
\label{subsec:prenorm}
% We interpret pre-normalization as a scale-invariance heuristic that approximates our transformation function by a linear transformation
We view loss normalization as a heuristic approach to scale invariance. To facilitate a direct comparison with SIMS, this strategy can be interpreted as applying a linear transformation to the task losses, where the associated linear coefficients are estimated from training statistics. We consider a Min–Max normalization baseline inspired by previous works~\cite{dai2023improvable,liu2019end,lakkapragada2023mitigating} that fixes the minimum to zero and varies the estimation scope of the maximum via three schemes: (i) a global estimate set to the mean loss at the second epoch, (ii) an epoch-wise estimate using the last epoch’s mean loss, and (iii) a step-wise estimate computed by an exponential moving average (EMA) over steps.

As shown in Table~\ref{table:trans}\footnote{***cy: As shown in Table XXX \edit{Done}}, all three variants substantially underperform compared to SIMS (i.e., $\Delta_b$ of $-1.20\%$, $-0.38\%$, and $-1.15\%$, respectively), with the epoch-wise variant being the strongest among them. This trend shows that purely statistic-driven normalization\footnote{***cy: `statistic-driven linear rescaling'? check. \edit{Done}} is fragile to the choice of reference, whereas SIMS removes the need for such auxiliary scale estimates by embedding scale invariance directly into the merit function construction.

\begin{figure}[!hpbt]
  \vspace{-5pt}
  \centering  \includegraphics[width=0.88\linewidth]{figs/var_nogrid.pdf}
  \captionsetup{skip=1pt}
  \caption{Trajectory Variance across Scales
  }
  \label{fig:var}
  \vspace{-10pt}
\end{figure}
\begin{add}

\subsection{Synthetic Experiments for Scale Invariance}
To further evaluate whether SIMS is invariant to task-wise loss rescaling, we conduct a quantitative synthetic experiment based on the motivating example in Section~\ref{sec:motivation}.
Here, we examine whether each scalarization method behaves consistently and produces similar trajectories when the relative scales of task losses are artificially changed.
% while the underlying Pareto front remains unchanged.
% We fix the initialization at $(0,0)$, use the distance to the Pareto front with threshold 0.05 as the stop criterion, and evaluate different methods under three fixed seeds and seven rescaling settings: 1:1, 1:1.25, 1.25:1, 1:10, 10:1, 100:1, and 1:100. As the quantitative metric, we compute the variance of the final converged solutions across the 7 rescaling settings. The results in the following table show that SIMS and GLS achieve the smallest variance, while other methods exhibit much larger variance. This provides direct empirical evidence that both SIMS and GLS are substantially more scale-invariant than existing baselines, consistent with the claim we make.
Specifically, we fix the initialization at $(0,0)$ and stop once the distance to the Pareto front is below $0.05$. We evaluate each method under three fixed random seeds and seven rescaling settings: $1:1$, $1:1.25$, $1.25:1$, $1:10$, $10:1$, $100:1$, and $1:100$. These task-wise rescaling settings only change the numerical magnitudes of the objectives but do not alter the Pareto front.

% , a scale-invariant method is expected to converge to similar final solutions across these settings.
% As the quantitative metric, we compute the variance of the final converged solutions across the seven rescaling settings. A smaller variance indicates that the method is less affected by arbitrary changes in objective scales and therefore has stronger scale-invariant behavior. We analyze this consistency from both the trajectory and final-solution perspectives. As shown in Fig.~\ref{fig:var}, the trajectory variances of EW, FOOPS, and STCH increase rapidly during optimization, suggesting that their optimization paths are strongly affected by task-wise rescaling. In contrast, GLS and SIMS keep the variance close to zero, indicating nearly consistent trajectories across different positive rescaling settings. At the final-solution level, Table~\ref{tab:variance_comparison} shows that SIMS and GLS achieve the smallest variances, while EW, STCH, and FOOPS exhibit much larger variances. These results provide direct empirical evidence that SIMS is substantially more scale-invariant than existing baselines, and further support the scale-invariant property of the proposed logarithmic transformation-induced merit function.
As the quantitative metric, we compute the variance across the seven rescaling settings, where a smaller variance indicates stronger scale-invariant behavior. We analyze the results from both trajectory and final solution perspectives. As shown in Fig.~\ref{fig:var}, the trajectory variances of EW, FOOPS, and STCH increase rapidly during optimization, indicating that their optimization trajectories are strongly affected by task-wise rescaling. In contrast, GLS and SIMS keep the variance close to zero, suggesting nearly consistent trajectories across rescaling settings. At the final solutions, Table~\ref{tab:variance_comparison} shows that SIMS and GLS achieve the smallest variances, while EW, STCH, and FOOPS exhibit much larger variances. These results provide direct empirical evidence that SIMS is substantially more scale-invariant than existing baselines, further supporting the scale-invariant property of the proposed logarithmic transformation-induced merit function.

\end{add}


\begin{table}[t]
\centering
\caption{Variance of final converged solutions under different task-wise rescaling settings. Lower ($\downarrow$) is better.}
\label{tab:variance_comparison}
\begin{tabular}{lccc}
\toprule
Method & Var. of $\ell_1\downarrow$ & Var. of $\ell_2\downarrow$ & Average Var.$\downarrow$ \\
\midrule
EW     & 57.322545 & 57.322545 & 57.322545 \\
FOOPS  & 57.318462 & 57.318462 & 57.318462 \\
STCH   & 50.89209  & 60.6234   & 55.757745 \\
GLS    & 0.020195  & 0.020195  & 0.020195 \\
% SVFair & 7.881019  & 11.375863 & 9.628441 \\
SIMS (Ours)   & 0.020105  & 0.020105  & 0.020105 \\
\bottomrule
\end{tabular}
\vspace{-10pt}
\end{table}
% \subsection{training loss}training loss to reflect the scale-free property


\section{Conclusion}
In this paper, we address the sensitivity of existing merit-function-based scalarization methods to the scales of task loss in multi-task learning. Specifically, we propose SIMS, a Scale-Invariant Merit-function-based Scalarization method, \edit{which employs a logarithmic transformation-induced merit function to convert the multi-objective problem into a single objective that is insensitive to loss magnitudes}. We theoretically prove that imposing scale invariance uniquely determines the transformation function to be logarithmic. Furthermore, we show that our formulation preserves weak Pareto optimality and admits a smooth surrogate with controllable approximation error. We demonstrate that the resulting surrogate objective can be efficiently solved via the TTGDA algorithm with established convergence guarantees. Extensive experiments on standard benchmarks demonstrate the effectiveness of SIMS. 

%%
%% The acknowledgments section is defined using the "acks" environment
%% (and NOT an unnumbered section). This ensures the proper
%% identification of the section in the article metadata, and the
%% consistent spelling of the heading.
\begin{acks}
This work was supported by National Natural Science Foundation of China under Grant no. 62136005 and Shenzhen fundamental research program JCYJ20250604144724032.
\end{acks}

%%
%% The next two lines define the bibliography style to be used, and
%% the bibliography file.
\bibliographystyle{ACM-Reference-Format}
\bibliography{reference}

\clearpage
%%
%% If your work has an appendix, this is the place to put it.
\appendix

\section{Proofs}
\subsection{Proofs in Section~\ref{section:formulation}}
\begin{proof}[Proof of Theorem~\ref{theo:u0_merit}]
Let $\theta\in \mathbb{R}^d$. By the definition Eq.~\eqref{eq:ubar} of $\bar{u}$, we have
\begin{align*}
\bar{u}(\theta)
&=
\sup_{\theta'\in \mathbb{R}^d}
\min_{i\in[m]}
\bigl\{\psi(\mathcal{L}_i(\theta))-\psi(\mathcal{L}_i(\theta'))\bigr\}\\
&\geq
\min_{i\in[m]}
\bigl\{\psi(\mathcal{L}_i(\theta))-\psi(\mathcal{L}_i(\theta'))\bigr\}
=0, 
\end{align*}
% \footnote{***cy:check here.}
which implies $\bar{u}(\theta)\ge 0$ for all $\theta\in \mathbb{R}^d$.

On the other hand, again by the definition Eq.~\eqref{eq:ubar} of $\bar{u}$, we obtain
\[
\bar{u}(\theta)=0
\;\Longleftrightarrow\;
\min_{i\in[m]}
\bigl\{\psi(\mathcal{L}_i(\theta))-\psi(\mathcal{L}_i(\theta'))\bigr\}\le 0,
\quad \forall\,\theta'\in \mathbb{R}^d.
\]
Hence, there does not exist any $\theta'\in \mathbb{R}^d$ such that
\[
\psi(\mathcal{L}_i(\theta))-\psi(\mathcal{L}_i(\theta'))>0,
\quad \forall\, i\in[m],
\]
or equivalently,
\[
\psi(\mathcal{L}_i(\theta'))<\psi(\mathcal{L}_i(\theta)),
\quad \forall\, i\in[m].
\]
Since $\psi(\cdot)$ a strictly increasing function, we have 
\[
\mathcal{L}_i(\theta')<\mathcal{L}_i(\theta),
\quad \forall\, i\in[m],
\]
which means that $\theta$ is weakly Pareto optimal for problem~\eqref{eq:1} by definition.
\end{proof}

\subsection{Proofs in Section~\ref{sec5}}
\begin{proof}[Proof of Theorem~\ref{thm:log_characterization}]
Fix any task $i$ and consider a special case in which the inner minimum is always attained at $i$ (e.g., other tasks have constant losses). Then the surrogate locally reduces to the difference form
\[
\bar u(\theta)=\psi(x)-\psi(y),
\]
for arbitrary positive values $x=\mathcal{L}_i(\theta)$ and
$y=\mathcal{L}_i(\theta')$. After rescaling by $c>0$, we obtain
\[
\bar u_c(\theta)=\psi(c x)-\psi(c y).
\]
By the assumed invariance property, we have that
$\bar u_c(\theta)-\bar u(\theta)=K(c)$ must be independent of $\theta$. Therefore, for all $x,y,c>0$,
\[
\psi(c x)-\psi(c y)=\psi(x)-\psi(y).
\]
Setting $y=1$ and defining $f(t)=\psi(t)-\psi(1)$ gives
\[
f(c x)=f(x)+f(c),\qquad \forall x,c>0,
\]
which is the multiplicative Cauchy functional equation
$f(xy)=f(x)+f(y)$.
Since $f$ is differentiable, its only solutions are $f(x)=a\ln x$ (~\citet{kuczma2009introduction}). Hence $\psi(x)=a\ln x+b$. Since $\psi$ is strictly increasing, we must have $a>0$.
\end{proof}


\subsection{Auxiliary Lemmas}
\label{app:aux_lemmas}
For convenience, we define the merit function $u_\lambda(\theta)$ and restate
the smoothed merit function $v_{\lambda,\tau}(\theta)$ below.

\begin{equation}
\label{eq:u_lambda}
u_\lambda(\theta) := \sup_{\theta'\in \mathbb{R}^d} \min_{i\in[m]}
\left\{\tilde{\mathcal{L}}_i(\theta)-\tilde{\mathcal{L}}_i(\theta')- \frac{\lambda}{2}\|\theta-\theta'\|^2\right\},
\end{equation}
and
\begin{equation}
\label{eq:v_lambda_tau}
v_{\lambda,\tau}(\theta) :=
-\min_{\theta'\in \mathbb{R}^d}\left\{\tau \ln
\left(\sum_{i=1}^m \exp\!\left(
\frac{\tilde{\mathcal{L}}_i(\theta')-\tilde{\mathcal{L}}_i(\theta)}{\tau}
\right)\right)+\frac{\lambda}{2}\|\theta-\theta'\|^2\right\}.\notag 
\end{equation}
% Correspondingly, we define the auxiliary function
% $h_{\lambda,\tau}(\theta,\theta')$ for analysis.
% Note that
% $
% v_{\lambda,\tau}(\theta)
% =
% -
% \min_{\theta'\in \mathbb{R}^d}
% h_{\lambda,\tau}(\theta,\theta').
% $

% \begin{equation}
% \label{eq:h_lambda_tau}
% h_{\lambda,\tau}(\theta,\theta')
% :=
% \tau \ln
% \left(
% \sum_{i=1}^m
% \exp\!\left(
% \frac{\tilde{\mathcal{L}}_i(\theta')-\tilde{\mathcal{L}}_i(\theta)}{\tau}
% \right)
% \right)
% +
% \frac{\lambda}{2}\|\theta-\theta'\|^2 .
% \end{equation}


\begin{lemma}[Theorems~3.1 and~3.3 in~\cite{tanabe2024new}]
\label{lem:u_lambda_merit}
For $\lambda\ge 0$, consider the merit function $u_\lambda(\theta)$ defined in Eq.~\eqref{eq:u_lambda}. Then $u_\lambda(\theta)\ge 0$ for all $\theta\in \mathbb{R}^d$. Moreover, $u_0(\theta)=0$ if and only if $\theta$ is weakly Pareto optimal. For $\lambda>0$, if $\theta$ is weakly Pareto optimal, then $u_\lambda(\theta)=0$; furthermore, if $\tilde{\mathcal{L}}_i(\theta)$ is convex for all $i\in[m]$, then $u_\lambda(\theta)=0$ implies that $\theta$ is weakly Pareto optimal.
\end{lemma}

% \begin{proposition}[Smoothness implies weak convexity]
% \label{prop:smooth_implies_weak_convex}
% If a locally Lipschitz function $f$ is $\ell_{f,1}$-smooth, then it is also
% $\ell_{f,1}$-weakly convex.
% \end{proposition}

\begin{lemma}[Log-sum-exp preserves weak convexity]
\label{lem:lse_weak_convex}
Let $\tilde{\mathcal{L}}_i(\theta)$, $i\in[m]$, be weakly convex with modulus
$\mu_i\in\mathbb{R}$. Let $\bar{\mu}=\max_{i\in[m]} \mu_i$. Then
\[
\ln\!\left(\sum_{i=1}^m e^{\tilde{\mathcal{L}}_i(\theta)}\right)
\]
is weakly convex with modulus $\bar{\mu}$.
\end{lemma}
\begin{proof}
By definition, since $\bar{\mu}=\max_{i\in[m]} \mu_i$, the function
\[
\tilde{\mathcal{L}}_i(\theta)+\frac{\bar{\mu}}{2}\|\theta\|^2
\]
is convex for all $i\in[m]$.
Because the log-sum-exp function preserves convexity, it follows that
\[
\ln\!\left(
\sum_{i=1}^m
e^{\tilde{\mathcal{L}}_i(\theta)+\frac{\bar{\mu}}{2}\|\theta\|^2}
\right)
\]
is convex.
Rearranging the above expression yields
\begin{align}
\ln\!\left(
\sum_{i=1}^m
e^{\tilde{\mathcal{L}}_i(\theta)+\frac{\bar{\mu}}{2}\|\theta\|^2}
\right)
=
\ln\!\left(\sum_{i=1}^m e^{\tilde{\mathcal{L}}_i(\theta)}\right)
+
\frac{\bar{\mu}}{2}\|\theta\|^2 \notag .
\end{align}
Hence,
\(
\ln\!\left(\sum_{i=1}^m e^{\tilde{\mathcal{L}}_i(\theta)}\right)
\)
is $\bar{\mu}$-weakly convex.
\end{proof}



% \begin{proof}[Proof of Corollary~\ref{cor:strict_convexity_h}]
% By $\mu_i$-weak convexity, $\frac{1}{\tau}\tilde{\mathcal{L}}_i(\theta')$ is
% $\frac{\mu_i}{\tau}$-weakly convex with respect to $\theta'$.
% Combining this with Lemma~\ref{lem:lse_weak_convex}, the function
% \[
% \tau\ln\!\left(
% \sum_{i=1}^m
% e^{\frac{\tilde{\mathcal{L}}_i(\theta)-\tilde{\mathcal{L}}_i(\theta')}{\tau}}
% \right)
% \]
% is $\bar{\mu}$-weakly convex with respect to $\theta'$, where
% $\bar{\mu}=\max_{i\in[m]} \mu_i$.
% Since $\lambda+\bar{\mu}>0$, it follows that
% $h_{\lambda,\tau}(\theta,\theta')$ is strictly convex with respect to
% $\theta'$, and hence admits a unique minimizer.
% \end{proof}

% \begin{lemma}
% \label{lem:existence_v}
% If $\tilde{\mathcal{L}}_i(\theta)$, $i\in[m]$, is continuous and $\mu_i$-weakly convex, and $\lambda\ge -\min_{i\in[m]} \mu_i$, then there exists $\theta\in \mathbb{R}^d$ such that
% \[
% v_{\lambda,\tau}(\theta)=-\tau\ln m .
% \]
% \end{lemma}

% \begin{proof}
% For the second claim, again from~\cite{beck2012smoothing}, we have for any
% $\theta,\theta'\in \mathbb{R}^d$,
% \begin{align}
% \tau \ln\!\left(\sum_{i=1}^m
% \exp\!\left(\frac{\tilde{\mathcal{L}}_i(\theta')-\tilde{\mathcal{L}}_i(\theta)}{\tau}\right)\right)
% +\frac{\lambda}{2}\|\theta-\theta'\|^2
% &\le
% \tau\ln m
% +
% \max_{i\in[m]}\{\tilde{\mathcal{L}}_i(\theta')-\tilde{\mathcal{L}}_i(\theta)\}
% +\frac{\lambda}{2}\|\theta-\theta'\|^2 .
% \end{align}
% Since taking $\min_{\theta'\in \mathbb{R}^d}$ preserves inequalities, it follows that
% \begin{align}
% \min_{\theta'\in \mathbb{R}^d}
% \Bigg\{
% \tau \ln\!\left(\sum_{i=1}^m
% \exp\!\left(\frac{\tilde{\mathcal{L}}_i(\theta')-\tilde{\mathcal{L}}_i(\theta)}{\tau}\right)\right)
% +\frac{\lambda}{2}\|\theta-\theta'\|^2
% \Bigg\}
% &\le
% \tau\ln m
% +
% \min_{\theta'\in \mathbb{R}^d}
% \Bigg\{
% \max_{i\in[m]}\{\tilde{\mathcal{L}}_i(\theta')-\tilde{\mathcal{L}}_i(\theta)\}
% +\frac{\lambda}{2}\|\theta-\theta'\|^2
% \Bigg\}.
% \label{eq:min_preserve_ineq}
% \end{align}
% By the definition of $v_{\lambda,\tau}$ in~\eqref{eq:v_lambda_tau},
% we thus have
% \begin{align}
% v_{\lambda,\tau}(\theta)
% &=
% -
% \min_{\theta'\in \mathbb{R}^d}
% \Bigg\{
% \tau \ln\!\left(\sum_{i=1}^m
% \exp\!\left(\frac{\tilde{\mathcal{L}}_i(\theta')-\tilde{\mathcal{L}}_i(\theta)}{\tau}\right)\right)
% +\frac{\lambda}{2}\|\theta-\theta'\|^2
% \Bigg\}
% \nonumber\\
% &\ge
% -
% \min_{\theta'\in \mathbb{R}^d}
% \Bigg\{
% \max_{i\in[m]}\{\tilde{\mathcal{L}}_i(\theta')-\tilde{\mathcal{L}}_i(\theta)\}
% +\frac{\lambda}{2}\|\theta-\theta'\|^2
% \Bigg\}
% -\tau\ln m
% \nonumber\\
% &=
% u_\lambda(\theta)-\tau\ln m
% \;\ge\;
% -\tau\ln m,
% \label{eq:v_lower_bound}
% \end{align}
% where the last inequality holds since $u_\lambda(\theta)\ge 0$
% (Lemma~\ref{lem:u_lambda_merit}).
% Furthermore, there exists $\theta\in \mathbb{R}^d$ such that
% $v_{\lambda,\tau}(\theta)=-\tau\ln m$ by Lemma~\ref{lem:existence_v}.
% Therefore,
% \(
% \min_{\theta\in \mathbb{R}^d} v_{\lambda,\tau}(\theta) = -\tau\ln m
% \),
% which completes the proof.
% \end{proof}

% \begin{proof}[Proof of Lemma~\ref{lem:existence_v}]
% From Corollary~\ref{cor:strict_convexity_h}, since
% $\lambda\ge -\min_{i\in[m]} \mu_i$, the function
% $h_{\lambda,\tau}(\theta,\theta')$ is convex with respect to $\theta'$.
% Let $P$ denote the indicator function of $\Theta$.
% Then
% \[
% 0\in\nabla_{\theta'} h_{\lambda,\tau}(\theta,\theta')
% +\partial P(\theta')
% \quad\Longleftrightarrow\quad
% \theta'=\arg\min_{\theta'\in \mathbb{R}^d} h_{\lambda,\tau}(\theta,\theta').
% \]
% \end{proof}

\begin{lemma}[Smoothness of log-sum-exp smoothing]\label{lem:lse_smoothness}
Suppose that for all $i\in[m]$,
$\tilde{\mathcal{L}}_i$ is $\zeta_i$-smooth and $\ell_i$-Lipschitz continuous on $\mathbb{R}^d$.
Define
\begin{align}
    g_{\mu}(\theta)=\mu\ln\!\Big(\sum_{i=1}^{m}\exp\big(\tilde{\mathcal{L}}_i(\theta)/\mu\big)\Big),
    \quad \mu>0 \notag.
\end{align}
Then $g_{\mu}$ is continuously differentiable on $\mathbb{R}^d$, and its gradient is Lipschitz on $\mathbb{R}^d$ with
\begin{align}
    \|\nabla g_{\mu}(\theta)\!-\!\nabla g_{\mu}(\theta')\|
    \le\!
    \Big(\frac{\bar{\ell}^2}{\mu}+\bar{\zeta}\Big)
    \|\theta-\theta'\|,\, \forall~\theta,\theta'\in \!\mathbb{R}^d \notag.
\end{align} %\footnote{***cy:check here.}
Equivalently, $g_{\mu}$ is $\Big(\frac{\bar{\ell}^2}{\mu}+\bar{\zeta}\Big)$-smooth on $\mathbb{R}^d$, where $\bar{\ell}=\max_{i}\ell_i$ and $\bar{\zeta}=\max_{i}\zeta_i$.
\end{lemma}
\begin{proof}
Define $\boldsymbol{\tilde{\mathcal{L}}}(\theta):=(\tilde{\mathcal{L}}_1(\theta),\ldots,\tilde{\mathcal{L}}_m(\theta))^\top$ and
\[
h_\mu(z):=\mu\ln\Big(\sum_{i=1}^m e^{z_i/\mu}\Big),\qquad z\in\mathbb{R}^m,
\]
so that $g_\mu(\theta)=h_\mu(\boldsymbol{\tilde{\mathcal{L}}}(\theta))$. Let $p(z):=\nabla h_\mu(z)$ be the softmax vector, i.e.,
$p_i(z)=\exp(z_i/\mu)/\sum_{j}\exp(z_j/\mu)$, hence $\|p(z)\|_1=1$.
Moreover, $h_\mu$ is $(1/\mu)$-smooth:
\begin{equation}\label{eq:hmu_smooth_icml}
\|\nabla h_\mu(z)-\nabla h_\mu(z')\|\le \frac{1}{\mu}\|z-z'\|,\qquad \forall z,z'\in\mathbb{R}^m .
\end{equation}

By the chain rule, with $J_{\boldsymbol{\tilde{\mathcal{L}}}}(\theta)=(\nabla \tilde{\mathcal{L}}_1(\theta),\ldots,\nabla \tilde{\mathcal{L}}_m(\theta))$,
\[
\nabla g_\mu(\theta)=J_{\boldsymbol{\tilde{\mathcal{L}}}}(\theta)\,\nabla h_\mu(\boldsymbol{\tilde{\mathcal{L}}}(\theta))
=\sum_{i=1}^m p_i(\boldsymbol{\tilde{\mathcal{L}}}(\theta))\,\nabla \tilde{\mathcal{L}}_i(\theta).
\]
For any $\theta,\theta'\in \mathbb{R}^d$, we decompose
\begin{align*}
\|\nabla g_\mu(\theta)-\nabla g_\mu(\theta')\|
\le& \|J_{\boldsymbol{\tilde{\mathcal{L}}}}(\theta)\|\,\|\nabla h_\mu(\boldsymbol{\tilde{\mathcal{L}}}(\theta))-\nabla h_\mu(\boldsymbol{\tilde{\mathcal{L}}}(\theta'))\| \\ 
&+ \|(J_{\boldsymbol{\tilde{\mathcal{L}}}}(\theta)-J_{\boldsymbol{\tilde{\mathcal{L}}}}(\theta'))\,p(\boldsymbol{\tilde{\mathcal{L}}}(\theta'))\|.
\end{align*}
% \footnote{***cy:check here.}

For the first term, From Eq.~\eqref{eq:hmu_smooth_icml}, we have
\[
\|J_{\boldsymbol{\tilde{\mathcal{L}}}}(\theta)\|\,\|\nabla h_\mu(\boldsymbol{\tilde{\mathcal{L}}}(\theta))-\nabla h_\mu(\boldsymbol{\tilde{\mathcal{L}}}(\theta'))\|
\le \frac{\|J_{\boldsymbol{\tilde{\mathcal{L}}}}(\theta)\|}{\mu}\,\|\boldsymbol{\tilde{\mathcal{L}}}(\theta)-\boldsymbol{\tilde{\mathcal{L}}}(\theta')\|.
\]
Since each $\tilde{\mathcal{L}}_i$ is $\ell_i$-Lipschitz on $\mathbb{R}^d$, we have
$\|\nabla \tilde{\mathcal{L}}_i(\theta)\|\le \ell_i$ and $\|\boldsymbol{\tilde{\mathcal{L}}}(\theta)-\boldsymbol{\tilde{\mathcal{L}}}(\theta')\|\le (\max_i \ell_i)\|\theta-\theta'\|$.
Thus $\|J_{\boldsymbol{\tilde{\mathcal{L}}}}(\theta)\|\le \max_i \ell_i$, and the first term is bounded by
$\frac{\max_i \ell_i^2}{\mu}\|\theta-\theta'\|$.

For the second term, using $\|p(\boldsymbol{\tilde{\mathcal{L}}}(\theta'))\|_1=1$ and the $\zeta_i$-smoothness of $\tilde{\mathcal{L}}_i$,
\begin{align}
\|(J_{\boldsymbol{\tilde{\mathcal{L}}}}(\theta)\!-\!J_{\boldsymbol{\tilde{\mathcal{L}}}}(\theta'))\,p(\boldsymbol{\tilde{\mathcal{L}}}(\theta'))\|&=\!
\Big\|\sum_{i=1}^m p_i(\boldsymbol{\tilde{\mathcal{L}}}(\theta'))\big(\nabla \tilde{\mathcal{L}}_i(\theta)-\nabla \tilde{\mathcal{L}}_i(\theta')\big)\Big\| \notag\\
&\le \max_i \zeta_i \,\|\theta-\theta'\|\notag.
\end{align}
% \[
% \|(J_{\boldsymbol{\tilde{\mathcal{L}}}}(\theta)-J_{\boldsymbol{\tilde{\mathcal{L}}}}(\theta'))\,p(\boldsymbol{\tilde{\mathcal{L}}}(\theta'))\|
% =
% \Big\|\sum_{i=1}^m p_i(\boldsymbol{\tilde{\mathcal{L}}}(\theta'))\big(\nabla \tilde{\mathcal{L}}_i(\theta)-\nabla \tilde{\mathcal{L}}_i(\theta')\big)\Big\|
% \le \max_i \zeta_i \,\|\theta-\theta'\|.
% \]
% \footnote{***cy:check here.}
Combining the two bounds proves
\[
\|\nabla g_\mu(\theta)-\nabla g_\mu(\theta')\|
\le
\Big(\frac{\max_i \ell_i^2}{\mu}+\max_i \zeta_i\Big)\,\|\theta-\theta'\|,
\qquad \forall \theta,\theta'\in \mathbb{R}^d.
\]
Let $\bar{\ell}=\max_{i}\ell_i$ and $\bar{\zeta}=\max_{i}\zeta_i$. Then, $g_{\mu}$ is $\Big(\frac{\bar{\ell}^2}{\mu}+\bar{\zeta}\Big)$-smooth on $\mathbb{R}^d$.
\end{proof}

\subsection{Proofs in Section~\ref{sec:app}}
\begin{proof}[Proof of Proposition~\ref{prop:smooth_approx}]
The first claim follows directly from Proposition~\ref{prop:bounded_approx} by taking the limit $\tau \rightarrow 0$, and we omit the proof for brevity. For the second claim, we define, for $\theta, \theta'\in \mathbb{R}^d$,
\[
\phi_i(\theta, \theta'):=\tilde{\mathcal{L}}_i(\theta')-\tilde{\mathcal{L}}_i(\theta),\qquad i\in[m].
\]
Since $\tilde{\mathcal{L}}_i$ is $\ell_i$-Lipschitz and $\zeta_i$-smooth on $\mathbb{R}^d$, $\phi_i$ is $\ell_i$-Lipschitz and $\zeta_i$-smooth w.r.t $\theta, \theta'\in \mathbb{R}^d$ as well.
Consider the log-sum-exp term
\[
h_{0,\tau}(\theta,\theta')=\tau\ln\!\Big(\sum_{i=1}^{m}\exp(\phi_i(\theta)/\tau)\Big).
\]
Applying Lemma~\ref{lem:lse_smoothness} with $\mu=\tau$ and $\{\phi_i\}$ gives that $h_{0,\tau}(\theta,\theta')$ is $\big(\frac{\bar{\ell}^2}{\tau}+\bar{\zeta}\big)$-smooth w.r.t $\theta, \theta'\in \mathbb{R}^d$. Since  $v_{0, \tau}(\theta) =-\min_{\theta' \in \mathbb{R}^d} h_{0,\tau}(\theta,\theta')$, we can also conclude that  $v_{0, \tau}(\theta)$ is also $\big(\frac{\bar{\ell}^2}{\tau}+\bar{\zeta}\big)$-smooth w.r.t $\theta$. Let $C=\bar{\zeta}$ and $\alpha = {\ell}^2$, we arrive the conclusion.
\end{proof}

\begin{proof}[Proof of Proposition~\ref{prop:bounded_approx}]
By the property of the log-sum-exp function (e.g.,~\cite{beck2012smoothing}), and since taking $\min$ preserves inequalities, we obtain
\begin{align}
\label{eq:prop1_bound}
u_\lambda(\theta) - \tau \ln m
\;\le\;
v_{\lambda,\tau}(\theta)
\;\le\;
u_\lambda(\theta),
\qquad \forall\,\theta\in \mathbb{R}^d .
\end{align}
When $\lambda=0,$ $u_0(\theta)= \bar{u}(\theta)$, we have 
$$v_{0,\tau}(\theta) \le \bar{u}(\theta)\leq v_{0,\tau}(\theta) + \tau\ln m. $$
Considering Eq.~\eqref{eq:prop1_bound}, it also implies that as $\tau\rightarrow 0$, $v_{\lambda,\tau}(\theta)$ uniformly converges to $u_\lambda(\theta)$. Also recall from Lemma~\ref{lem:u_lambda_merit} that $\theta$ is weakly Pareto optimal if and only if $u_0(\theta)=0$. Therefore, $\theta$ is weakly Pareto optimal if and only if $\lim_{\tau\rightarrow 0} v_{0,\tau}(\theta)=0$. which completes the proof.
\end{proof}



\begin{proof}[Proof of Lemma~\ref{lemmsc}]
By $\mu_i$-weak convexity, $\frac{1}{\tau}\tilde{\mathcal{L}}_i(\theta')$ is $\frac{\mu_i}{\tau}$-weakly convex with respect to $\theta'$. Therefore, combining with Lemma~\ref{lem:lse_weak_convex}, the function
\[
\tau\ln\!\left(
\sum_{i=1}^m \exp\!\left(\frac{\tilde{\mathcal{L}}_i(\theta)-\tilde{\mathcal{L}}_i(\theta')}{\tau}\right)
\right)
\]
is $\bar{\mu}$-weakly convex with respect to $\theta'$, where $\bar{\mu}=\max_{i\in[m]}\mu_i$. Since $\lambda+\max_{i\in[m]}\mu_i=\lambda+\bar{\mu}>0$, the function $h_{\lambda,\tau}(\theta,\theta')$ is strictly convex with respect to $\theta'$. Hence the minimizer of $\min_{\theta'\in \mathbb{R}^d} h_{\lambda,\tau}(\theta,\theta')$ is unique.
\end{proof}

\begin{proof}[Proof of Proposition~\ref{theov}]
\noindent\textbf{(1)} For the first claim, by Eq.~\eqref{eq:prop1_bound} and $u_\lambda(\theta)\geq 0$, we have $v_{\lambda,\tau}(\theta) \ge -\tau \ln m$ for any $\theta \in \mathbb{R}^d$. 

\noindent\textbf{(2)} For the second claim, by Lemma~\ref{lem:u_lambda_merit}, if $\theta$ is weakly Pareto optimal, then $u_{\lambda}(\theta)=0$. Furthermore, by Proposition~\ref{prop:bounded_approx}, we have $u_{\lambda}(\theta)\ge v_{\lambda,\tau}(\theta)$, which implies $v_{\lambda,\tau}(\theta)\le 0$.
% For the first claim, by Eq.~\eqref{eq:prop1_bound} and $u_\lambda(\theta)\geq 0$, we have $v_{\lambda,\tau}(\theta) \ge -\tau \ln m$ for any $\theta \in \mathbb{R}^d$. 

% Lemma~\ref{lem:u_lambda_merit}, if $\theta$ is weakly Pareto optimal, then $u_{\lambda}(\theta)=0$. Furthermore, by Proposition~\ref{prop:bounded_approx}, we have $u_{\lambda}(\theta)\ge v_{\lambda,\tau}(\theta)$, which implies $v_{\lambda,\tau}(\theta)\le 0$.


% For the first claim, by Lemma~\ref{lem:u_lambda_merit}, if $\theta$ is weakly Pareto optimal, then $u_{\lambda}(\theta)=0$. Furthermore, by Proposition~\ref{prop:bounded_approx}, we have $u_{\lambda}(\theta)\ge v_{\lambda,\tau}(\theta)$, which implies $v_{\lambda,\tau}(\theta)\le 0$.

% Conversely, we prove the second claim under the two conditions.
\noindent\textbf{(3) $\lambda=0$.} If $v_{0,\tau}(\theta)\le 0$, then by Proposition~\ref{prop:bounded_approx} we have
\begin{equation}
\label{eq:prop2_casea_1}
\bar{u}(\theta) \le v_{0,\tau}(\theta)+\tau\ln m \le \tau\ln m.
\end{equation}
By the definition of $\bar{u}$, for all $\hat{\theta}\in \mathbb{R}^d$, it holds that
\begin{equation}
\label{eq:prop2_casea_2}
\min_{i\in[m]} \bigl\{\tilde{\mathcal{L}}_i(\theta)-\tilde{\mathcal{L}}_i(\hat{\theta})\bigr\} \le \bar{u}(\theta) \le \tau\ln m.
\end{equation}
Equivalently, there does not exist any $\hat{\theta}\in \mathbb{R}^d$ such that
\[
\tilde{\mathcal{L}}_i(\hat{\theta}) < \tilde{\mathcal{L}}_i(\theta) - \tau\ln m, \qquad \forall i\in[m].
\]
Since $\psi(\cdot)$ is strictly increasing, this implies that there is no $\hat{\theta}\in \mathbb{R}^d$ satisfying
\[
\mathcal{L}_i(\hat{\theta})< \mathcal{L}_i(\theta)-\epsilon, \qquad \forall i\in[m], \quad\text{with }\epsilon=\tau\ln m,
\]
i.e., $\theta$ is $\epsilon$-weakly Pareto optimal.

\noindent\textbf{(4) $\lambda>0$.}
If $v_{\lambda,\tau}(\theta)\le -\tau\ln m$, then Proposition~\ref{prop:bounded_approx} yields
\begin{align}
0 \le u_{\lambda}(\theta) \le v_{\lambda,\tau}(\theta)+\tau\ln m \le 0\notag,
\end{align}
hence $u_{\lambda}(\theta)=0$. By the definition of $u_{\lambda}$, this implies that
\begin{equation}
\label{eq:prop2_caseb_2}
\min_{i\in[m]}
\bigl\{\tilde{\mathcal{L}}_i(\theta)-\tilde{\mathcal{L}}_i(\theta')\bigr\}
-\frac{\lambda}{2}\|\theta-\theta'\|^2
\le 0,
\qquad \forall\,\theta'\in \mathbb{R}^d.
\end{equation}
For any $\hat{\theta}\in \mathbb{R}^d$ and any $t\in(0,1)$. Let
\[
\theta_t := (1-t)\theta+t\hat{\theta} \in \mathbb{R}^d.
\]
Plugging $\theta'=\theta_t$ into Eq.~\eqref{eq:prop2_caseb_2}, we obtain
\begin{equation}
\label{eq:prop2_caseb_3}
\min_{i\in[m]} \bigl\{\tilde{\mathcal{L}}_i(\theta)-\tilde{\mathcal{L}}_i(\theta_t)\bigr\} -\frac{\lambda}{2}\|\theta-\theta_t\|^2 \le 0.
\end{equation}
Assume that for all $i\in[m]$, $\tilde{\mathcal{L}}_i$ is convex at $\theta$. Then for each $i\in[m]$,
\[
\tilde{\mathcal{L}}_i(\theta_t) \le t\,\tilde{\mathcal{L}}_i(\hat{\theta})+(1-t)\tilde{\mathcal{L}}_i(\theta),
\]
which implies
\[
\tilde{\mathcal{L}}_i(\theta)-\tilde{\mathcal{L}}_i(\theta_t) \ge t\bigl(\tilde{\mathcal{L}}_i(\theta)-\tilde{\mathcal{L}}_i(\hat{\theta})\bigr).
\]
Substituting into Eq.~\eqref{eq:prop2_caseb_3} and using $\|\theta-\theta_t\|^2=t^2\|\hat{\theta}-\theta\|^2$ yields
\begin{align}
\min_{i\in[m]}
\Bigl\{t\bigl(\tilde{\mathcal{L}}_i(\theta)-\tilde{\mathcal{L}}_i(\hat{\theta})\bigr)\Bigr\} -\frac{\lambda}{2}\,t^2\|\hat{\theta}-\theta\|^2 \le 0 \notag.
\end{align}
Dividing both sides by $t>0$ and letting $t\rightarrow 0$, we obtain for all $\hat{\theta}\in \mathbb{R}^d$,
\begin{align}
\min_{i\in[m]} \bigl\{\tilde{\mathcal{L}}_i(\theta)-\tilde{\mathcal{L}}_i(\hat{\theta})\bigr\} \le 0 \notag,
\end{align}
which means that there does not exist $\hat{\theta}\in \mathbb{R}^d$ such that $\tilde{\mathcal{L}}_i(\hat{\theta})<\tilde{\mathcal{L}}_i(\theta)$ for all $i\in[m]$. Therefore, $\theta$ is weakly Pareto optimal.
\end{proof}


% \begin{lemma}
% \label{lem:existence_v_full}
% If $\tilde{\mathcal{L}}_i(\theta)$, $i\in[m]$, is continuous and weakly convex with modulus $\mu_i$,
% and
% \[
% \lambda \ge -\min_{i\in[m]}\mu_i,
% \]
% then there exists $\theta^\star\in \mathbb{R}^d$ such that
% \[
% v_{\lambda,\tau}(\theta^\star)=-\tau\ln m.
% \]
% \end{lemma}

% \begin{proof}[Proof of Lemma~\ref{lem:existence_v_full}]
% From Lemma~\ref{lemmsc}, since
% $\lambda \ge -\min_{i\in[m]}\mu_i$,
% the function $h_{\lambda,\tau}(\theta,\theta')$ is convex with respect to $\theta'$.
% Let $P$ be the indicator function of $\Theta$. Then
% \[
% 0\in \nabla_{\theta'} h_{\lambda,\tau}(\theta,\theta') + \partial P(\theta')
% \quad\Longleftrightarrow\quad
% \theta' \in \arg\min_{\theta'\in \mathbb{R}^d} h_{\lambda,\tau}(\theta,\theta').
% \]

% By the definition of $h_{\lambda,\tau}(\theta,\theta')$, the gradient
% $\nabla_{\theta'} h_{\lambda,\tau}(\theta,\theta')$ can be written as
% \begin{equation}
% \label{eq:grad_theta_prime_h}
% \nabla_{\theta'} h_{\lambda,\tau}(\theta,\theta')
% =
% \sum_{i=1}^m
% \frac{\exp\!\left(\frac{\tilde{\mathcal{L}}_i(\theta')-\tilde{\mathcal{L}}_i(\theta)}{\tau}\right)}
% {\sum_{j=1}^m \exp\!\left(\frac{\tilde{\mathcal{L}}_j(\theta')-\tilde{\mathcal{L}}_j(\theta)}{\tau}\right)}
% \,\nabla \tilde{\mathcal{L}}_i(\theta')
% +
% \lambda(\theta'-\theta).
% \end{equation}

% When $\theta'=\theta$, we further have
% \begin{equation}
% \label{eq:grad_diag}
% \nabla_{\theta'} h_{\lambda,\tau}(\theta,\theta)
% =
% \frac{1}{m}\sum_{i=1}^m \nabla \tilde{\mathcal{L}}_i(\theta).
% \end{equation}

% Now recall that for any $\alpha\in\Delta^m$ (the probability simplex),
% the scalarization $\alpha^\top \psi(\theta)=\sum_{i=1}^m \alpha_i \tilde{\mathcal{L}}_i(\theta)$
% is lower bounded and continuous since $\tilde{\mathcal{L}}_i$ are continuous for all $i\in[m]$.
% Assume either $\Theta$ is compact, or $\Theta=\mathbb{R}^d$ and each $\tilde{\mathcal{L}}_i$
% is coercive. Then the minimizer of $\alpha^\top \psi(\theta)$ exists.
% In particular, let
% \[
% \theta^\star \in \arg\min_{\theta\in \mathbb{R}^d}\frac{1}{m}\sum_{i=1}^m \tilde{\mathcal{L}}_i(\theta).
% \]
% The first-order optimality condition yields
% \begin{equation}
% \label{eq:opt_cond_avg}
% 0\in \frac{1}{m}\sum_{i=1}^m \nabla \tilde{\mathcal{L}}_i(\theta^\star) + \partial P(\theta^\star).
% \end{equation}

% Combining~\eqref{eq:opt_cond_avg} with~\eqref{eq:grad_diag}, we obtain
% \[
% 0\in \nabla_{\theta'} h_{\lambda,\tau}(\theta,\theta') + \partial P(\theta')
% \Big|_{(\theta,\theta')=(\theta^\star,\theta^\star)}.
% \]
% Therefore,
% \[
% \theta^\star \in \arg\min_{\theta'\in \mathbb{R}^d} h_{\lambda,\tau}(\theta^\star,\theta').
% \]
% It follows that
% \begin{align}
% \min_{\theta'\in \mathbb{R}^d} h_{\lambda,\tau}(\theta^\star,\theta')
% &=
% h_{\lambda,\tau}(\theta^\star,\theta^\star)
% =
% \tau \ln\!\left(\sum_{i=1}^m
% \exp\!\left(\frac{\tilde{\mathcal{L}}_i(\theta^\star)-\tilde{\mathcal{L}}_i(\theta^\star)}{\tau}\right)\right)
% \nonumber\\
% &=
% \tau \ln m.
% \label{eq:min_h_star}
% \end{align}

% By definition,
% \begin{equation}
% \label{eq:v_star}
% v_{\lambda,\tau}(\theta^\star)
% =
% -
% \min_{\theta'\in \mathbb{R}^d} h_{\lambda,\tau}(\theta^\star,\theta')
% =
% -\tau\ln m,
% \end{equation}
% which completes the proof.
% \end{proof}

\subsection{Proofs in Section~\ref{analysis}}
\begin{proof}[Proof of Lemma~\ref{smoothnessh}]
According to proof of Proposition~\ref{prop:smooth_approx}, we know $h_{0,\tau}(\theta,\theta')$ is $\big(\frac{\bar{\ell}^2}{\tau}+\bar{\zeta}\big)$-smooth w.r.t $\theta, \theta'\in \mathbb{R}^d$. Then, we have $h_{\lambda,\tau}(\theta,\theta')$ is $\big(\frac{\bar{\ell}^2}{\tau}+\bar{\zeta}+ \lambda \big) $-smooth, where the $\lambda$ term comes from the quadratic regularizer $\frac{\lambda}{2}\|\theta-\theta'\|^2$.
\end{proof}


\begin{lemma}[Lemma 4.3 in~\cite{lin2020gradient}]
For TTGDA, let $\delta^k=\bigl\|\theta^{\star}_{\lambda,\tau}(\theta^k)-\theta'^k\bigr\|^2$,
Then the following statements hold true.
\begin{equation}
\label{eq:delta}
\delta^k\le\Bigl(1-\frac{1}{2\kappa}+4\kappa^3\ell^2\eta_{\theta}^2\Bigr)\delta^{k-1}+4\kappa^3\eta_{\theta}^2\bigl\|\nabla v_{\lambda,\tau}(\theta^{k-1})\bigr\|^2,
\end{equation}
and 
\begin{equation}
\label{eq:vdescent}
v_{\lambda,\tau}(\theta^k)\le v_{\lambda,\tau}(\theta^{k-1})-\frac{7\eta_{\theta}}{16}\bigl\|\nabla v_{\lambda,\tau}(\theta^{k-1})\bigr\|^2+\frac{9\eta_{\theta}\ell^2}{16}\delta^{k-1}.
\end{equation}
\end{lemma}


\begin{proof}[Proof of Theorem~\ref{thm:gda_complexity}]
Let
\[
\delta^k
=
\bigl\|
\theta^{\star}_{\lambda,\tau}(\theta^k)
-
\theta'^k
\bigr\|^2,
\qquad
g^{k}
:=
\bigl\|
\nabla v_{\lambda,\tau}(\theta^{k})
\bigr\|^2.
\]
Define
\[
\gamma
:=
1-\frac{1}{2\kappa}+4\kappa^3\ell^2\eta_\theta^2.
\]
By Eq.~\eqref{eq:delta}, for all $k\ge 1$,
\begin{equation}
\label{eq:rec_delta_C1}
\delta^k \le \gamma\,\delta^{k-1}+4\kappa^3\eta_\theta^2\, g^{k-1}.
\end{equation}
Unrolling Eq.~\eqref{eq:rec_delta_C1} yields
\begin{equation}
\label{eq:delta_unroll_C1}
\delta^k
\le
\gamma^k \delta^0
+
4\kappa^3\eta_\theta^2
\sum_{j=0}^{k-1}\gamma^{k-1-j} g^{j}.
\end{equation}

Next, by Eq.~\eqref{eq:vdescent}, for all $k\ge 1$,
\begin{equation}
\label{eq:descent_C1}
v_{\lambda,\tau}(\theta^k)
\le
v_{\lambda,\tau}(\theta^{k-1})
-
\frac{7\eta_\theta}{16}\, g^{k-1}
+
\frac{9\eta_\theta\ell^2}{16}\,\delta^{k-1}.
\end{equation}
Summing Eq.~\eqref{eq:descent_C1} over $k=1,\dots,K$ gives
\begin{align}
\frac{7\eta_\theta}{16}\sum_{k=0}^{K-1} g^{k}
\le
v_{\lambda,\tau}(\theta^{0})-v_{\lambda,\tau}(\theta^{K})
+
\frac{9\eta_\theta\ell^2}{16}\sum_{k=0}^{K-1}\delta^{k}\notag.
\end{align}
Using $v_{\lambda,\tau}(\theta^{K})\ge \min_{\theta}v_{\lambda,\tau}(\theta)$ and
$\Delta_v:=v_{\lambda,\tau}(\theta^{0})-\min_{\theta}v_{\lambda,\tau}(\theta)$,
we obtain
\begin{equation}
\label{eq:sum_descent_C1b}
\frac{7\eta_\theta}{16}\sum_{k=0}^{K-1} g^{k}
\le
\Delta_v
+
\frac{9\eta_\theta\ell^2}{16}\sum_{k=0}^{K-1}\delta^{k}.
\end{equation}

We now bound $\sum_{k=0}^{K-1}\delta^{k}$ using Eq.~\eqref{eq:delta_unroll_C1}. Summing Eq.~\eqref{eq:delta_unroll_C1} over $k=1,\dots,K$ and exchanging the order of summation yield
\begin{align}
\sum_{k=1}^{K}\delta^{k}
&\le
\delta^{0}\sum_{k=1}^{K}\gamma^{k}
+
4\kappa^3\eta_\theta^2
\sum_{k=1}^{K}\sum_{j=0}^{k-1}\gamma^{k-1-j} g^{j} \nonumber\\
&=
\delta^{0}\sum_{k=1}^{K}\gamma^{k}
+
4\kappa^3\eta_\theta^2
\sum_{j=0}^{K-1} g^{j}\sum_{k=j+1}^{K}\gamma^{k-1-j} \nonumber\\
&\le
\delta^{0}\frac{\gamma}{1-\gamma}
+
4\kappa^3\eta_\theta^2
\sum_{j=0}^{K-1} g^{j}\frac{1}{1-\gamma},
\label{eq:sum_delta_C1}
\end{align}
where we used $\sum_{k=1}^{K}\gamma^{k}\le \frac{\gamma}{1-\gamma}$ and $\sum_{k=j+1}^{K}\gamma^{k-1-j}\le \frac{1}{1-\gamma}$.

Substituting Eq.~\eqref{eq:sum_delta_C1} into Eq.~\eqref{eq:sum_descent_C1b} gives
\begin{align}
\frac{7\eta_\theta}{16}\sum_{k=0}^{K-1} g^{k}
&\le
\Delta_v
+
\frac{9\eta_\theta\ell^2}{16}
\left(
\delta^{0}\frac{\gamma}{1-\gamma}
+
4\kappa^3\eta_\theta^2\frac{1}{1-\gamma}
\sum_{k=0}^{K-1} g^{k}
\right).
\label{eq:pre_rearrange_C1}
\end{align}
Rearranging Eq.~\eqref{eq:pre_rearrange_C1} yields
\begin{equation}
\label{eq:rearranged_C1}
\left(
\frac{7\eta_\theta}{16}
-
\frac{9\eta_\theta\ell^2}{16}
\cdot
\frac{4\kappa^3\eta_\theta^2}{1-\gamma}
\right)
\sum_{k=0}^{K-1} g^{k}
\le
\Delta_v
+
\frac{9\eta_\theta\ell^2}{16}
\cdot
\frac{\gamma}{1-\gamma}\,\delta^{0}.
\end{equation}

Finally, choose $\eta_\theta=\Theta(1/(\kappa^2\ell))$ so that
$4\kappa^3\ell^2\eta_\theta^2 \le \frac{7}{32\kappa}$, which implies
\[
1-\gamma
=
\frac{1}{2\kappa}-4\kappa^3\ell^2\eta_\theta^2
\ge
\frac{9}{32\kappa},
\qquad
\frac{1}{1-\gamma}\le \frac{32}{9}\kappa,
\qquad
\frac{4\kappa^3\eta_\theta^2}{1-\gamma}\le \frac{128}{9}\kappa^4\eta_\theta^2.
\]
Under this choice, the coefficient on the left-hand side of Eq.~\eqref{eq:rearranged_C1} is lower bounded by a constant multiple of $\eta_\theta$, and thus
\begin{align}
\sum_{k=0}^{K-1} g^{k}
\le
\mathcal{O}\!\left(
\frac{\Delta_v}{\eta_\theta}
+
\kappa\ell^2\,\delta^{0}
\right)\notag.
\end{align}
Let $\widehat{k}$ be sampled uniformly from $\{0,\dots,K-1\}$. Then
\[
\mathbb{E}\bigl\|\nabla v_{\lambda,\tau}(\theta^{\widehat{k}})\bigr\|^2
=
\frac{1}{K}\sum_{k=0}^{K-1} g^{k}
\le
\mathcal{O}\!\left(
\frac{\Delta_v}{K\eta_\theta}
+
\frac{\kappa\ell^2\,\delta^{0}}{K}
\right).
\]
Setting the right-hand side to be at most $\epsilon^2$ and using
$\eta_\theta=\Theta(1/(\kappa^2\ell))$ gives the iteration (and gradient)
complexity
\[
K
=
\mathcal{O}\!\left(
\frac{\kappa^2\ell\,\Delta_v+\kappa\ell^2 \delta^{0}}{\epsilon^2}
\right).
\]
This concludes the proof.
\end{proof}





% \section{Nesterov’s Acceleration and Adam Updates}
% To improve the convergence speed of the alternating min--max updates, we can
% incorporate Nesterov's acceleration and Adam update into GDA.
% \paragraph{GDA with Nesterov's acceleration.}
% At iteration $k$, we update $\theta'$ and $\theta$ by the Nesterov accelerated updates:
% \begin{align}
% g_{\theta'}^{k} 
% &:= \nabla_{\theta'} h_{\lambda,\tau}(\theta^{k},\theta'^{k}), \notag\\
% v_{\theta'}^{k+1} 
% &:= \beta v_{\theta'}^{k} + g_{\theta'}^{k}, \notag\\
% \theta'^{k+1} 
% &\leftarrow \Pi_{\Theta}\!\left(
% \theta'^{k} - \eta_{\theta'}\bigl(g_{\theta'}^{k}+\beta v_{\theta'}^{k+1}\bigr)
% \right), \notag\\
% g_{\theta}^{k} 
% &:= \nabla_{\theta} h_{\lambda,\tau}(\theta^{k},\theta'^{k+1}), \notag\\
% v_{\theta}^{k+1} 
% &:= \beta v_{\theta}^{k} + g_{\theta}^{k}, \notag\\
% \theta^{k+1} 
% &\leftarrow \Pi_{\Theta}\!\left(
% \theta^{k} + \eta_{\theta}\bigl(g_{\theta}^{k}+\beta v_{\theta}^{k+1}\bigr)
% \right)\notag,
% \end{align}
% where $\beta\in(0,1)$ is the momentum parameter.



% \paragraph{GDA with Adam updates.}
% For $k=0,1,\ldots,K-1$, perform the following updates:
% \begin{align}
% g_{\theta'}^{k} &:= \nabla_{\theta'} h_{\lambda,\tau}(\theta^{k},\theta'^{k}), \notag\\
% m_{\theta'}^{k+1} &:= \beta_1 m_{\theta'}^{k} + (1-\beta_1) g_{\theta'}^{k}, \qquad
% s_{\theta'}^{k+1} := \beta_2 s_{\theta'}^{k} + (1-\beta_2) (g_{\theta'}^{k})^2, \notag\\
% \hat m_{\theta'}^{k+1} &:= \frac{m_{\theta'}^{k+1}}{1-\beta_1^{k+1}}, \qquad
% \hat s_{\theta'}^{k+1} := \frac{s_{\theta'}^{k+1}}{1-\beta_2^{k+1}}, \notag\\
% \theta'^{k+1} &\leftarrow \Pi_{\Theta}\!\left(
% \theta'^{k} - \eta_{\theta'} \frac{\hat m_{\theta'}^{k+1}}{\sqrt{\hat s_{\theta'}^{k+1}}+\epsilon}
% \right), \notag\\
% g_{\theta}^{k} &:= \nabla_{\theta} h_{\lambda,\tau}(\theta^{k},\theta'^{k+1}), \notag\\
% m_{\theta}^{k+1} &:= \beta_1 m_{\theta}^{k} + (1-\beta_1) g_{\theta}^{k}, \qquad
% s_{\theta}^{k+1} := \beta_2 s_{\theta}^{k} + (1-\beta_2) (g_{\theta}^{k})^2, \notag\\
% \hat m_{\theta}^{k+1} &:= \frac{m_{\theta}^{k+1}}{1-\beta_1^{k+1}}, \qquad
% \hat s_{\theta}^{k+1} := \frac{s_{\theta}^{k+1}}{1-\beta_2^{k+1}}, \notag\\
% \theta^{k+1} &\leftarrow \Pi_{\Theta}\!\left(
% \theta^{k} + \eta_{\theta} \frac{\hat m_{\theta}^{k+1}}{\sqrt{\hat s_{\theta}^{k+1}}+\epsilon}
% \right)\notag,
% \end{align}
% where $\beta_1,\beta_2\in(0,1)$
% are decay rates, and $\epsilon>0$ is a small constant.


\begin{table*}[t]
\caption{Ablation studies on the impact of $\tau$ on the NYUv2 dataset. The best results for each task are shown in \textbf{bold}. $\uparrow$($\downarrow$) means that the higher (lower) the value, the better the performance.}
\label{table:ablation_tau_detail}
\centering
\setlength{\tabcolsep}{5pt}
\renewcommand{\arraystretch}{1.15}
\begin{tabular}{l cc cc cc ccc c}
\toprule
\multirow{3}{*}[-1ex]{$\tau$}
& \multicolumn{2}{c}{\textbf{Segmentation}}
& \multicolumn{2}{c}{\textbf{Depth}}
& \multicolumn{5}{c}{\textbf{Surface Normal}}
& \multirow{3}{*}[-1ex]{$\Delta_b \uparrow$} \\
\cmidrule(lr){2-3}\cmidrule(lr){4-5}\cmidrule(lr){6-10}
& \multirow{2}{*}[-0.4ex]{\textbf{mIoU}$\uparrow$}
& \multirow{2}{*}[-0.4ex]{\textbf{Pix Acc}$\uparrow$}
& \multirow{2}{*}[-0.4ex]{\textbf{Abs Err}$\downarrow$}
& \multirow{2}{*}[-0.4ex]{\textbf{Rel Err}$\downarrow$}
& \multicolumn{2}{c}{\textbf{Angle Distance}}
& \multicolumn{3}{c}{\textbf{Within $t^\circ$}}
& \\
\cmidrule(lr){6-7}\cmidrule(lr){8-10}
& & & & 
& \textbf{Mean}$\downarrow$
& \textbf{Median}$\downarrow$
& \textbf{11.25}$\uparrow$
& \textbf{22.5}$\uparrow$
& \textbf{30}$\uparrow$
& \\
\midrule
$\tau=0.1$
& 65.21 & 83.46
& 0.2982 & 0.1141
& 18.71 & 13.63
& 42.59 & 71.28 & 81.38
& +19.10\% \\
$\tau=0.5$
& 66.58 & 84.41
& 0.2822 & 0.1100
& 17.05 & 11.87
& 48.49 & 75.30 & 84.20
& +23.84\% \\
$\tau=1$
& \textbf{68.24} & \textbf{84.78}
& 0.2788 & 0.1094
& 16.93 & 11.81
& 48.71 & 75.62 & 84.49
& +24.80\% \\
$\tau=3$
& 66.92 & 84.43
& 0.2738 & 0.1068
& 16.79 & 11.69
& 49.14 & 76.07 & 84.76
& +25.06\% \\
$\tau=5$
& 67.12 & 84.70
& \textbf{0.2722} & \textbf{0.1050}
& 16.85 & 11.65
& 49.35 & 75.85 & 84.54
& +\textbf{25.44\%} \\
$\tau=10$
& 67.35 & 84.32
& 0.2772 & 0.1088
& \textbf{16.73} & \textbf{11.56}
& \textbf{49.59} & \textbf{76.15} & \textbf{84.77}
& +24.97\% \\
\bottomrule
\end{tabular}
\end{table*}

\begin{table*}[t]
\caption{The impact of $\lambda$ on the NYUv2 dataset. The best results for each task are shown in \textbf{bold}. $\uparrow$($\downarrow$) means that the higher (lower) the value, the better the performance.}
\label{table:ablation_lambda_detail}
\centering
\setlength{\tabcolsep}{5pt}
\renewcommand{\arraystretch}{1.15}
\begin{tabular}{c cc cc cc ccc c}
\toprule
\multirow{3}{*}[-1ex]{$\lambda$}
& \multicolumn{2}{c}{\textbf{Segmentation}}
& \multicolumn{2}{c}{\textbf{Depth}}
& \multicolumn{5}{c}{\textbf{Surface Normal}}
& \multirow{3}{*}[-1ex]{$\Delta_b \uparrow$} \\
\cmidrule(lr){2-3}\cmidrule(lr){4-5}\cmidrule(lr){6-10}
& \multirow{2}{*}[-0.4ex]{\textbf{mIoU}$\uparrow$}
& \multirow{2}{*}[-0.4ex]{\textbf{Pix Acc}$\uparrow$}
& \multirow{2}{*}[-0.4ex]{\textbf{Abs Err}$\downarrow$}
& \multirow{2}{*}[-0.4ex]{\textbf{Rel Err}$\downarrow$}
& \multicolumn{2}{c}{\textbf{Angle Distance}}
& \multicolumn{3}{c}{\textbf{Within $t^\circ$}}
& \\
\cmidrule(lr){6-7}\cmidrule(lr){8-10}
& & & & 
& \textbf{Mean}$\downarrow$
& \textbf{Median}$\downarrow$
& \textbf{11.25}$\uparrow$
& \textbf{22.5}$\uparrow$
& \textbf{30}$\uparrow$
& \\
\midrule
0
& 67.12 & 84.70
& \textbf{0.2722} & \textbf{0.1050}
& 16.85 & 11.65
& 49.35 & 75.85 & 84.54
& +25.44\% \\

$1.00\times10^{-6}$
& 67.57 & 84.67
& 0.2759 & 0.1080
& 16.84 & 11.69
& 49.14 & 75.98 & \textbf{84.70}
& +25.06\% \\

$1.00\times10^{-5}$
& \textbf{68.02} & 84.73
& 0.2742 & 0.1063
& \textbf{16.77} & \textbf{11.62}
& \textbf{49.42} & \textbf{76.00} & \textbf{84.70}
& \textbf{+25.57}\% \\

$1.00\times10^{-4}$
& 67.39 & \textbf{84.75}
& 0.2768 & 0.1080
& 16.83 & 11.66
& 49.26 & 75.86 & 84.57
& +25.00\% \\

$1.00\times10^{-3}$
& 67.21 & 84.64
& 0.2805 & 0.1091
& 17.00 & 11.77
& 48.85 & 75.51 & 84.32
& +24.41\% \\

\bottomrule
\end{tabular}
\end{table*}


\section{Experimental Details}

\subsection{Illustrative Example}
\label{appdix:example}
We detail the illustrative example in Section~\ref{sec:motivation}.
Let $\theta=(\theta_1,\theta_2)\in\mathbb{R}^2$, and consider the following intermediate objectives:
\begin{align*}
\tilde{\ell}_1(\theta) &= c_1(\theta) f_1(\theta) + c_2(\theta) g_1(\theta), \\
\tilde{\ell}_2(\theta) &= c_1(\theta) f_2(\theta) + c_2(\theta) g_2(\theta),
\end{align*}
where
\begin{align*}
f_1(\theta) &= \ln\!\Big(\max\big(\big|0.5(-\theta_1-7)-\tanh(-\theta_2)\big|,\;5\times 10^{-6}\big)\Big) + 6, \\
f_2(\theta) &= \ln\!\Big(\max\big(\big|0.5(-\theta_1+3)-\tanh(-\theta_2)+2\big|,\;5\times 10^{-6}\big)\Big) + 6, \\
g_1(\theta) &= \frac{(-\theta_1+7)^2 + 0.1\,(-\theta_2-8)^2}{10} - 20, \\
g_2(\theta) &= \frac{(-\theta_1-7)^2 + 0.1\,(-\theta_2-8)^2}{10} - 20, \\
c_1(\theta) &= \max\!\big(\tanh(0.5\theta_2),\,0\big), \qquad
c_2(\theta) = \max\!\big(\tanh(-0.5\theta_2),\,0\big).
\end{align*}

This example follows~\cite{liu2021conflict} and its modification in~\cite{navon2022multi}.
Building on the latter modification, we further adjust the objectives by adding a constant shift to keep both losses positive. Specifically, we define
\begin{equation*}
\hat{\ell}_1(\theta)=\tilde{\ell}_1(\theta)+30,\qquad \hat{\ell}_2(\theta)=\tilde{\ell}_2(\theta)+30,
\end{equation*}
and use the final objectives
\begin{equation*}
\ell_1(\theta)=0.1\,\hat{\ell}_1(\theta),\qquad \ell_2(\theta)=\hat{\ell}_2(\theta). 
\end{equation*}
\begin{equation*}
\text{and }\ell_1(\theta)=\,\hat{\ell}_1(\theta),\qquad \ell_2(\theta)=0.1\hat{\ell}_2(\theta). 
\end{equation*}

We use five initialization points $(-8.5,7.5)$, $(0.0,0.0)$, $(9.0,9.0)$, $(-7.5,-0.5)$, and $(9.0,-1.0)$. All methods use Adam with learning rate $10^{-3}$ for up to $35\mathrm{K}$ iterations. 
We approximate the Pareto front by evaluating an $800\times800$ uniform grid over $[-12,12]$ along each dimension and retaining the non-dominated points. At iteration $t$, we compute the distance from the current objective-space point $\mathbf{z}_t$ to the approximated Pareto set $\mathcal{P}$:
\[
d_t \;=\; \min_{\mathbf{p}\in\mathcal{P}} \|\mathbf{z}_t-\mathbf{p}\|_2 .
\]
If $d_t<0.05$, we stop early.


\begin{table*}[t]
\caption{The baseline with normalization on the NYUv2 dataset. $\uparrow$($\downarrow$) means that the higher (lower) the value, the better the performance.}
\label{table8}
\centering
\setlength{\tabcolsep}{5pt}
\renewcommand{\arraystretch}{1.15}
\begin{tabular}{c cc cc cc ccc c}
\toprule
\multirow{3}{*}[-1ex]{Method}
& \multicolumn{2}{c}{\textbf{Segmentation}}
& \multicolumn{2}{c}{\textbf{Depth}}
& \multicolumn{5}{c}{\textbf{Surface Normal}}
& \multirow{3}{*}[-1ex]{$\Delta_b \uparrow$} \\
\cmidrule(lr){2-3}\cmidrule(lr){4-5}\cmidrule(lr){6-10}
& \multirow{2}{*}[-0.4ex]{\textbf{mIoU}$\uparrow$}
& \multirow{2}{*}[-0.4ex]{\textbf{Pix Acc}$\uparrow$}
& \multirow{2}{*}[-0.4ex]{\textbf{Abs Err}$\downarrow$}
& \multirow{2}{*}[-0.4ex]{\textbf{Rel Err}$\downarrow$}
& \multicolumn{2}{c}{\textbf{Angle Distance}}
& \multicolumn{3}{c}{\textbf{Within $t^\circ$}}
& \\
\cmidrule(lr){6-7}\cmidrule(lr){8-10}
& & & & 
& \textbf{Mean}$\downarrow$
& \textbf{Median}$\downarrow$
& \textbf{11.25}$\uparrow$
& \textbf{22.5}$\uparrow$
& \textbf{30}$\uparrow$
& \\
\midrule
EW + EMA
& 67.19 & 84.32
& 0.2750 & 0.1076
& 17.04 & 11.93
& 48.24 & 75.35 & 84.27
& +24.52\% \\

STCH + Epoch-wise
& 66.83 & 84.24
& 0.2777 & 0.1095
& 17.23 & 12.07
& 47.76 & 74.92 & 83.94
& +23.79\% \\

SIMS (Ours)  & 67.12 & 84.70 & 0.2722 & 0.1050 & 16.85   & 11.65   & 49.35 & 75.85 & 84.54 & +25.44\% \\
\bottomrule
\end{tabular}
\end{table*}

\begin{table}[t]
\centering
\caption{Comparison of different methods on Ali-CCP.}
\label{tab:auc_results}
\begin{tabular}{lccc}
\toprule
Method & AUC/T1 & AUC/T2 & Avg AUC \\
\midrule
EW           & 0.5828 & 0.5774 & 0.5801 \\
FOOPS        & 0.5742 & 0.5779 & 0.5761 \\
STCH         & 0.5653 & 0.5979 & 0.5816 \\
SIMS (Ours)  & 0.5757 & 0.6012 & 0.5884 \\
\bottomrule
\end{tabular}
\end{table}


\begin{table}[t]
\centering
\caption{Effect of $\lambda$ on convergence and Pareto-front quality.}\label{table10}
\begin{tabular}{c c c}
\toprule
$\lambda$ & Mean Iter. to $r_k < 0.1$ & Mean Rank of dist to pf \\
\midrule
0      & 2463 & 1.750 \\
$1\times10^{-4}$ & 2526 & 2.333 \\
$1\times10^{-3}$ & 1385 & 2.917 \\
$1\times10^{-2}$ & 711  & 2.917 \\
$1\times10^{-1}$ & 198  & 2.833 \\
\bottomrule
\end{tabular}
\end{table}

\begin{table}[t]
\centering
\caption{Effect of different $\lambda$ values on convergence speed and distance to the Pareto front for problem F11.}
\label{tab:f11_lambda}
\resizebox{\columnwidth}{!}{
\begin{tabular}{cccccc}
\toprule
Problem & Init Point & $\lambda$ & Iter. to $r_k < 0.1$ & Dist. to PF & Rank \\
\midrule
F11 & $(-2.0, 2.0)$ & $0$       & $98$ & $0.000453952$ & $1$ \\
F11 & $(-2.0, 2.0)$ & $10^{-4}$ & $97$ & $0.000637059$ & $2$ \\
F11 & $(-2.0, 2.0)$ & $10^{-3}$ & $95$ & $0.002296451$ & $3$ \\
F11 & $(-2.0, 2.0)$ & $10^{-2}$ & $15$ & $0.00230228$  & $4$ \\
F11 & $(-2.0, 2.0)$ & $10^{-1}$ & $2$  & $0.004238246$ & $5$ \\
\bottomrule
\end{tabular}
}
\end{table}


\begin{table*}[t]
\caption{Performance on the NYUv2 dataset when the depth loss is rescaled by $100\times$. $\uparrow$($\downarrow$) means that the higher (lower) the value, the better the performance.}
\label{tab:rescale}
\centering
\setlength{\tabcolsep}{5pt}
\renewcommand{\arraystretch}{1.15}
\begin{tabular}{c cc cc cc ccc c}
\toprule
\multirow{3}{*}[-1ex]{Method}
& \multicolumn{2}{c}{\textbf{Segmentation}}
& \multicolumn{2}{c}{\textbf{Depth}}
& \multicolumn{5}{c}{\textbf{Surface Normal}}
& \multirow{3}{*}[-1ex]{$\Delta_b \uparrow$} \\
\cmidrule(lr){2-3}\cmidrule(lr){4-5}\cmidrule(lr){6-10}
& \multirow{2}{*}[-0.4ex]{\textbf{mIoU}$\uparrow$}
& \multirow{2}{*}[-0.4ex]{\textbf{Pix Acc}$\uparrow$}
& \multirow{2}{*}[-0.4ex]{\textbf{Abs Err}$\downarrow$}
& \multirow{2}{*}[-0.4ex]{\textbf{Rel Err}$\downarrow$}
& \multicolumn{2}{c}{\textbf{Angle Distance}}
& \multicolumn{3}{c}{\textbf{Within $t^\circ$}}
& \\
\cmidrule(lr){6-7}\cmidrule(lr){8-10}
& & & & 
& \textbf{Mean}$\downarrow$
& \textbf{Median}$\downarrow$
& \textbf{11.25}$\uparrow$
& \textbf{22.5}$\uparrow$
& \textbf{30}$\uparrow$
& \\
\midrule
EW
& 57.58 & 78.94
& 0.2755 & 0.1085
& 18.95 & 13.46
& 43.32 & 70.89 & 80.64
& +17.39\% \\


STCH
& 67.01 & 84.12
& 0.2835 & 0.1122
& 17.35 & 12.24
& 47.12 & 74.65 & 83.80
& +23.01\% \\

foops
& 65.23 & 83.26
& 6.0287 & 2.9620
& 19.03 & 14.02
& 41.19 & 70.40 & 80.87
& -536.99\% \\

SIMS
& 67.84 & 84.85
& 0.2778 & 0.1078
& 16.80 & 11.60
& 49.40 & 76.02 & 84.66
& +25.21\% \\
\bottomrule
\end{tabular}
\end{table*}

\subsection{Implementation Details}
\label{app:mtldetails}
We employ the following loss functions:

\noindent \textbf{Semantic $\&$ Human Parts Segmentation.}
Both tasks use pixel-wise Cross-Entropy Loss over valid pixels:
$$\mathcal{L}_{seg} = - \frac{1}{N_{valid}} \sum_{p \in \Omega} \mathbf{1}_{[y_p \neq \text{ignore}]} \log(P(y_p | x_p))$$
where $\Omega$ is the image domain, $y_p$ is the ground truth label, and the loss ignores the void class (index 255).

\noindent \textbf{Depth Estimation.}
We use the $L_1$ Loss (Mean Absolute Error) restricted to pixels with valid depth annotations:
$$\mathcal{L}_{depth} = \frac{1}{N_{valid}} \sum_{p \in \Omega} \mathbf{1}_{[d_p > 0]} | d_p - \hat{d}_p |$$where $d_p$ and $\hat{d}_p$ denote the ground truth and predicted depth, respectively.

\noindent \textbf{Surface Normal Prediction.}
We align predicted normals with ground truth via Inverse Cosine Similarity:
$$\mathcal{L}_{normal} = 1 - \frac{1}{N_{valid}} \sum_{p \in \Omega} \mathbf{1}_{[m_p = 1]} \langle \mathbf{n}_p, \hat{\mathbf{n}}_p \rangle$$where $\mathbf{n}_p, \hat{\mathbf{n}}_p$ are unit vectors, and $\langle \cdot, \cdot \rangle$ denotes the dot product.

\noindent \textbf{Edge Detection.}
To handle the extreme sparsity of edge pixels, we use a Class-Balanced Binary Cross-Entropy Loss:
$$\mathcal{L}_{edge} = - \sum_{p \in \Omega} \left[ \beta y_p \log(\hat{y}_p) + (1-y_p) \log(1-\hat{y}_p) \right]$$
where $\beta$ is a positive weight (set to 0.95) to increase the penalty for missing edge pixels.



\subsection{Metric for Each Task}
\label{appdix:metric}
\noindent \textbf{NYUv2 dataset.} For semantic segmentation, we use mean Intersection over Union (mIoU) and Pixel Accuracy (Pix Acc); for depth prediction, Absolute Error (Abs Err) and Relative Error (Rel Err); and for surface normal estimation, mean and median angular errors in degrees, pixel percentages within 11.25, 22.5, and 30 degrees.


\noindent \textbf{CityScapes dataset.} For semantic segmentation, using mean Intersection over Union (mIoU) and Pixel Accuracy (Pix Acc); for depth prediction, Absolute Error (Abs Err) and Relative Error (Rel Err).

\noindent \textbf{PASCAL-Context dataset.} For semantic segmentation, human parts segmentation, and saliency estimation, we use mean Intersection over Union (mIoU); for surface normal estimation, mean angular error in degrees.


\noindent \textbf{The selection of $\tau$.}
As indicated by our ablation study, the performance is relatively insensitive to the choice of $\tau$ within a reasonable range. Specifically, for the reported results, we adopt $\tau=5$ for NYUv2 and $\tau=3$ for both Cityscapes and PASCAL-Context.

\section{Additional Experimental Studies}

\subsection{Effect of $\tau$}

\label{appdix:tau}
The complete table is shown in the Table.~\ref{table:ablation_tau_detail}.



\subsection{Effect of $\lambda$}
\label{appdix:lambda}
The complete table is shown in the Table.~\ref{table:ablation_lambda_detail}.

\begin{add}


\section{Additional Experiment}
\subsection{Normalization baselines}
We add comparisons with normalization baselines. \zbedit{Table~\ref{table8} shows the results of the baseline (EW and STCH) with normalization (EMA and Epoch-wise) on the NYUv2 dataset.} We see that EW + EMA improves over vanilla EW, but it still underperforms SIMS (24.52\% vs. 25.44\%). In contrast, STCH + Epoch-wise shows no gain (gain=-0.02\%), suggesting that simply attaching online normalization to a merit-based objective does not necessarily yield a better trade-off and may even disturb the original learning objective. Overall, these results indicate that the gain of SIMS comes from its principled scale-invariant scalarization design with weak Pareto property.
\subsection{Non-vision MTL scenarios}
We add a non-vision experiment on Ali-CCP \cite{ma2018entire}, where the CVR loss is 1–2 orders of magnitude smaller than CTR. \zbedit{Table~\ref{tab:auc_results} shows the results.} On a standard Shared-Bottom model, SIMS achieves the best average AUC (0.5884), mainly by improving the smaller-scale task. This supports the generalization of SIMS beyond vision tasks.
\subsection{Synthetic Experiments on $\lambda$}
To better address the concern on the gap between the theory and the empirical setting of $\lambda$, we added additional synthetic experiments. According to the bi-objective BBOB \cite{brockhoff2022using} benchmark, we selected F1, F2, and F11 because they provide a controlled progression in difficulty and geometric structure.
% F1 is the simplest and most symmetric case, F2 introduces heterogeneous conditioning by combining two objectives with different properties, and F11 further increases the difficulty by making both objectives highly ill-conditioned. This choice allows us to examine our analysis under problems with different structures and synthetic setting.
Specifically, F1 is formed by two Sphere functions, F2 by a Sphere function and a separable Ellipsoid function, and F11 by two separable Ellipsoid functions. We set centers as $(-3.0,-3.0)$ and $(3.0,3.0)$ for F1, $(-3.0,-2.0)$ and $(3.0,2.0)$ for F2, and $(-3.0,-2.0)$ and $(3.0,2.0)$ for F11, and all constant shift to be 1.

We choice $\lambda$ from the set $\{0, 10^{-1}, 10^{-2}, 10^{-3}, 10^{-4}\}$ and the initial points from the set $\{(-2.0,2.0), (2.0,-2.0), (3.0,1.0),(5.0,-5.0)\}$ to examine both convergence behavior and solution quality.
% On these three benchmarks, we compare SIMS under different choices of $\lambda \in {0, 10^{-1}, 10^{-2}, 10^{-3}, 10^{-4}}$ from multiple initial points, i.e., $(-2.0,2.0)$, $(2.0,-2.0)$, $(3.0,1.0)$, and $(5.0,-5.0)$, to examine both convergence behavior and solution quality. 
Empirically, we find that all tested values of $\lambda$ converge on these problems. Moreover, under the overall ranking across the three benchmarks and multiple initializations, where the ranking is defined by ordering different $\lambda$ values according to the distance between their final converged solutions and the Pareto front (pf), with smaller distance corresponding to a lower rank value, $\lambda=0$ achieves the smallest distance to the Pareto front in most cases (Mean Rank=1.75, \zbedit{as shown in Table~\ref{table10}}). This is consistent with our analysis: while a larger positive $\lambda$ improves the conditioning of the inner problem, it also changes the surrogate itself and may introduce a bias relative to the original merit-function characterization, whereas $\lambda=0$ remains the closest smooth approximation to the original objective.

To better characterize the role of $\lambda$, we further introduce the stationarity residual $$ r_k=\sqrt{|\nabla_{\theta} h_{\lambda,\tau}(\theta_k,\theta_k')|^2+|\nabla_{\theta'} h_{\lambda,\tau}(\theta_k,\theta_k')|^2 }.$$ Here $r_k$ can be used to measure how close the current solution to a first-order stationary point. We use $r_k$ to study the convergence to an $\epsilon$-stationary point in the $\lambda> \bar\mu$ regime. \zbedit{From the results in Table~\ref{table10},} we observe that as $\lambda$ increases, the mean first hitting iteration for $r_k<0.1$ becomes smaller, indicating that larger $\lambda$ generally leads to faster stabilization of the optimization dynamics. This is also in line with our theory, since a larger positive $\lambda$ improves the curvature of the inner problem and makes the optimization procedure better conditioned.

In summary, these observations indicate that choosing $\lambda$ is essentially a trade-off (Table~\ref{tab:f11_lambda} provides a representative example on problem F11 with initialization at $(-2.0,, 2.0)$). A larger positive $\lambda$ improves the conditioning of the inner problem and tends to accelerate the convergence, while a smaller $\lambda$ preserves a more faithful characterization of the original merit-function objective and often yields a smaller distance to the Pareto front (Sec.\ref{sec:ablation}). Hence, $\lambda$ balances the convergence rate and fidelity to the original problem formulation. How to choose $\lambda$ in a more principled or adaptive way is an interesting direction that deserves further investigation in our future work.




\subsection{Rescale task losses on a real benchmark}
We additionally rescale the depth loss by $100\times$ on NYUv2 to evaluate the robustness of different methods under real-data task-wise rescaling. As shown in Table \ref{tab:rescale} and \ref{table:nyu}, EW is clearly affected by the enlarged depth loss, with $\Delta_b$ dropping from $22.34\%$ to $17.39\%$. Although STCH adopts objective normalization, it still exhibits a mild performance decline, with $\Delta_b$ decreasing from $23.81\%$ to $23.01\%$. FOOPS is even more severely affected by the rescaling, the optimization on the depth task collapses, leading to extremely large depth errors and also degrading the performance of the other tasks, which results in a substantially worse overall score. In contrast, SIMS remains highly stable, with $\Delta_b$ only changing from $25.44\%$ to $25.21\%$. These results provide direct real-data evidence that artificial task-wise rescaling can significantly affect the scalarization methods, while SIMS remains robust under loss rescaling, consistent with our scale-invariance claim.

\end{add}




\end{document}
\endinput
%%
%% End of file `sample-sigconf.tex'.
